\documentclass[11pt]{article}
\usepackage[a4paper,margin=1in]{geometry}
\usepackage{amsmath,amssymb}
\usepackage{booktabs}
\usepackage[numbers,sort&compress]{natbib}
\usepackage{hyperref}

\title{Mathieu-Harmonic Mass Spectrum of Elementary Particles\\
in Induced Scalar-Pressure Gravity}
\author{George Lotishvili}
\date{Version 4.7}

\begin{document}
\maketitle

\section*{Introduction}

\textit{Physical picture.}---Consider a guitar string.
When plucked, it vibrates at a definite note
(frequency) with large amplitude. The vibration
transfers energy to the guitar body, which radiates
it as sound. Over time the string grows quieter---but
the \textbf{note never changes}. Only the amplitude
(energy) decreases; the frequency is invariant.

The ISPG universe behaves the same way. Space is a
medium vibrating at a single fundamental carrier frequency
$\nu_0 \approx 1.91 \times 10^{30}$~Hz---the
MASS/Mathieu-calibrated ``note'' of the cosmos.
This vibration is energy;
its macroscopic manifestation is \emph{pressure}.
Through nonlinear self-interaction, the vibration
spontaneously condenses into localized resonances
(oscillons)---these are particles. The particle
spectrum is treated as a Mathieu-organized resonance
ladder built on this carrier; this quantitative ladder is
not derived by the frequency-invariance theorem alone.
The oscillon, in turn, generates a static gravitational
field---the ``sound'' radiated into the surrounding
medium. The original $\nu_0$-mode vibration
depletes as it condenses into matter, but the
matter itself radiates faint resonant echoes back
into the medium. The accumulated echoes \emph{are}
dark energy---a slowly growing stochastic
background ($w$ slightly below $-1$) whose growth
decelerates via negative feedback. The note $\nu_0$
remains the same; the partition of energy between
the fundamental mode, matter, and the echo
background evolves (\S\ref{sec:dark_energy}).

\medskip

In the ISPG (Induced Scalar-Pressure Gravity) framework
\cite{Lotishvili2026,Lotishvili2026Q,Horndeski1974}, space
is described as a scalar-metric medium possessing pressure, flow,
and the capacity for wave propagation. Particles (electrons, quarks,
bosons) are \textbf{oscillons} of this medium---localized resonant
excitations \cite{Copeland1995,GleiserWatkins1993}.

The central question of this work is: \textbf{if all particles are
excitations of a single medium, are their masses determined by the
fundamental properties of that medium?}

The answer is affirmative. We show that:
\begin{itemize}
  \item the medium possesses a single universal proper-time
    carrier frequency $\nu_0$;
  \item the massive-particle spectrum is organized as a
    Mathieu resonance ladder built on this carrier; neutrinos are identified as
    \emph{polarisation-rotation waves}---helical
    excitations whose oscillation plane rotates as
    they propagate---whose tiny mass reflects the
    suppressed torsional rigidity of the nonlinear
    $G_2(X)$ coupling (\S\ref{sec:neutrinos});
  \item the mass formula $m_N = m_1\,b_N(q)$ (Mathieu
    characteristic value~\cite{McLachlan1947,DLMFMathieu})
    compresses the charged leptons, quarks, and selected
    heavy bosons onto a common integer ladder; in the
    current PDG-style scan, $10/12$ entries are below
    $1\%$ and $11/12$ below $5\%$, with the $u$ quark
    ($5.62\%$) as the outlier and $b,H,c,t$ all below
    $0.07\%$~\cite{PDG2024}; $e$ and $\mu$ are calibration
    anchors, and the physical rule selecting each
    particle's $N$ remains an open derivation target;
  \item the electron is the first fully stable occupied charged
    Mathieu mode under the adopted gap-threshold criterion:
    $N=1$--$3$ modes lie in the wide-gap instability zone,
    while $N=4$ is a barely-stable resonance candidate
    predicted at $\sim 329$~keV (Sec.~\ref{sec:N4_predictions});
  \item the parameter $q$ follows from the shift-symmetric
    $G_2(X)$ extension of the ISPG action \cite{Nicolis2009,Deffayet2009}---the same
    extension independently established for ghost-free
    UV completion in Appendix~11 (Stability)---without
    introducing new fields or breaking any symmetry;
  \item the integer-harmonic structure of the scalar sector
    ($j_0$ zeros $= n\pi$) is the spectral signature of a
    one-dimensional substrate; the requirement for tensor
    ($\ell \geq 2$) modes to exist uniquely selects three
    spatial dimensions (\S\ref{sec:dimensions}).
\end{itemize}

\medskip
\noindent\textbf{Epistemic status of the above.}---The mass formula
reproduces the observed spectrum numerically, but several ingredients
are working hypotheses rather than first-principles derivations:
(i) the reduction from the full three-dimensional cavity problem of
the Quantum paper to the one-dimensional Mathieu equation used here
is not formally established and is left as an open calculation
(\S\ref{sec:q_derivation});
(ii) the numerical value of $\nu_0$ is a calibrated
MASS/Mathieu value, while its first-principles origin remains
open (\S\ref{sec:nu0_computation});
(iii) the identification of neutrinos as polarisation-rotation waves
(\S\ref{sec:neutrinos}) and of the $W^\pm/Z^0$ bosons as bridge
modes (\S\ref{sec:WZ_bosons}) is heuristic, with the Standard-Model
quantum numbers not derived here;
(iv) the three-dimensional selection argument relies on the
assumption that tensor ($\ell\geq 2$) modes are physically required,
which the main theory invokes but does not prove from the action
alone (\S\ref{sec:dimensions}).
Each of these items is flagged locally in the corresponding section;
their resolution is deferred to future work.

\section{Frequency Invariance and Derivation of Time Dilation}
\label{sec:freq_invariance}

\subsection{Starting Point in the Theory}

In the ISPG framework the bi-conformal metric is built from two
operational postulates (Appendix~3):
\begin{align}
  d\tau &= e^{\phi/2}\,dt
  \qquad\text{(clock postulate)},
  \label{eq:clock_here}\\[4pt]
  d\ell &= e^{-\phi/2}\,d\sigma
  \qquad\text{(rod postulate)},
  \label{eq:rod_here}
\end{align}
while the effective mass obeys:
\begin{equation}
  m_{\mathrm{eff}} = m_0\,e^{\phi/2}.
  \label{eq:mass_here}
\end{equation}
All three---time, length, mass---are manifestations of a single
conformal factor ($e^{\phi/2}$). Currently all three are treated
as \emph{parallel} consequences of a single metric.

\subsection{Logical Necessity: Why an Invariant Must Exist}

Before calling the invariant frequency an \emph{assumption}, let us
consider a more fundamental question: is ``time slowing down'' even
possible if \textbf{nothing in the universe remains unchanged}?

The answer is no.

``Slowing down'' is a \emph{comparative} notion. To say ``this clock
runs slowly'' one needs a reference---something that \emph{did not}
slow down, against which to compare. If everything changes equally,
no one notices anything. This is the deep content of the local
undetectability theorem.

Therefore time dilation logically requires \textbf{two elements}:
\begin{enumerate}
  \item \textbf{A variable element}---something that changes with
    the pressure state ($\phi$). In ISPG this is the oscillon's
    \emph{size} and \emph{mass}.
  \item \textbf{An invariant element}---something that \emph{never}
    changes, relative to which the variable element's change becomes
    perceptible.
\end{enumerate}

What is this invariant element? Consider the candidates:
\begin{itemize}
  \item Size? No, it varies: $d\ell = e^{-\phi/2}\,d\sigma$.
  \item Mass? No, it varies:
    $m_{\mathrm{eff}} = m_0\,e^{\phi/2}$.
  \item Clock rate? No, that is precisely what we seek to explain.
  \item Local speed of light? It is invariant
    ($c_{\mathrm{local}} = c$), but this is a \emph{consequence},
    not a cause.
\end{itemize}

One candidate remains: the \textbf{fundamental vibration frequency}
of the medium, $\nu_0$. If it does not change, then variations in
size and mass become \emph{perceptible} relative to it---and this
perception is what we call ``time dilation.''

\subsection{Circle Analogy}

Imagine a circular track of fixed circumference $L$. Two runners
stand on it: one tall (stride length $\lambda_1$), the other
exactly half as tall (stride $\lambda_2 = \lambda_1/2$). Both run
at the same cadence---the same number of strides per second.

The tall runner completes the lap in $N_1 = L/\lambda_1$ strides.
The short one in $N_2 = L/\lambda_2 = 2N_1$ strides---twice as
many strides, hence twice as much time.

Now imagine each runner drags the hand of a clock. The tall runner
turns the hand fast---``time passes quickly.'' The short runner
turns it slowly---``time passes slowly.''

Three elements are present:
\begin{itemize}
  \item \textbf{Variable:} the runner's size (stride length)---this
    is $m_{\mathrm{eff}}$.
  \item \textbf{Invariant:} the track length $L$---this is the
    fundamental frequency $\nu_0$ (the ``cycle'' that is the same
    for everyone).
  \item \textbf{Result:} the clock rate---this is $d\tau/dt$.
\end{itemize}

If the track length also changed (proportionally to the runner's
size), both runners would finish the lap simultaneously---no
``slowing down'' would occur. \textbf{It is precisely the
invariance of the track (frequency invariance) that makes slowing
down possible.}

\subsection{Theorem: Invariance of the Fundamental Frequency}
\label{sec:freq_inv_theorem}

The logical necessity argument (\S\ref{sec:freq_invariance})
identified a fundamental carrier frequency $\nu_0$ as the invariant
element required by the very concept of time dilation.  The theorem
below proves the narrower statement needed here: once a localized
matter-sector resonator with proper-time frequency $\nu_0$ is given,
a static background $\phi_{\rm bg}$ does not change that proper-time
frequency.  The theorem does not derive the existence, uniqueness,
or numerical value of $\nu_0$; the value used in this paper is the
separate MASS/Mathieu calibration discussed in
\S\ref{sec:nu0_computation}.

\medskip
\textbf{Theorem (frequency invariance).}
\textit{The proper-time oscillation frequency of any localized
resonant excitation in $\Phi$ (oscillon, atomic clock, or other
matter-sector resonator) is position-independent: if such a
resonator oscillates at frequency $\nu_0$ in proper time at
$\phi = 0$, it oscillates at the same frequency $\nu_0$ (measured
in local proper time of an observer co-located with the resonator)
in any region with static background potential
$\phi_{\rm bg} \neq 0$.  The substrate itself is not assigned
its own clock; the statement is about matter-sector resonators
living in the substrate.}

\medskip
\textit{Proof.}---Consider a region where the static gravitational
background is approximately constant:
$\phi_{\rm bg} \approx \mathrm{const}$.  The bi-conformal metric
is
\begin{equation}
  ds^2 = -e^{\phi_{\rm bg}}c^2\,dt^2
  + e^{-\phi_{\rm bg}}\,d\sigma^2.
  \label{eq:metric_bg}
\end{equation}
Introduce locally Minkowski coordinates of an observer co-located
with the resonator at fixed $\phi_{\rm bg}$:
\begin{equation}
  \tau \equiv e^{\phi_{\rm bg}/2}\,t, \qquad
  \xi^i \equiv e^{-\phi_{\rm bg}/2}\,\sigma^i.
  \label{eq:proper_frame}
\end{equation}
In these coordinates the metric becomes Minkowski:
$ds^2 = -c^2 d\tau^2 + d\xi^2$.  The transformation is a
mathematical convenience for analysing the field equation of the
oscillon solution; it does not assign a hidden time variable to
the substrate.

The kinetic scalar
$X = -\tfrac{1}{2}\,g^{\mu\nu}\partial_\mu\phi\,
\partial_\nu\phi$
is invariant under coordinate transformations (it is a scalar).
Therefore the $G_2(X)$ field equation
\begin{equation}
  \nabla_\mu\!\left[
    \frac{\partial G_2}{\partial X}\,\nabla^\mu\phi
  \right] = 0
\end{equation}
takes, in the $(\tau,\xi)$ chart, exactly its flat-space form:
\begin{equation}
  \partial_\mu\!\left[
    \frac{\partial G_2}{\partial X}\,
    \eta^{\mu\nu}\partial_\nu\phi
  \right] = 0.
  \label{eq:eom_proper}
\end{equation}
Equation~\eqref{eq:eom_proper} contains no reference to
$\phi_{\rm bg}$: its coefficients are the coupling constants
$M_{\rm Pl}$, $\alpha$, $M$ appearing in $G_2(X)$.  A localized
oscillon therefore keeps the same flat-space proper-time
eigenfrequency $\omega_0 = 2\pi\nu_0$ when placed in a static
background.  This proves background independence of the local
proper-time frequency, not a first-principles computation of its
numerical value.  $\square$

\medskip
\textit{Coordinate frequency and gravitational redshift.}---In
coordinate time $t$, the oscillon appears to oscillate at frequency
\begin{equation}
  \nu_{\rm coord} = \nu_0\,e^{\phi_{\rm bg}/2},
  \label{eq:nu_coord}
\end{equation}
since $\tau = e^{\phi_{\rm bg}/2}\,t$.  An asymptotic observer
($\phi_{\rm bg} = 0$) receiving signals from a region with
$\phi_{\rm bg} < 0$ detects a lower frequency---this is the
gravitational redshift, here \emph{derived} from the resonator's
intrinsic frequency invariance rather than postulated.

\medskip
\textit{Epoch independence.}---For cosmological evolution,
$\phi_{\rm bg}(t)$ varies on timescales
$\sim H_0^{-1} \sim 10^{17}$~s, while a typical oscillon resonates
at $\nu_0 \sim 10^{30}$~Hz.  The adiabatic ratio
$\nu_0/H_0 \sim 10^{47}$ ensures that the local proper-time
frequency tracks its flat-space value to fractional accuracy
$\sim (H_0/\nu_0)^2 \sim 10^{-94}$.  Over the entire Hubble
time, the cumulative phase drift is
$\Delta\omega/\omega_0 \sim H_0/(2\pi\nu_0) \sim 10^{-48}$
---undetectable by any conceivable experiment.
This is an adiabatic extension of the static theorem, not an
independent proof for arbitrary dynamical backgrounds.

\medskip
\textit{Physical content.}---What changes is not the frequency but
the ``stride size'' of each vibration---the oscillon's effective
mass and spatial scale.  The frequency is the invariant ``track''
against which size changes are perceived as ``time dilation.''

\bigskip
\begin{center}
\fbox{\parbox{0.85\textwidth}{\centering\large
\textbf{Central thesis.}\quad
Time is the balance between oscillon size and the medium's
intrinsic frequency.  This frequency, measured in local proper
time, is universal---it is the same at every point in space and
is adiabatically stable over cosmological evolution.  The particle
mass spectrum is then modeled by the Mathieu resonance program
built on this carrier.
}}
\end{center}
\bigskip

\subsection{Quantum Origin: The ISPG Vacuum as a
  Relativistic Condensate}
\label{sec:condensate}

The local proof of \S\ref{sec:freq_inv_theorem}
established static-background invariance of a given proper-time
carrier frequency.  A deeper question remains: \emph{why}
does the medium have a single, universal carrier frequency at all?
The answer is that the ISPG vacuum is a macroscopic quantum
coherent state---a relativistic Bose--Einstein condensate
\cite{PethickSmith2008,Leggett2001} of
the archion field.

\medskip
\textbf{Quantization of the archion field.}
Promote the classical scalar $\phi$ to a quantum field
operator and decompose into Fourier modes in a large
comoving volume~$V$:
\begin{equation}
  \hat\phi(\mathbf{x},t)
  = \sum_{\mathbf{k}}
    \frac{1}{\sqrt{2\omega_k V}}
    \!\left[
      a_{\mathbf{k}}\,
      e^{\,i(\mathbf{k}\cdot\mathbf{x}-\omega_k t)}
      + a_{\mathbf{k}}^\dagger\,
      e^{-i(\mathbf{k}\cdot\mathbf{x}-\omega_k t)}
    \right],
  \label{eq:phi_quantized}
\end{equation}
with canonical commutation relation
$[a_{\mathbf{k}},\,a_{\mathbf{k}'}^\dagger]
= \delta_{\mathbf{k}\mathbf{k}'}$.

\medskip
\textbf{Condensate ground state.}
The ISPG vacuum is defined as a coherent state of the
zero-momentum ($\mathbf{k}=0$) mode:
\begin{equation}
  a_{\mathbf{0}}\,|\alpha_0\rangle
  = \alpha_0\,|\alpha_0\rangle,
  \qquad
  N_0 \equiv |\alpha_0|^2 \gg 1,
  \label{eq:coherent_state}
\end{equation}
where $N_0$ is the macroscopic occupation number.  All
$N_0$ quanta share the single-particle ground-state
energy $\varepsilon_0 = \hbar\omega_0$.  The expectation
value of the field operator is
\begin{equation}
  \langle\hat\phi\rangle
  = \frac{2\,|\alpha_0|}{\sqrt{2\omega_0 V}}\;
    \cos(\omega_0 t + \theta_0)
  \;\equiv\;
  \Phi_{\rm bg}\cos(\omega_0 t),
  \label{eq:mean_field}
\end{equation}
which is precisely the classical oscillating background
assumed in the oscillon ansatz~\eqref{eq:oscillon_ansatz}.
The classical field used throughout this work is the
condensate mean field.

\medskip
\textbf{Frequency as chemical potential.}
In a non-relativistic BEC with contact interaction
$g|\Psi|^2$, the chemical potential $\mu = g\,n_0$
determines the condensate frequency: $\omega = \mu/\hbar$.
In ISPG the analogous role is played by the nonlinear
kinetic self-interaction
$G_2(X) \supset \alpha X^2/M^2$: the $X^2$ term acts
as the self-interaction that sets the oscillation
frequency.  From the self-consistency
condition~\eqref{eq:amplitude_G2},
\begin{equation}
  \Phi_{\rm bg}^2\,\omega_0^2
  = \frac{2\,M^2\,M_{\rm Pl}^2}{3\,\alpha}\,,
  \label{eq:self_consistency_condensate}
\end{equation}
the product $\Phi_{\rm bg}^2\omega_0^2$ is fixed
entirely by the coupling constants $\alpha$, $M$,
$M_{\rm Pl}$.  Given the total number of quanta
$N_0 = \Phi_{\rm bg}^2\omega_0 V/(2\hbar)$ (from
equating the classical energy density
$\rho = \tfrac{1}{2}\Phi_{\rm bg}^2\omega_0^2$ with
$N_0\hbar\omega_0/V$), the frequency is
\begin{equation}
  \boxed{\;
  \omega_0
  = \frac{M^2\,M_{\rm Pl}^2\,V}
         {3\,\alpha\,\hbar\,N_0}
  \;\propto\; \frac{1}{n_0}\,,
  \;}
  \label{eq:omega0_condensate}
\end{equation}
where $n_0 = N_0/V$ is the archion number density.
This is a \emph{structural} analog of $\mu = g\,n_0$ in a
standard BEC: in both systems the oscillation frequency is
set by the condensate density and coupling constants---not
by position.  The functional dependence is, however,
\emph{inverted}: in a repulsive BEC the chemical potential
\emph{increases} with density ($\mu \propto n_0$), whereas
in ISPG the proper-time frequency \emph{decreases}
($\omega_0 \propto 1/n_0$) because the $X^2$
self-interaction raises the effective inertia of the
oscillating field at higher amplitude.  Since $\alpha$, $M$, $M_{\rm Pl}$ are
spacetime-constant parameters and $n_0$ is homogeneous on
scales where $\phi_{\rm bg} \approx \mathrm{const}$, the
frequency $\omega_0$ is the same at every point.

\medskip
\textbf{Why $\nu_0$ invariance is automatic.}
In a coherent state~\eqref{eq:coherent_state}, every
quantum carries the same energy $\hbar\omega_0$.
Spatial variation of the condensate (Bernoulli depletion
near a massive body) changes the \emph{local amplitude}
$\Phi_{\rm bg}(x)$ and the local number density $n_0(x)$,
but not the single-particle energy.  This is exact: in a
superfluid, a density dip (vortex, depletion layer) does
not shift the chemical potential---it is the chemical
potential that is globally uniform, while the density
adjusts.  In ISPG, the chemical potential is $\omega_0$:
\begin{equation}
  \nu_0 \;=\; \frac{\omega_0}{2\pi}
  \;=\; \text{const (in proper time)}
  \qquad\text{everywhere.}
  \label{eq:nu0_condensate}
\end{equation}
The coordinate-frame manifestation
$\nu_{\rm coord} = \nu_0\,e^{\phi_{\rm bg}/2}$
(Eq.~\ref{eq:nu_coord}) is the standard relation
between proper-time and coordinate frequencies.

\medskip
\textbf{Active-existence guardrail.}
The active-existence picture gives an operational reading of
condensate depletion and measurable-content redistribution: a local
change in active carrier content is seen by matter standards as a
change of measured scale, density, or mass.  It does not replace
the separate MASS/Mathieu calibration of $\nu_0$.  In particular,
active-existence does not derive $\nu_0$, does not derive
$1.91\times 10^{30}$~Hz, does not prove the mass ladder, and does
not close the Mathieu $q$ derivation.

\medskip
\textbf{Oscillons as condensate excitations.}
Particles (oscillons) are localized, nonlinear excitations
on top of the condensate background, analogous to vortices
in He-II, solitons in nonlinear optics, or rotons in
superfluid helium.  In the second-quantized language:
\begin{equation}
  \hat\phi(x,t) = \underbrace{
    \Phi_{\rm bg}\cos(\omega_0 t)
  }_{\text{condensate (vacuum)}}
  +\; \underbrace{
    \delta\hat\phi(x,t)
  }_{\text{excitations (oscillons, waves)}}.
  \label{eq:field_split}
\end{equation}
The condensate determines the arena---background pressure,
$c_{\rm coord}$, effective mass scaling.  The excitations
are what we call matter.  This split is the quantum analog
of the classical decomposition already used in
\S\ref{sec:q_derivation}.

\medskip
\textbf{Summary of the condensate identification.}
\begin{center}
\begin{tabular}{lll}
\hline
\textbf{Standard BEC} & \textbf{ISPG condensate}
  & \textbf{Physical meaning} \\
\hline
$\Psi(\mathbf{x},t)$
  & $\phi(\mathbf{x},t)$
  & Macroscopic order parameter \\
$\mu = g\,n_0$
  & $\hbar\omega_0 \propto 1/n_0$
  & Chemical potential = frequency \\
$|\Psi|^2 = n_0$
  & $\Phi_{\rm bg}^2\omega_0/(2\hbar V)$
  & Number density \\
$g|\Psi|^2$
  & $\alpha X^2/M^2$
  & Nonlinear self-interaction \\
Phonons, vortices
  & Oscillons, gravitational waves
  & Elementary excitations \\
Local $n$ varies
  & Local $\Phi_{\rm bg}(x)$ varies
  & Density adjusts; $\omega_0$ does not \\
\hline
\end{tabular}
\end{center}
The ``one archion everywhere'' of the intuitive picture
is, in formal language, the statement that $N_0$ quanta
occupy a single macroscopic quantum state with one
frequency.  This is neither a postulate nor a metaphor---it
is the standard quantum-field-theoretic description of a
condensate, realized here for the ISPG scalar sector.

\subsection{Practical Note: Searching for $\nu_0$}

This subsection gives a search heuristic, not an independent proof.
Every particle has a Compton frequency
$\nu_C = mc^2/h$ (de~Broglie, 1924; M\"uller, 2013).  If a single
carrier frequency underlies the MASS/Mathieu ladder, the observed
Compton frequencies should display a commensurate pattern when
expressed through the ladder's dimensionless indices.  The exact
integer/harmonic status of that pattern belongs to the Mathieu
analysis, not to the frequency-invariance theorem itself.

Concretely: rational or near-rational mass ratios are evidence for
a resonant organization of the spectrum and guide the calibration
of $\nu_0$ in the MASS/Mathieu model.  They do not by themselves
derive $\nu_0$ from first principles.

\subsection{Guitar String Analogy}

Recall a guitar string. If we loosen it (reduce tension), the same
fret produces a different note. To maintain the same note
(frequency) one must:
\begin{itemize}
  \item shorten the vibrating length (size changes),
  \item reduce the amplitude (mass/energy changes),
  \item but keep the frequency the same.
\end{itemize}

Analogously: when the medium's background pressure decreases, the
oscillon ``shrinks in size and mass'' in order to maintain the
fundamental frequency.

\subsection{Tall and Short Person Analogy}

Two people run at the same cadence (one frequency $\nu_0$). The
tall person has long legs (large $m_{\mathrm{eff}}$, high
pressure)---covers a large distance per stride. The short person
has short legs (small $m_{\mathrm{eff}}$, low pressure)---covers
less distance at the same cadence.

``Time rate'' $=$ cadence $\times$ stride size. The cadence
($\nu_0$) is the same. The stride ($m_{\mathrm{eff}}$) differs.
Therefore the rate differs.

\subsection{Derivation of the Clock Postulate}

\textbf{Given:}
\begin{enumerate}
  \item Rod postulate: $d\ell = e^{-\phi/2}\,d\sigma$.
  \item Effective mass scaling:
    $m_{\mathrm{eff}} = m_0\,e^{\phi/2}$
    (follows from the bi-conformal metric's matter coupling).
  \item Frequency invariance
    (Theorem~\ref{sec:freq_inv_theorem}): proper-time
    frequency $\nu_0 = \mathrm{const}$; coordinate
    frequency $\nu_{\rm coord} = \nu_0\,e^{\phi/2}$.
\end{enumerate}

\textbf{To derive:} Clock postulate
$d\tau = e^{\phi/2}\,dt$.

\medskip
\textbf{Derivation.}

An oscillon-based clock counts oscillations.  In coordinate
time $dt$, the number of oscillations is
$N = \nu_{\rm coord}\,dt = \nu_0\,e^{\phi/2}\,dt$
(Eq.~\ref{eq:nu_coord}).  Each oscillation takes proper time
$1/\nu_0$ (since $\nu_0$ is the proper-time frequency).
Therefore:
\begin{equation}
  d\tau = N \times \frac{1}{\nu_0}
  = \nu_0\,e^{\phi/2}\,dt \times \frac{1}{\nu_0}
  = e^{\phi/2}\,dt.
  \label{eq:clock_derived}
\end{equation}

Hence:
\begin{equation}
  \boxed{d\tau = e^{\phi/2}\,dt}
  \label{eq:clock_from_freq}
\end{equation}
which is exactly the clock postulate \eqref{eq:clock_here}.
$\square$

\medskip
\textbf{Note on circularity.}
The coordinate frequency $\nu_{\rm coord} = \nu_0\,e^{\phi/2}$
(Eq.~\ref{eq:nu_coord}) was derived using the bi-conformal metric,
in which the clock postulate is already embedded.  This derivation
is therefore a self-consistency check rather than an independent
proof.

This problem is fully resolved in
Section~\ref{sec:two_plate_derivation}, where mass scaling is
\textbf{independently} derived from two-plate collision physics.

\subsection{Derivation of the Local Speed of Light}

The locally measured speed of light:
\begin{equation}
  c_{\mathrm{local}}
  = \frac{d\ell}{d\tau}
  = \frac{e^{-\phi/2}\,d\sigma}{e^{\phi/2}\,dt}
  = e^{-\phi}\,\frac{d\sigma}{dt}.
  \label{eq:clocal_step1}
\end{equation}
The coordinate speed of light in the bi-conformal metric
($ds^2 = 0$):
\begin{equation}
  \frac{d\sigma}{dt} = c\,e^{\phi}.
  \label{eq:ccoord_here}
\end{equation}
Substituting \eqref{eq:ccoord_here} into
\eqref{eq:clocal_step1}:
\begin{equation}
  c_{\mathrm{local}} = e^{-\phi}\cdot c\,e^{\phi} = c.
  \label{eq:clocal_result}
\end{equation}

\begin{equation}
  \boxed{c_{\mathrm{local}} = c}
  \label{eq:c_invariant}
\end{equation}

The locally measured speed of light is invariant once the rod and
clock scalings are used.  The clock scaling's status is the
self-consistency check discussed above and qualified below. $\square$

\subsection{Explanation of Redshift}

Light is emitted at pressure state $\phi_{\mathrm{emit}}$ and
received at $\phi_{\mathrm{recv}}$.

The fundamental frequency $\nu_0$ does not change. But a local
observer measures frequency with their own clock, whose ``stride
size'' is proportional to $e^{\phi/2}$. Therefore:

\begin{equation}
  \frac{\nu_{\mathrm{obs}}}{\nu_{\mathrm{emit}}}
  = \frac{e^{\phi_{\mathrm{emit}}/2}}
         {e^{\phi_{\mathrm{recv}}/2}}
  = e^{(\phi_{\mathrm{emit}}-\phi_{\mathrm{recv}})/2}.
  \label{eq:redshift_freq}
\end{equation}

\textbf{Gravitational redshift:}
Light is emitted at low pressure (in a gravitational well,
$\phi_{\mathrm{emit}} < \phi_{\mathrm{recv}}$). Then
$\nu_{\mathrm{obs}} < \nu_{\mathrm{emit}}$---redshift.

Light ``loses'' nothing. The fundamental oscillation is the same.
What changed is the \emph{measuring ruler's size}: at the emission
point the clock ``stride'' was small (small $m_{\mathrm{eff}}$),
at the reception point it is large. We perceive this difference as
``redshift.''

\textbf{Cosmological redshift:}
In the past, background pressure was high---oscillons were
``large,'' the clock stride was large. Today background pressure
has decreased---oscillons are ``small,'' the clock stride is small.
Distant light was emitted in the ``large stride'' epoch and arrives
in the ``small stride'' epoch---we measure a large-ruler result
with a small ruler, and we see redshift.

\subsection{Summary: One Postulate Instead of Two}

If the scoped proper-time frequency invariance above is used,
then:
\begin{center}
\renewcommand{\arraystretch}{1.3}
\begin{tabular}{@{}ll@{}}
\textbf{Given} & Rod postulate + mass scaling
  + $\nu_0 = \mathrm{const}$ \\
\textbf{Self-consistency check} & Clock postulate:
  $d\tau = e^{\phi/2}\,dt$ \\
\textbf{Consequence} & $c_{\mathrm{local}} = c$ once the rod
  and clock scalings are used \\
\textbf{Consequence} & Gravitational and cosmological
  redshift \\
\end{tabular}
\end{center}

Instead of treating the clock law as an unrelated operational
postulate, this section records its oscillon self-consistency route
from the rod scaling and frequency invariance.  The caveat below
spells out the status: this is a self-consistency check, not an
independent kinematic derivation from frequency data alone.

\textbf{Time does not ``slow down.'' The stride shrinks. The
rhythm is the same.}

\subsection{Status}

The derivation \eqref{eq:clock_derived} above relies on the
interpretation that proper-time accumulation is proportional to the
oscillation's energetic ``weight.''

\textbf{Caveat (self-consistency, not independent derivation).}
Step (3) of the derivation uses
$\nu_{\rm coord} = \nu_0\,e^{\phi/2}$, which itself follows
from the metric/clock law of the bi-conformal background.
Equation~\eqref{eq:clock_derived} should therefore be read as
a \emph{self-consistency check} of the clock postulate within
the oscillon picture, not as an independent derivation of
$d\tau = e^{\phi/2}\,dt$ from kinematics alone.

\textbf{Update.}
The two-plate model (\S\ref{sec:two_plate_derivation})
\textbf{confirmed} this interpretation:
$m_{\mathrm{eff}} = m_0 e^{\phi/2}$ follows from collision
kinematics (not analogical reasoning), and the clock postulate
$d\tau = e^{\phi/2}\,dt$ follows logically from it. The numerical
Mathieu solution (\S\ref{sec:mathieu_numerical}) provided
quantitative confirmation: muon mass to $< 0.001\%$ accuracy.

\section{Harmonic Resonance: The Universe as a Chord}
\label{sec:harmonic}

\subsection{From One Frequency to Many}

In the previous section a single proper-time carrier frequency
$\nu_0$ was identified as invariant under static backgrounds.  This
does not by itself determine the particle mass ladder.  Nature
contains different particles with different masses---electrons,
muons, quarks, bosons. If mass and frequency are related
($E = mc^2 = h\nu$), different masses imply different Compton
frequencies. How is this compatible with a single carrier
frequency?

The answer comes from the same guitar string considered above, but
now examined more deeply.

\subsection{Guitar String Harmonics}

When a guitar string vibrates, it does not produce only one note.
Alongside the fundamental, an entire spectrum of
\emph{harmonics}---accompanying frequencies that are integer
multiples of the fundamental---sounds:
\begin{equation}
  \nu_n = n\,\nu_{\mathrm{fund}}, \qquad n = 1, 2, 3, \ldots
  \label{eq:harmonics}
\end{equation}

But string physics is richer. Pressing the string at different
points yields \emph{musical intervals}:
\begin{itemize}
  \item At the midpoint---\textbf{octave}: ratio $2{:}1$.
  \item At $2/3$---\textbf{fifth}: ratio $3{:}2$.
  \item At $3/4$---\textbf{fourth}: ratio $4{:}3$.
  \item At $4/5$---\textbf{major third}: ratio $5{:}4$.
\end{itemize}
All these intervals are \emph{compatible} with the fundamental: they
are rational fractions, i.e.\ every $n$-th vibration coincides
exactly with the fundamental cycle.

This is precisely why these intervals sound ``harmonious''---there
is no dissonance between them because they are \emph{commensurate
sub-modes} of the same fundamental vibration.

\subsection{The Medium as a String---Oscillons as Harmonics}

The ISPG medium is far richer than a one-dimensional string---it is
three-dimensional, has pressure, compressibility, and topological
structure (vortices). But the fundamental principle is the same: the
medium can support only certain resonant modes---those oscillations
that ``fit'' local conditions and do not vanish through destructive
interference.

This means that alongside the fundamental $\nu_0$, the medium may
host \emph{harmonically compatible} resonances:
\begin{equation}
  \nu_{n,m} = \frac{n}{m}\,\nu_0,
  \qquad n, m \in \mathbb{Z}^{+},
  \label{eq:rational_freq}
\end{equation}
where $n/m$ is a rational fraction, analogous to a musical
interval.

Different particles would be different ``notes'' in this spectrum:
\begin{itemize}
  \item One oscillon ``lives'' on every fundamental click---the
    fundamental mode.
  \item Another responds every \emph{second} click---half
    frequency, twice the mass.
  \item A third responds every \emph{third} click.
  \item Some respond every 2 out of 3 clicks (the analogue of a
    fifth), or 3 out of 4 (a fourth).
\end{itemize}

All modes ``live'' in the same fundamental rhythm, but each
``listens'' to it differently---one on every beat, another on every
second, a third on every third.

\subsection{Two Plates and Bouncing Balls}

The guitar-string analogy has one weakness: a string possesses
\emph{all} harmonics simultaneously. In nature, however, only
specific particles with specific masses exist. What \emph{selection
rule} determines which harmonics are ``allowed'' and which are not?

A different, more precise analogy answers this question.

\medskip
\textbf{Picture: two planes of a wave.}
The medium's fundamental vibration $\nu_0$ creates a standing wave.
This wave has two extrema: the \textbf{crest} (upper plate) and the
\textbf{trough} (lower plate). The two plates stand at a fixed
``distance'' $L$---the space defined by the wave amplitude.

A particle (oscillon) bounces \emph{between} these two plates: the
lower plate throws it upward, it crosses the distance, the upper
plate ``throws'' it back down, it returns, the lower plate throws
it up again---and so on indefinitely.

The particle \textbf{never rests}. It is in perpetual motion between
the plates---corresponding to the fact that real elementary
particles never ``stop'': their internal oscillation (Compton
frequency, Zitterbewegung) is constant.

\medskip
\textbf{Resonance condition.}
A ball bouncing between two plates is stable only when its
round-trip time (bottom to top $+$ top to bottom $=$ one full
cycle) \emph{exactly} fits an integer multiple of the plate
vibration period:
\begin{equation}
  t_{\mathrm{cycle}} = t_{\uparrow} + t_{\downarrow}
  = \frac{n}{\nu_0},
  \qquad n = 1, 2, 3, \ldots
  \label{eq:flight_time}
\end{equation}

If the ball reaches a plate at the ``wrong'' moment---the plate is
already moving the other way---the ball loses energy, and after a
few cycles the bouncing ceases. This is an \textbf{unstable mode}
that ``dies.''

\medskip
\textbf{Mass $=$ collision intensity.}
Different balls cross the distance between the two plates at
different speeds:
\begin{itemize}
  \item A \textbf{fast} ball collides frequently with the plates
    $\Rightarrow$ large impulse per collision $\Rightarrow$ strong
    interaction with the external world $\Rightarrow$
    \textbf{large mass}.
  \item A \textbf{slow} ball collides rarely $\Rightarrow$ small
    impulse $\Rightarrow$ weak interaction $\Rightarrow$
    \textbf{small mass}.
\end{itemize}
In this picture \textbf{mass is not a property of a ``thing''---mass
is the intensity of the bouncing}. It is the rate of impulse
transfer to the plates. A light particle (electron) crosses slowly,
collides rarely. A heavy particle (top quark) crosses fast, collides
often.

Meanwhile \emph{all} plates vibrate at the same frequency $\nu_0$.
One ball keeps up with every vibration; another needs several
vibrations. But all share the same rhythm.

\medskip
\textbf{Energy conservation: the necessity of balance.}
The upper plate ``throws'' the ball downward with exactly the same
force the lower plate threw it upward. The ball cannot ``stick'' to
the upper plate---it necessarily returns. This symmetry is the
energy conservation law in this picture.

If the lower impulse $\neq$ the upper impulse, the ball gradually
drifts to one side and bouncing ceases. Only \textbf{perfect
balance} is stable.

\medskip
\textbf{Different balls---different sub-harmonics.}
\begin{itemize}
  \item Fast ball: completes a full cycle on every vibration
    ($n = 1$, frequency $\nu_0$, maximum mass in this class).
  \item Medium ball: completes a cycle every \emph{second}
    vibration ($n = 2$, frequency $\nu_0/2$).
  \item Slow ball: every \emph{third} vibration
    ($n = 3$, frequency $\nu_0/3$).
  \item Very slow: every $n$-th vibration.
\end{itemize}

\medskip
\textbf{Complex mass ratios.}
The ball's transit time depends not only on its size but also on
elasticity, the character of plate impulses, and wave amplitude.
This complex dependence naturally produces complicated mass ratios
(e.g.\ $m_\mu/m_e \approx 206.8$---not a simple 2 or 3), yet
\emph{all} are still determined by a single frequency between two
plates.

\medskip
\textbf{Physical precedents.}
The two-plate bouncing picture has direct parallels in known
physics:
\begin{enumerate}
  \item \textbf{Quantum particle in a box.} A particle between two
    walls---the most fundamental model in quantum mechanics. Energy
    levels $E_n = n^2\pi^2\hbar^2/(2mL^2)$ are discrete and
    determined by the box size $L$.
  \item \textbf{Zitterbewegung.} Schr\"odinger (1930) showed that
    the Dirac electron oscillates between positive and negative
    energy states---effectively ``bouncing'' between matter and
    antimatter at frequency $\omega = 2mc^2/\hbar$. The factor 2
    arises because one cycle involves \emph{two} crossings: up and
    down.
  \item \textbf{Casimir effect.} Between two metal plates, vacuum
    fluctuations produce a measurable force. Only certain
    wavelengths ``fit'' between the plates---the rest are
    forbidden. This is a natural selection rule between two
    boundaries.
\end{enumerate}

\medskip
\textbf{Experimental confirmation: Couder's bouncing droplets.}
In 2005 Yves Couder and Emmanuel Fort (Paris) demonstrated an
experiment in which an oil droplet bounces on a vibrating oil bath.
The droplet bounces on every second vibration (sub-harmonic
$\nu_0/2$) and each bounce creates a surface wave---a
\emph{pilot wave}.

This classical system \textbf{reproduced quantum-mechanical
phenomena}: double-slit interference, quantized orbits, tunneling,
and spin analogues, as reported in the experimental literature
in 2005--2006.

\medskip
\textbf{ISPG analogue.}
\begin{center}
\renewcommand{\arraystretch}{1.3}
\begin{tabular}{@{}lll@{}}
\textbf{Two-plate model} & $\longleftrightarrow$
  & \textbf{ISPG} \\
Upper/lower plate & & Wave crest/trough in medium \\
Plate frequency $\nu_0$ & & Fundamental frequency $\nu_0$ \\
Bouncing ball & & Oscillon (particle) \\
Transit speed & & Effective mass $m_{\mathrm{eff}}$ \\
Collision impulse & & Inertia / mass \\
Stable bouncing & & Stable particle \\
Unstable bouncing & & Non-existent particle \\
Plate balance & & Energy conservation \\
\end{tabular}
\end{center}

\medskip
\textbf{Mathematical formalization.}
The dynamics of a particle oscillating between two boundaries is
described by the \textbf{Mathieu equation} \cite{McLachlan1947,DLMFMathieu}:
\begin{equation}
  \frac{d^2 x}{dt^2}
  + \bigl(a - 2q\cos(2\nu_0 t)\bigr)\,x = 0,
  \label{eq:mathieu}
\end{equation}
where $x$ is the particle's position between the two plates, and
$a$, $q$ depend on the particle's properties and plate vibration
amplitude.

The Mathieu equation possesses \textbf{stability zones}---specific
$(a, q)$ combinations where the solution is periodic (the ball
bounces stably), and \textbf{instability zones} where the solution
grows exponentially or decays (the ball ``dies'').

Stability zones $=$ allowed harmonics $=$ existing particles.
Instability zones $=$ forbidden modes $=$ non-existent particles.

If the oscillon wave equation in the ISPG medium yields a
Mathieu-type structure, the particle mass spectrum follows
automatically from the stability zones.

\subsection{What Happens When the Plates Move Closer}

Now let us pose the decisive question: if the plates move closer
($L \to L' < L$) and the ball \emph{proportionally} shrinks, what
changes?

\medskip
\textbf{From the ball's perspective: nothing.}
The ball's size and its ``habitat'' (distance between plates) shrank
equally. The ball cannot tell it is smaller---its ruler is smaller
too. This is \textbf{local undetectability}: local physics is
unchanged.

\medskip
\textbf{From an external observer's perspective: everything.}
The external observer's plates are unchanged. They see:
\begin{enumerate}
  \item \textbf{Rhythm is the same}---$\nu_0$ did not change. The
    ball still completes a cycle every $n$-th vibration.
  \item \textbf{Distance is less}---plates are closer; each jump
    covers \emph{less} distance.
  \item \textbf{Speed is less}---less distance in the same time
    $=$ less speed.
  \item \textbf{Impulse is less}---slow collision $=$ less impulse
    transfer $=$ \textbf{less mass}.
  \item \textbf{Time rate is less}---same rhythm $\times$ smaller
    stride $=$ \textbf{slower clock}.
\end{enumerate}

\medskip
\textbf{This is time dilation.}
Denote the plate-separation reduction factor by $\alpha$:
$L' = \alpha L$, where $\alpha < 1$. Then:
\begin{itemize}
  \item Stride length: $\lambda' = \alpha\,\lambda$
    \quad(rod postulate),
  \item Speed: $v' = L'/t = \alpha\,v$,
  \item Collision impulse: $p' = \alpha\,p$
    $\;\Rightarrow\;$ $m' = \alpha\,m$
    \quad(mass scaling),
  \item Time rate: $\nu_0 \times \lambda'
    = \alpha\,(\nu_0 \times \lambda)$
    $\;\Rightarrow\;$ $d\tau' = \alpha\,d\tau$
    \quad(clock postulate).
\end{itemize}

In ISPG, $\alpha = e^{\phi/2}$. Everything follows from a single
factor---the plate separation.

\medskip
\textbf{Resolving the circularity.}
This picture resolves the circularity problem described earlier.
When deriving the clock postulate we used
$m_{\mathrm{eff}} = m_0\,e^{\phi/2}$, which was itself derived
from the clock postulate---a circle.

In the two-plate model, mass scaling \textbf{does not require the
clock postulate}. It follows directly from collision physics:
\begin{quote}
Plates closer $\;\Rightarrow\;$ less distance $\;\Rightarrow\;$
less speed $\;\Rightarrow\;$ less impulse $\;\Rightarrow\;$ less
mass.
\end{quote}
This is an \emph{independent} causal chain that does not require
the clock postulate. Mass scaling becomes a \emph{consequence}, not
an \emph{assumption}.

Then: mass $+$ invariant frequency $\Rightarrow$ clock
postulate---and this chain is no longer circular.

\medskip
\textbf{Complete causal chain.}
\begin{center}
\fbox{\parbox{0.85\textwidth}{
$\nu_0 = \mathrm{const}$\quad(fundamental frequency)\\[4pt]
$+$\quad Plate separation $L(\phi)$\quad(rod postulate)\\[4pt]
$\Downarrow$\\[2pt]
Less $L$ $\;\Rightarrow\;$ less impulse
  $\;\Rightarrow\;$ $m_{\mathrm{eff}}
  = m_0\,e^{\phi/2}$\quad(mass scaling)\\[4pt]
$\Downarrow$\\[2pt]
$\nu_0 \times m_{\mathrm{eff}}$ $\;\Rightarrow\;$
  $d\tau = e^{\phi/2}\,dt$\quad(clock postulate)\\[4pt]
$\Downarrow$\\[2pt]
$d\ell/d\tau = c$\quad(light-speed invariance)
}}
\end{center}

Everything flows from two elements: invariant frequency $+$ plate
separation. Everything else---mass, clock, speed of light---is a
\emph{consequence}.

\subsection{Formal Derivation from the Two-Plate Model}
\label{sec:two_plate_derivation}

\textbf{Given:}
\begin{enumerate}
  \item Fundamental frequency: $\nu_0 = \mathrm{const}$.
  \item Rod postulate: the oscillon's local size $\ell_0$ is
    unchanged; its coordinate size (plate separation) is
    \begin{equation}
      L(\phi) = \ell_0\,e^{\phi/2}.
      \label{eq:L_phi}
    \end{equation}
    (From $d\ell = e^{-\phi/2}\,d\sigma$ inversion:
    $d\sigma = e^{\phi/2}\,d\ell$.)
    In the active-existence reading we may also write
    $L_{\rm std}(\phi)\equiv L(\phi)$ and
    \begin{equation}
      \frac{\mathcal N_{\rm act}(\phi)}
           {\mathcal N_{\rm act}(0)}
      =
      \frac{L_{\rm std}(\phi)}{L_{\rm std}(0)}
      = e^{\phi/2}.
      \label{eq:Nact_two_plate}
    \end{equation}
    This is an operational definition of $\mathcal N_{\rm act}$ through
    the two-plate matter-standard ratio, not an independent number-density
    count of medium nodes and not a new observable.
  \item Resonance condition: the ball completes a full cycle (up $+$
    down) every $n/\nu_0$ coordinate seconds, where
    $n \in \mathbb{Z}^{+}$ is a fixed harmonic number.
\end{enumerate}

\textbf{To derive:}
$m_{\mathrm{eff}}(\phi) = m_0\,e^{\phi/2}$
and $d\tau = e^{\phi/2}\,dt$.

\bigskip
\textbf{Step 1. Coordinate velocity.}

In one cycle the ball covers $2L(\phi)$ coordinate distance
(up $L$ $+$ down $L$) in coordinate time $T = n/\nu_0$:
\begin{equation}
  v(\phi) = \frac{2L(\phi)}{n/\nu_0}
  = \frac{2\ell_0\,\nu_0}{n}\,e^{\phi/2}
  \equiv v_0\,e^{\phi/2},
  \label{eq:v_phi}
\end{equation}
where $v_0 \equiv 2\ell_0\nu_0/n$.

\bigskip
\textbf{Step 2. Impulse rate on the plates (effective mass
  derivation).}

Each collision transfers impulse $\Delta p = 2\mu\,v(\phi)$ to a
plate, where $\mu$ is the ball's internal inertia---an oscillon
``core'' parameter determined by \emph{local} physics and therefore
$\phi$-independent (local undetectability).

One cycle involves exactly 2 collisions (upper $+$ lower plate).
Collision frequency: $f_{\mathrm{coll}} = 2\nu_0/n$.

Average impulse transfer rate (force):
\begin{equation}
  F(\phi) = \Delta p \times f_{\mathrm{coll}}
  = 2\mu\,v(\phi) \times \frac{2\nu_0}{n}
  = \frac{4\mu\,\nu_0}{n}\,v(\phi).
  \label{eq:F_phi}
\end{equation}
Substituting $v(\phi)$ from \eqref{eq:v_phi}:
\begin{equation}
  F(\phi) = \frac{8\mu\,\ell_0\,\nu_0^2}{n^2}
  \,e^{\phi/2}
  \equiv F_0\,e^{\phi/2}.
  \label{eq:F_scaling}
\end{equation}

The effective mass is what the external world ``feels''---the
impulse rate.  Equivalently, it measures the same operational
active-existence content that controls the effective matter-standard
extension in \eqref{eq:Nact_two_plate}.  By normalization:
\begin{equation}
  \boxed{m_{\mathrm{eff}}(\phi) =
  m_0\,e^{\phi/2}}
  \label{eq:mass_from_plates}
\end{equation}
where $m_0 \equiv F_0 \times (\text{normalization factor})$.

\textbf{The clock postulate was not used anywhere in this
derivation.} Mass scaling follows solely from:
(a)~the coordinate plate separation \eqref{eq:L_phi} and
(b)~collision kinematics.

\bigskip
\textbf{Step 3. Derivation of the clock postulate.}

An oscillon-based clock counts oscillations at frequency $\nu_0/n$.
Each oscillation's ``weight'' in proper-time accumulation is
proportional to the oscillation energy:
$E(\phi) = m_{\mathrm{eff}}(\phi)\,c^2$.

Proper-time accumulation rate:
\begin{equation}
  \frac{d\tau}{dt}
  = \frac{E(\phi)}{E(0)}
  = \frac{m_{\mathrm{eff}}(\phi)}{m_{\mathrm{eff}}(0)}
  = e^{\phi/2}.
  \label{eq:clock_from_plates}
\end{equation}
\begin{equation}
  \boxed{d\tau = e^{\phi/2}\,dt}
  \label{eq:clock_final}
\end{equation}

\bigskip
\textbf{Step 4. Local speed of light.}

Locally measured speed:
\begin{equation}
  c_{\mathrm{local}}
  = \frac{d\ell}{d\tau}
  = \frac{e^{-\phi/2}\,d\sigma}{e^{\phi/2}\,dt}
  = e^{-\phi}\,\frac{d\sigma}{dt}.
\end{equation}
Coordinate speed of light in the bi-conformal metric:
$d\sigma/dt = c\,e^{\phi}$. Therefore:
\begin{equation}
  \boxed{c_{\mathrm{local}}
  = e^{-\phi}\cdot c\,e^{\phi} = c}
  \label{eq:c_final}
\end{equation}

\bigskip
\textbf{Step 5. Local undetectability check.}

Local plate separation:
\begin{equation}
  \ell_0 = e^{-\phi/2}\cdot L(\phi)
  = e^{-\phi/2}\cdot \ell_0\,e^{\phi/2} = \ell_0.
\end{equation}

\textbf{(a) Mass-ratio invariance.}
Consider two oscillons---harmonics $n_1$ and $n_2$. From Step~2:
\begin{equation}
  \frac{m_1(\phi)}{m_2(\phi)}
  = \frac{F_1(\phi)}{F_2(\phi)}
  = \frac{v_1(\phi)\,f_1}{v_2(\phi)\,f_2}
  = \frac{(2L\nu_0/n_1)\cdot(2\nu_0/n_1)}
         {(2L\nu_0/n_2)\cdot(2\nu_0/n_2)}
  = \frac{n_2^2}{n_1^2}.
  \label{eq:mass_ratio}
\end{equation}
$L(\phi)$ and $\nu_0$ cancel! The ratio is a function of integers
only and is $\phi$-invariant.

\textbf{(b) Compton-ratio invariance.}
In the bi-conformal frame, dimensionful constants scale as:
\begin{alignat}{2}
  m_{\mathrm{eff}} &= m_0\,e^{\phi/2}
    &\qquad& [\text{mass}],\\
  \hbar_{\mathrm{eff}} &= \hbar_0\,e^{-\phi}
    &\qquad& [\text{action} = ML^2T^{-1}],\\
  \nu_{\mathrm{proper}} &= (\nu_0/n)\,e^{-\phi/2}
    &\qquad& [\text{proper frequency}].
\end{alignat}
Checking the Compton ratio:
\begin{equation}
  \hbar_{\mathrm{eff}}\,(2\pi\nu_{\mathrm{proper}})
  = \hbar_0\,e^{-\phi}\cdot
  \frac{2\pi\nu_0}{n}\,e^{-\phi/2}
  = \frac{2\pi\hbar_0\nu_0}{n}\,e^{-3\phi/2}.
\end{equation}
On the other hand:
\begin{equation}
  m_{\mathrm{eff}}\,c^2
  = m_0\,c^2\,e^{\phi/2}.
\end{equation}
The local observer measures a \emph{dimensionless} quantity:
\begin{equation}
  \frac{\hbar_{\mathrm{eff}}\,\omega_{\mathrm{proper}}}
       {m_{\mathrm{eff}}\,c^2}
  = \frac{2\pi\hbar_0\nu_0}{n\,m_0\,c^2}\,
  e^{-2\phi}.
\end{equation}
At first glance this depends on $\phi$.

\textbf{Resolution:} $\nu_0$ is a \emph{coordinate} frequency. A
local observer cannot measure this number---they measure only
\emph{ratios} among local quantities. All local frequencies (of all
oscillons and photons) scale identically as $e^{-\phi/2}$ in proper
time, so their \emph{ratios} are $\phi$-invariant.

\textbf{(c) Summary.} A local observer cannot distinguish the value
of $\phi$ by any experiment, because:
\begin{itemize}
  \item Plate separation: $\ell_0$ (unchanged),
  \item Mass ratios: $n_2^2/n_1^2$ (integers),
  \item Speed of light: $c$ (unchanged),
  \item All frequencies co-scale---mutual ratios are fixed.
\end{itemize}
This is the local undetectability theorem in the two-plate model.
$\square$

\bigskip
\textbf{Summary.}
\begin{center}
\renewcommand{\arraystretch}{1.4}
\begin{tabular}{@{}lll@{}}
\textbf{Quantity} & \textbf{Formula} & \textbf{Source}
  \\[2pt]\hline\\[-6pt]
Plate separation
  & $L(\phi) = \ell_0\,e^{\phi/2}$
  & Rod postulate \\
Ball velocity
  & $v(\phi) = v_0\,e^{\phi/2}$
  & $L$ + resonance \\
Impulse rate
  & $F(\phi) = F_0\,e^{\phi/2}$
  & Collision physics \\
Effective mass
  & $m_{\mathrm{eff}} = m_0\,e^{\phi/2}$
  & $\equiv F/(\text{norm.})$ \\
Active measurable existence
  & $\mathcal N_{\rm act}(\phi)/\mathcal N_{\rm act}(0)
     = L_{\rm std}(\phi)/L_{\rm std}(0)=e^{\phi/2}$
  & Operational reading of $L(\phi)$ \\
Clock rate
  & $d\tau/dt = e^{\phi/2}$
  & $\nu_0 \times m_{\mathrm{eff}}$ \\
Speed of light
  & $c_{\mathrm{local}} = c$
  & $d\ell/d\tau$ \\
\end{tabular}
\end{center}

\textbf{One assumption} (rod) $+$ \textbf{one principle}
($\nu_0 = \mathrm{const}$) $+$ \textbf{collision kinematics}
$\Rightarrow$ the entire bi-conformal structure.

\subsection{Why This Is Physically Non-Trivial}

If particle masses are harmonics of a fundamental frequency, then
the mass spectrum is determined not by accident but by the medium's
resonant structure. This potentially answers one of the deepest
unanswered questions in physics: \textbf{why do particles have the
masses they have?}

In the Standard Model, particle masses (electron, muon, quarks, $W$
and $Z$ bosons) are 20+ free parameters inserted ``by hand'' from
experiment---the theory does not explain their origin.

If these masses are \emph{different harmonics of a single
fundamental frequency}, then 20+ parameters may follow from a
single principle---just as all harmonics of a guitar string follow
from one fundamental note.

\subsection{Precedents in Physics}

The harmonic resonance principle is not new in physics---it has been
discovered repeatedly at different levels:
\begin{enumerate}
  \item \textbf{Bohr atom (1913).} The electron in hydrogen occupies
    only ``allowed'' orbits---those fitting an integer number of
    wavelengths. Energy levels $E_n = E_0/n^2$ form a harmonic
    spectrum yielding hydrogen's entire spectral series (Balmer,
    Lyman, Paschen).
  \item \textbf{Nuclear shell model.} Nucleons organize into
    energetic ``shells,'' analogous to electronic orbitals. ``Magic
    numbers'' (2, 8, 20, 28, 50, 82, 126) mark especially stable
    configurations---resonant modes of the nuclear medium.
  \item \textbf{Koide formula \cite{Koide1982}.} Among the electron ($m_e$),
    muon ($m_\mu$), and tau lepton ($m_\tau$) masses a remarkably
    precise empirical relationship exists:
    \begin{equation}
      \frac{m_e + m_\mu + m_\tau}
           {\bigl(\sqrt{m_e} + \sqrt{m_\mu}
            + \sqrt{m_\tau}\bigr)^2}
      = \frac{2}{3},
      \label{eq:koide}
    \end{equation}
    holding to $0.01\%$ accuracy. The ratio $2/3$ is the inverse of
    the musical fifth ($3{:}2$). This formula remains unexplained
    within the Standard Model but naturally fits the harmonic
    resonance picture.
  \item \textbf{String theory.} Particles are interpreted as
    different vibrational modes of a single ``string.'' Particle
    mass $=$ frequency ($E = mc^2 = h\nu$). The entire mass
    spectrum is a spectrum of harmonics---exactly the same
    philosophy, except in ISPG the ``string'' is a scalar-metric
    \emph{medium}.
  \item \textbf{De~Broglie (1924).} Every massive particle has a
    wave of frequency $\nu = mc^2/h$. Matter vibrates. Matter $=$
    wave. This foundational quantum-mechanical discovery directly
    states: \textbf{mass is frequency.}
  \item \textbf{M\"uller's experiment (2013).} Holger M\"uller
    (UC Berkeley) experimentally demonstrated that a cesium atom's
    matter wave can serve as a clock---a ``Compton clock'' counting
    oscillations at frequency $mc^2/h$. His famous phrase: ``a rock
    is a clock.'' This experimentally confirms that every massive
    particle ``ticks'' at a fundamental frequency.
\end{enumerate}

\subsection{ISPG Context: Topology and Harmonics}

In the ISPG framework oscillons are already classified by
topological winding number $n$: fermions (vortex winding
$n = \pm 1$), opposite-winding antiparticle labels, and
bosons ($n = 0$).  The topological classification labels the
particle's \emph{topological class}. The harmonic resonance
principle complements this:
within a single topological class, different particles (e.g.\
electron, muon, and tau---all fermions with $n = 1$) may be
\emph{different harmonics} of the same topology:
\begin{equation}
  \nu_{\mathrm{electron}} : \nu_{\mathrm{muon}}
  : \nu_{\mathrm{tau}}
  = n_e : n_\mu : n_\tau,
  \label{eq:lepton_harmonics}
\end{equation}
where $n_e$, $n_\mu$, $n_\tau$ are integers or rationals determined
by the medium's resonant structure.

Topology $+$ harmonic together give a particle its full identity:
topology says \emph{what} it is (fermion or boson); the harmonic
says \emph{which} it is (electron, muon, or tau).

\subsection{Metaphor: The Universe as a Chord}

The universe is not a single note---it is a \textbf{chord}.

All notes come from one ``string'' (the medium). Different
particles are different notes in this chord. All notes are
harmonically compatible because all are multiples or rational
fractions of one fundamental rhythm.

Particle interactions are the chord's \emph{consonance}. Consonance
(stable binding) occurs when frequency ratios form simple rational
fractions. Dissonance (instability, decay) occurs when the ratio is
irrational or very complex.

For a musician this is intuitive: a fifth ($3{:}2$) is harmonious;
a tritone ($\sqrt{2}{:}1$) is dissonant. Perhaps nature also
``chooses'' which particle combinations are stable by the same
principle.

\subsection{Status}

\textbf{Program completed:}
\begin{enumerate}
  \item[\checkmark] The oscillon wave equation has been solved---
    particle masses: $m_N = m_1 b_N(q)$, where $b_N(q)$ is a
    Mathieu characteristic value
    (\S\ref{sec:numerical_results},
    \S\ref{sec:mathieu_numerical}).
  \item[\checkmark] The mass spectrum has been compared with
    experiment as a nearest-$N$ scan: in the current
    PDG-style table, $10/12$ entries lie below $1\%$ and
    $11/12$ below $5\%$; the muon anchor is exact at the
    fitted $q$, while the physical $N$-selection rule remains
    an open derivation target.
  \item[\checkmark] The Koide formula \eqref{eq:koide}
    \textbf{automatically} holds from the integer harmonics
    $(5, 72, 295)$ ($0.028\%$ accuracy).
\end{enumerate}

The ISPG framework not only describes particles' gravitational
behavior (which it already does) but \textbf{explains} their
\emph{existence}, their \emph{mass spectrum}, and the
\emph{electron's special status as the first fully stable charged
Mathieu mode in the adopted gap-threshold scheme}---
all from the resonant structure of a single medium.

%% ==========================================================
\section{Calibration and Historical Estimates of the Universal Frequency $\nu_0$}
\label{sec:nu0_computation}
%% ==========================================================

Goal: separate the historical classical estimate of $\nu_0$ from
the calibrated MASS/Mathieu value used in the quantitative mass
spectrum.  The frequency-invariance theorem fixes the background
independence of a given proper-time carrier frequency; it does not
by itself compute this numerical value.

\subsection{Bohr--Sommerfeld Condition in the Two-Plate Model}

In Section~\ref{sec:two_plate_derivation} mass scaling was derived
from classical collision kinematics. We now add a \textbf{quantum
condition}: the ball's action over one full cycle is quantized.

A ball (at $\phi = 0$ background) moves between the plates
(local separation $\ell_0$) at speed
$v = 2\ell_0\nu_0/r$ (where $r$ is the harmonic parameter). A
full cycle (up $+$ down) covers distance $2\ell_0$.

Bohr--Sommerfeld quantization:
\begin{equation}
  \oint p\,dx = N\,h, \qquad N = 1, 2, 3, \ldots
  \label{eq:bohr_sommerfeld}
\end{equation}
The ball's momentum is $p = \mu\,v$, where $\mu$ is its internal
inertia (a local parameter). Over one cycle:
$\oint p\,dx = \mu\,v \times 2\ell_0$. Hence:
\begin{equation}
  \mu \cdot \frac{2\ell_0\nu_0}{r}\cdot 2\ell_0
  = N\,h
  \qquad\Longrightarrow\qquad
  \nu_0 = \frac{N\,r\,h}{4\mu\,\ell_0^2}.
  \label{eq:nu0_from_BS}
\end{equation}

\subsection{Quantum Mass Formula}

The classical mass formula
(from Section~\ref{sec:two_plate_derivation}):
$m = 8\mu\,\ell_0^2\,\nu_0^2 / (r^2 c^2)$.
Substituting $\nu_0$ from \eqref{eq:nu0_from_BS}:
\begin{align}
  m &= \frac{8\mu\,\ell_0^2}{r^2\,c^2}
  \cdot\left(\frac{N\,r\,h}{4\mu\,\ell_0^2}\right)^2
  = \frac{8\mu\,\ell_0^2}{r^2\,c^2}
  \cdot\frac{N^2\,r^2\,h^2}{16\mu^2\,\ell_0^4}
  \notag\\
  &= \frac{N^2\,h^2}{2\mu\,\ell_0^2\,c^2}.
  \label{eq:mass_quantum_box}
\end{align}

This is the standard quantum ``particle in a box'' result.
\textbf{Mass is determined by the Bohr--Sommerfeld quantum number
$N$}, not the harmonic number $r$.

\medskip
\textbf{Important clarification:} In the classical model, $r$
determined mass ($m \propto 1/r^2$). In the quantum model, $r$
and $N$ are linked: for each $r$, only specific $N$ values are
stable (from the Mathieu equation's stability zones). Therefore the
\emph{allowed} $(r, N)$ pairs produce a \textbf{discrete mass
spectrum}.

\subsection{Relation Between Mass and Frequency}

From the Compton relation every particle has a fundamental
frequency:
\begin{equation}
  \nu_C = \frac{mc^2}{h}.
  \label{eq:compton_def}
\end{equation}

In the two-plate model the ball's bounce frequency is $\nu_0/r$.
Via the Zitterbewegung interpretation, this bouncing \textbf{is}
the particle's internal oscillation:
\begin{equation}
  \frac{\nu_0}{r} = \frac{2mc^2}{h}
  = 2\nu_C.
  \label{eq:bounce_compton}
\end{equation}
(The factor of 2 comes from the Zitterbewegung
$\omega_Z = 2mc^2/\hbar$: in one cycle the ball bounces
\emph{twice}---up and down.)

Therefore:
\begin{equation}
  \nu_0 = 2r\,\frac{mc^2}{h}
  = 2r\,\nu_C
  \label{eq:nu0_from_mass}
\end{equation}

This general relation links $\nu_0$ to any particle's
mass $m$ and its classical harmonic number $r$.

\subsection{Historical Planck-Scale Estimate}

The following calculation records the older classical estimate.  In
that classical picture, the fundamental mode ($r = 1$) corresponds to the
\textbf{maximum mass}. Physically, this is the mass above which
an oscillon collapses into a black hole---the
\textbf{Planck mass}:
\begin{equation}
  m_{r=1} = m_P = \sqrt{\frac{\hbar c}{G}}
  = 2.176 \times 10^{-8}\;\mathrm{kg}
  = 1.221 \times 10^{19}\;\mathrm{GeV}/c^2.
  \label{eq:planck_mass}
\end{equation}

From \eqref{eq:nu0_from_mass} with $r = 1$:
\begin{equation}
  \nu_0^{(\mathrm{cl})} = \frac{2m_P c^2}{h}
  = \frac{\nu_P}{\pi}
  \label{eq:nu0_value}
\end{equation}
where $\nu_P = 1/t_P = 1.855 \times 10^{43}$~Hz
is the Planck frequency.

Numerical value:
\begin{equation}
  \nu_0^{(\mathrm{cl})} = 5.904 \times 10^{42}\;\mathrm{Hz}.
  \label{eq:nu0_numerical}
\end{equation}

\textit{Historical status.}---This classical estimate
$\nu_0^{(\mathrm{cl})}$ is based on the linear
dispersion $m \propto 1/r$.  The quantum Mathieu model
(\S\ref{sec:mathieu_numerical}) replaces this with
$m \propto N^2$, yielding a corrected value
$\nu_0 = 1.91 \times 10^{30}$~Hz---thirteen orders of
magnitude lower (Eq.~\ref{eq:nu0_quantum_value}).
\textbf{Only the quantum value $\nu_0$ is used in all
subsequent calculations}; $\nu_0^{(\mathrm{cl})}$ serves
solely as a historical/pedagogical stepping stone.

Corresponding period of the classical estimate:
\begin{equation}
  T_0^{(\mathrm{cl})}
  = 1/\nu_0^{(\mathrm{cl})} = \pi\,t_P
  = 1.694 \times 10^{-43}\;\mathrm{s}.
\end{equation}
Corresponding wavelength:
\begin{equation}
  \lambda_0^{(\mathrm{cl})}
  = c/\nu_0^{(\mathrm{cl})} = \pi\,\ell_P
  = 5.08 \times 10^{-35}\;\mathrm{m}.
\end{equation}

\subsection{Classical Particle Harmonic Map}

In the historical classical map, $r = 1$ fundamental mode
$=$ Planck mass.
Each particle's harmonic number:
\begin{equation}
  r = \frac{m_P}{m}
  \label{eq:r_from_mass}
\end{equation}
(follows from \eqref{eq:nu0_from_mass}
with $\nu_0^{(\mathrm{cl})} = 2m_P c^2/h$.)

\begin{center}
\renewcommand{\arraystretch}{1.3}
\begin{tabular}{@{}lrlrl@{}}
\textbf{Particle}
  & \textbf{Mass (MeV/$c^2$)}
  & \textbf{$\nu_C$ (Hz)}
  & \textbf{$r = m_P/m$}
  & \textbf{$\nu_0^{(\mathrm{cl})}/r$ (Hz)} \\[2pt]
\hline\\[-6pt]
Top quark
  & $172\,760$
  & $4.18 \times 10^{25}$
  & $7.07 \times 10^{16}$
  & $8.35 \times 10^{25}$ \\
Higgs boson
  & $125\,100$
  & $3.02 \times 10^{25}$
  & $9.76 \times 10^{16}$
  & $6.05 \times 10^{25}$ \\
$Z$ boson
  & $91\,188$
  & $2.20 \times 10^{25}$
  & $1.34 \times 10^{17}$
  & $4.41 \times 10^{25}$ \\
$W$ boson
  & $80\,379$
  & $1.94 \times 10^{25}$
  & $1.52 \times 10^{17}$
  & $3.89 \times 10^{25}$ \\
$b$ quark
  & $4\,180$
  & $1.01 \times 10^{24}$
  & $2.92 \times 10^{18}$
  & $2.02 \times 10^{24}$ \\
$\tau$ lepton
  & $1\,776.9$
  & $4.30 \times 10^{23}$
  & $6.87 \times 10^{18}$
  & $8.59 \times 10^{23}$ \\
$c$ quark
  & $1\,270$
  & $3.07 \times 10^{23}$
  & $9.61 \times 10^{18}$
  & $6.14 \times 10^{23}$ \\
$\mu$ lepton
  & $105.66$
  & $2.55 \times 10^{22}$
  & $1.16 \times 10^{20}$
  & $5.11 \times 10^{22}$ \\
$s$ quark
  & $93.4$
  & $2.26 \times 10^{22}$
  & $1.31 \times 10^{20}$
  & $4.51 \times 10^{22}$ \\
$d$ quark
  & $4.67$
  & $1.13 \times 10^{21}$
  & $2.62 \times 10^{21}$
  & $2.26 \times 10^{21}$ \\
$u$ quark
  & $2.16$
  & $5.23 \times 10^{20}$
  & $5.65 \times 10^{21}$
  & $1.05 \times 10^{21}$ \\
Electron
  & $0.511$
  & $1.24 \times 10^{20}$
  & $\mathbf{2.39 \times 10^{22}}$
  & $2.47 \times 10^{20}$ \\[4pt]
\hline\\[-6pt]
Planck mass
  & $1.221 \times 10^{19}$
  & $2.95 \times 10^{42}$
  & $\mathbf{1}$
  & $5.90 \times 10^{42}$ \\
\end{tabular}
\end{center}

\medskip
\textbf{The hierarchy problem in ISPG language.}
The table shows that all known particles have very large $r$
($10^{16}$--$10^{22}$). The fundamental mode ($r = 1$, Planck
mass) versus the lightest known particle (electron)---$r \sim
10^{22}$.

This is particle physics' well-known \textbf{hierarchy problem}:
why are elementary particles $10^{17}$--$10^{22}$ times lighter
than the fundamental scale?

In the historical ISPG map this means: particles are
\textbf{extremely high classical harmonics} of the medium---very
slow bouncing between very rapidly vibrating plates. Modes close to
the fundamental ($r \sim 1$--$10^{10}$) are unstable---they
cannot maintain resonance outside the Mathieu equation's stability
zones.

\textit{Note:} the harmonic numbers $r$ in this table are
computed from the classical estimate
$\nu_0^{(\mathrm{cl})}$.  In the quantum Mathieu model
(\S\ref{sec:mathieu_numerical}), the relevant quantum
numbers are $N = 1, 2, 3, \ldots$ with a corrected
$\nu_0 = 1.91 \times 10^{30}$~Hz.  The physical content
(discrete spectrum, hierarchy) is retained as a qualitative
motivation, but the final labeling and numerical carrier frequency
come from the calibrated MASS/Mathieu model.

\subsection{Harmonic Structure of Lepton Masses}

Lepton mass ratios:
\begin{alignat}{2}
  \frac{m_\mu}{m_e} &= 206.768
  &\qquad&
  \sqrt{\frac{m_\mu}{m_e}} = 14.38,
  \label{eq:ratio_mu_e}\\[4pt]
  \frac{m_\tau}{m_e} &= 3\,477.5
  &&
  \sqrt{\frac{m_\tau}{m_e}} = 58.97,
  \label{eq:ratio_tau_e}\\[4pt]
  \frac{m_\tau}{m_\mu} &= 16.82
  &&
  \sqrt{\frac{m_\tau}{m_\mu}} = 4.101.
  \label{eq:ratio_tau_mu}
\end{alignat}

These numbers are close to integers but \textbf{not exactly
integers}. We seek the best approximations:

\medskip
\textbf{Approximation 1:}
$\sqrt{m_\tau/m_e} \approx 59$
\quad$\Rightarrow$\quad
$m_\tau/m_e \approx 59^2 = 3481$
\quad vs.\ actual $3477.5$
\quad($0.1\%$ error).

\textbf{Approximation 2:}
$m_\mu/m_e \approx (72/5)^2 = 207.36$
\quad vs.\ actual $206.77$
\quad($0.3\%$ error).

\textbf{Approximation 3:}
$m_\tau/m_\mu \approx (41/10)^2 = 16.81$
\quad vs.\ actual $16.82$
\quad($0.06\%$ error).

\medskip
These approximations show that $\sqrt{m_1/m_2}$ tends toward
\emph{rational} numbers, consistent with the harmonic-rational
resonance model. Exact values require solving the Mathieu equation.

\subsection{Mathieu Equation: Stability Zones and Mass Selection}

The two-plate model is analogous to the Fermi--Ulam problem
(a ball between oscillating walls). In the linear approximation,
the dynamics is described by the Mathieu equation:
\begin{equation}
  \frac{d^2\psi}{d\tau^2}
  + \bigl[a - 2q\cos(2\tau)\bigr]\,\psi = 0,
  \label{eq:mathieu_general}
\end{equation}
where
\begin{equation}
  a = \left(\frac{E}{h\nu_0}\right)^2
  = \left(\frac{E}{\hbar\omega_0}\right)^2
  = \left(\frac{mc^2}{h\nu_0}\right)^2,
  \qquad
  \omega_0 \equiv 2\pi\nu_0,
  \qquad
  q \approx 1 + \delta q,
  \label{eq:mathieu_params}
\end{equation}
where $q$ arises from the nonlinear kinetic self-interaction
$G_2(X) \supset \alpha X^2/M^2$ of the ISPG scalar field.
\textbf{$q$ does not depend on $N$}---this is a consequence of
ISPG's conformal coupling
(detailed derivation: \S\ref{sec:q_derivation}).

The Mathieu equation's \textbf{stability map} in the $(a, q)$
plane creates zones where solutions are bounded (stable oscillon)
and zones where solutions grow (unstable---particle decays).

\textbf{Allowed particle masses $=$ stability zone $a$-values:}
\begin{equation}
  m_n = \frac{h\nu_0}{c^2}\,\sqrt{a_n(q)}.
  \label{eq:mass_from_mathieu}
\end{equation}

The Planck-mass form belongs only to the historical
Planck-anchored estimate of $\nu_0^{(\mathrm{cl})}$ below.
The calibrated MASS/Mathieu construction uses the cyclic
frequency $\nu_0$ with $h\nu_0=\hbar\omega_0$ and does not
assume the historical $m_P/\pi$ normalization.

After fixing $q$ by calibration, or eventually deriving it from
the full oscillon dynamics, the Mathieu characteristic values
$a_n(q)$ yield the \textbf{particle mass spectrum} in this
one-dimensional EFT limit,
analogous to Bohr's formula $E_n = E_0/n^2$ for hydrogen.

\textbf{Estimate of $\varepsilon$.}
If $\varepsilon \ll 1$ (weak fluctuation), the principal stability
zones are centered at $a = n^2$ ($n = 0, 1, 2, \ldots$) and mass
ratios are approximately $n_2/n_1$.
If $\varepsilon \sim 1$ (strong fluctuation), stability zones shift
dramatically and produce a \emph{complex} mass spectrum---possibly
explaining why elementary particle masses appear in ``unexpected''
ratios.

\subsection{Historical Interpretation of $\nu_0^{(\mathrm{cl})}$}

\begin{center}
\fbox{\parbox{0.88\textwidth}{\centering\large
\textbf{Historical classical estimate}\\[3pt]
$\displaystyle\nu_0^{(\mathrm{cl})} = \frac{\nu_P}{\pi}
\approx 5.9 \times 10^{42}\;\text{Hz}$

\medskip\normalsize
This is not the current calibrated carrier frequency.  It is the
older Planck-anchored estimate obtained from the classical
$m \propto 1/r$ map.  The present paper uses the calibrated
MASS/Mathieu value $\nu_0 = 1.91\times 10^{30}$~Hz for
oscillonic physics.
}}
\end{center}

\begin{itemize}
  \item \textbf{Historical particle $=$ sub-harmonic picture.}
    The electron ``bounces'' every
    $r_e \approx 2.4 \times 10^{22}$ medium vibrations. The top
    quark---every $r_t \approx 7 \times 10^{16}$ vibrations.
    This is the classical intuition, not the final Mathieu
    calibration.
  \item \textbf{Classical Planck limit.} $r = 1$---bouncing on every
    vibration---yields
    $m_P = 1.22 \times 10^{19}$~GeV.
    A heavier classical oscillon would collapse in this picture.
  \item \textbf{Classical wavelength $=$
    $\pi\,\ell_P \approx 5 \times 10^{-35}$~m.}
    This belongs to the historical estimate
    $\lambda_0^{(\mathrm{cl})}=c/\nu_0^{(\mathrm{cl})}$.
    It should not be read as the wavelength of the calibrated
    $\nu_0$.
  \item \textbf{Historical mass-frequency map.}
    In the classical picture, $m=m_P/r$.  The quantitative
    particle spectrum used in this paper is instead the calibrated
    MASS/Mathieu spectrum.
\end{itemize}

\subsection{Status of the Program}

The historical value $\nu_0^{(\mathrm{cl})}$
\eqref{eq:nu0_numerical} relies on the assumption that
$r = 1$ corresponds to the Planck mass. The quantum
MASS/Mathieu model performs the recalibration used below.
The following status summary separates completed results from
open closure tasks:
\begin{enumerate}
  \item[\checkmark] \textbf{Derivation of $q$ from ISPG
    $G_2(X)$ extension at leading order}---
    $q_0 = 1$ follows from flat-space self-consistency.
  \item[$\circ$] \textbf{Nonlinear correction to the fitted
    Mathieu parameter}---
    the fitted reference value is
    $q_{\mathrm{fit}} = 1.853$; the present semi-analytic
    correction bundle gives $q_{\mathrm{semi}}\simeq 1.907$,
    within $2.9\%$ of that value. Exact closure of
    $\delta q$ requires the full lattice/PDE oscillon
    calculation (\S\ref{sec:q_derivation}).
  \item[\checkmark] \textbf{Solution of Mathieu stability
    zones}---
    $b_N(q)$ for all particles
    (\S\ref{sec:mathieu_numerical}).
  \item[\checkmark] \textbf{Mass ratio comparison}---
    muon: $< 0.001\%$; leptons: $< 0.2\%$
    (\S\ref{sec:mathieu_numerical}).
  \item[\checkmark] \textbf{Reproduction of the Koide
    formula}---
    $0.028\%$ accuracy
    (\S\ref{sec:mathieu_numerical}).
\end{enumerate}

%% ==========================================================
\subsection{Numerical Results of the Quantum Model}
\label{sec:numerical_results}
%% ==========================================================

By Bohr--Sommerfeld quantization
(Eq.~\ref{eq:mass_quantum_box}), mass scales as $N^2$:
\begin{equation}
  m = m_1 \times N^2,
  \qquad m_1 \equiv \frac{h^2}{2\mu\,\ell_0^2\,c^2},
  \qquad N \in \mathbb{Z}^{+}.
  \label{eq:mass_N_squared}
\end{equation}

The electron is one of the harmonics: $m_e = m_1 N_e^2$.
\textbf{Goal:} find $N_e$ and then all other particles' $N$.

\medskip
\textbf{Determining $N_e$ from lepton masses.}

Mass ratio: $m_\mu/m_e = N_\mu^2/N_e^2$.
Hence $N_\mu/N_e = \sqrt{m_\mu/m_e}
= \sqrt{206.768} = 14.380$.

We seek a rational approximation:
\begin{equation}
  14.380 \approx \frac{72}{5} = 14.400
  \qquad(\text{error}\;0.14\%).
  \label{eq:mu_ratio_approx}
\end{equation}

This gives: $N_e = 5$, $N_\mu = 72$.

For the tau lepton:
$N_\tau/N_e = \sqrt{m_\tau/m_e} = \sqrt{3477.5}
= 58.97 \approx 59$.
Hence $N_\tau = 5 \times 59 = 295$.

\begin{equation}
  \boxed{N_e = 5, \qquad N_\mu = 72, \qquad N_\tau = 295}
  \label{eq:lepton_N}
\end{equation}

\medskip
\textbf{Prediction vs.\ experiment (leptons):}

\begin{center}
\renewcommand{\arraystretch}{1.3}
\begin{tabular}{@{}lcccr@{}}
  & $N$ & $m_{\mathrm{predicted}}$ (MeV)
  & $m_{\mathrm{actual}}$ (MeV) & Error \\[2pt]
\hline\\[-6pt]
$e$ & 5 & $0.511$ (reference) & $0.5110$ & --- \\
$\mu$ & 72 & $105.96$ & $105.658$ & $0.29\%$ \\
$\tau$ & 295 & $1778.8$ & $1776.86$ & $0.11\%$ \\
\end{tabular}
\end{center}

\medskip
\textbf{Koide consistency on the selected lepton ladder.}

The Koide formula \eqref{eq:koide} with $m \propto N^2$ becomes:
\begin{equation}
  \frac{N_e^2 + N_\mu^2 + N_\tau^2}
       {(N_e + N_\mu + N_\tau)^2}
  = \frac{25 + 5184 + 87025}{(5 + 72 + 295)^2}
  = \frac{92\,234}{138\,384}
  = 0.66654.
  \label{eq:koide_from_N}
\end{equation}
Exact value: $2/3 = 0.66667$.
\textbf{Error: $0.02\%$!}

This means: once the lepton integers $(5, 72, 295)$ are
selected from the mass pattern, the Koide formula---an
unexplained empirical relation since 1981---is reproduced by
the $N^2$ ladder to $0.024\%$.  A rank-1/E1D consistency
argument supporting this assignment is given in
ISPG\_Quantum, Sec.~7.2; the derivation of the lepton
integers $(5, 72, 295)$ from oscillon dynamics, and the
formal 3D-to-1D bridge to first-principles charged-lepton
masses, remain open.

\medskip
\textbf{Harmonic map of all elementary particles.}

Fundamental mass: $m_1 = m_e/N_e^2
= m_e/25 = 20.44$~keV.
The harmonic quantum number of each particle is
$N = \mathrm{round}(5\sqrt{m/m_e})$.

\textit{Methodological note.}---The integer $N$
is extracted from measured masses; the table
below is therefore a \emph{one-parameter
compression} of the mass spectrum, not yet a
blind prediction.  The non-trivial content is
threefold: (i)~twelve quark, lepton and boson
masses, which in the Standard Model require
twelve independent Yukawa couplings, here
collapse onto a single fundamental scale $m_1$
and a single nonlinearity parameter $q$;
(ii)~the integers $N$ are not arbitrary---they
coincide with the Mathieu stability chart
(unstable/wide-gap for $N = 1$--$3$, barely stable for
$N = 4$, fully stable for $N \ge 5$ under the adopted
threshold convention), which explains why the electron is
the first fully stable occupied charged Mathieu mode in this
scheme;
(iii)~the residual errors are case-by-case after the
nearest-$N$ selection, with the current PDG-style scan giving
$10/12$ entries below $1\%$ and $11/12$ below $5\%$.
Light-quark masses depend on the renormalization scheme
(PDG running $\overline{\mathrm{MS}}$ masses at $2$~GeV
versus the paper-target values), while the high-$N$ heavy
entries are empirical placements until a physical
$N$-selection principle is derived.
The programme becomes fully predictive once the fitted
reference value of $q$ is replaced by an independently
derived value from the lattice simulation of $G_2(X)$
(future work, \S\ref{sec:future}).

\begin{center}
\renewcommand{\arraystretch}{1.3}
\begin{tabular}{@{}lrrrr@{}}
\textbf{Particle}
  & \textbf{$m$ (MeV)}
  & \textbf{$N$}
  & \textbf{$m_{\mathrm{pred}}$ (MeV)}
  & \textbf{Error} \\[2pt]
\hline\\[-6pt]
$e$ & $0.511$ & $5$ & $0.511$ & --- \\
$u$ quark & $2.16$ & $10$ & $2.04$ & $5.4\%$\rlap{$^*$} \\
$d$ quark & $4.67$ & $15$ & $4.60$ & $1.5\%$\rlap{$^*$} \\
$s$ quark & $93.4$ & $68$ & $94.5$ & $1.2\%$\rlap{$^*$} \\
$\mu$ & $105.66$ & $72$ & $105.96$ & $0.29\%$ \\
$c$ quark & $1\,270$ & $249$ & $1\,267$ & $0.24\%$ \\
$\tau$ & $1\,776.9$ & $295$ & $1\,778.8$ & $0.11\%$ \\
$b$ quark & $4\,180$ & $452$ & $4\,176$ & $0.10\%$ \\
$W$ & $80\,379$ & $1\,983$ & $80\,363$ & $0.02\%$ \\
$Z$ & $91\,188$ & $2\,112$ & $91\,215$ & $0.03\%$ \\
$H$ & $125\,100$ & $2\,474$ & $125\,132$ & $0.03\%$ \\
$t$ quark & $172\,760$ & $2\,907$ & $172\,686$ & $0.04\%$ \\
\end{tabular}
\end{center}
\rlap{$^*$}\;Light quark masses are poorly known
experimentally ($\sim 10$--$30\%$ uncertainty).

\medskip
\textbf{Assessment:} The table is a compact empirical
ladder placement: one scale ($m_1$), one fitted
nonlinearity ($q$), and one integer per particle compress
the mass data, but the integer choice is not yet predicted.
The light-quark entries require an explicit mass convention,
and the heavy high-$N$ entries are close nearest-$N$ scan
hits rather than first-principles electroweak/Higgs mass
derivations.

\medskip
\textbf{Physical interpretation of $m_1$.}
The fundamental mass $m_1 = 20.44$~keV corresponds to the
``slowest bouncing'' ($N = 1$) energy in the box.
This is $\sim 40$ times lighter than the electron.
In the adopted gap-threshold scheme, $N=1$ and $N=2$ are
unstable (gaps $\sim3.5$ and $\sim1.3$), $N=3$ is
marginal/short-lived (gap $\sim0.18$), $N=4$ is the first
barely-stable candidate (gap $\sim10^{-2}$), and $N=5$ is
the first fully stable mode (gap $<10^{-3}$):
\begin{alignat}{2}
  m_{N=2} &= 4\,m_1 = 81.8\;\text{keV}
  &\qquad& (\text{unstable Mathieu mode}),\\
  m_{N=3} &= 9\,m_1 = 184\;\text{keV}
  && (\text{marginal/short-lived resonance}).
\end{alignat}

\subsection{Recalibration of $\nu_0$ via the Quantum Model}

From the Bohr--Sommerfeld condition
\eqref{eq:nu0_from_BS}: $\nu_0 = Nrh/(4\mu\ell_0^2)$,
where $Nr = K$ (constant for all particles).

The Planck limit ($m_P$ = maximum mass) now corresponds to
$N = N_{\max}$:
\begin{equation}
  N_{\max} = \frac{N_e\,m_P^{1/2}}{m_e^{1/2}}
  = 5\,\sqrt{\frac{m_P}{m_e}}
  = 5 \times 1.546 \times 10^{11}
  = 7.73 \times 10^{11}.
\end{equation}

$\nu_0$ in the quantum model:
\begin{equation}
  \nu_0
  = \frac{\sqrt{m_P\,m_e}\;c^2}{10\,h}
  = \frac{\sqrt{m_P\,m_e}}{10}\,\nu_C\!\left(
  \frac{1}{\sqrt{m_P m_e}}\right).
  \label{eq:nu0_quantum}
\end{equation}

Numerically:
\begin{align}
  \sqrt{m_P\,m_e}
  &= \sqrt{2.176 \times 10^{-8}
    \times 9.109 \times 10^{-31}}
  = 1.408 \times 10^{-19}\;\text{kg}
  \notag\\
  &= 7.90 \times 10^{4}\;\text{MeV}/c^2
  \approx 79\;\text{GeV}/c^2.
  \label{eq:geometric_mean_mass}
\end{align}

\begin{equation}
  \boxed{\nu_0
  \equiv \nu_0^{(\mathrm{quantum})}
  = 1.91 \times 10^{30}\;\mathrm{Hz}}
  \label{eq:nu0_quantum_value}
\end{equation}

\textbf{This is the calibrated MASS/Mathieu value of $\nu_0$ used
throughout this paper.}  It supersedes the historical classical
estimate $\nu_0^{(\mathrm{cl})} = 5.9 \times 10^{42}$~Hz
(Eq.~\ref{eq:nu0_numerical}), which is $10^{13}$ times
larger.  The difference arises because
the classical model has $m \propto 1/r$, while the quantum
model has $m \propto N^2 \propto 1/r^2$. The quantum model
yields \textbf{significantly better} mass ratios.  The
frequency-invariance theorem does not compute this number; it
establishes the background independence of the proper-time carrier
once the carrier is given.  A first-principles derivation of the
calibrated value remains open.

\medskip
\textbf{Interesting coincidence:}
The geometric mean $\sqrt{m_P\,m_e}
\approx 79$~GeV is very close to the electroweak scale
($v_{\mathrm{EW}} = 246$~GeV, $m_W = 80$~GeV).
This is not accidental: the electroweak scale is
expected, in the ISPG completion, to be fixed by the
same oscillon-amplitude normalization that determines
$\Phi_0$; the explicit coarse-graining map from the
dimensionless amplitude $\Phi_0$ to the imported energy
scale $v_{\mathrm{EW}}=246$~GeV remains open
(\S\ref{sec:higgs}).

\subsection{Numerical Solution of the Mathieu Equation}
\label{sec:mathieu_numerical}

The Mathieu equation's \eqref{eq:mathieu_general}
quantum characteristic values $b_N(q)$ are computed using the
\texttt{scipy} library.

\medskip
\textbf{Fitting the reference value of $q$ from the muon mass.}

We define the fitted reference value $q_{\mathrm{fit}}$ by requiring
$b_{72}(q)/b_5(q) = m_\mu/m_e = 206.768$:
\begin{equation}
  \boxed{q_{\mathrm{fit}} = 1.853}
  \label{eq:q_value}
\end{equation}

At this fitted value the muon/electron ratio is reproduced
by construction ($< 0.001\%$ numerical error). The independent
theory task is then to derive a value of $q$ from the oscillon
dynamics that closes on this reference value.

\medskip
\textbf{Stability analysis: why $N_e = 5$?}

Instability zone widths at $q_{\mathrm{fit}} = 1.853$:

\begin{center}
\renewcommand{\arraystretch}{1.3}
\begin{tabular}{@{}crrrll@{}}
$N$ & $b_N(q)$ & $a_N(q)$ & Gap width & Status
  & \\[2pt]\hline\\[-6pt]
$1$ & $-1.19$ & $2.33$ & $3.52$
  & \textbf{Unstable}
  & $b_1 < 0$: does not exist\\
$2$ & $3.72$ & $5.04$ & $1.33$
  & \textbf{Unstable} & \\
$3$ & $9.13$ & $9.31$ & $0.185$
  & Marginal & \\
$4$ & $16.11$ & $16.12$ & $0.010$
  & Stable & \\
$\mathbf{5}$ & $\mathbf{25.07}$ & $\mathbf{25.07}$
  & $\mathbf{0.0003}$
  & \textbf{Stable}
  & $\leftarrow$ \textbf{Electron}\\
$6$ & $36.05$ & $36.05$ & $5 \times 10^{-6}$
  & Stable & \\
$7$ & $49.04$ & $49.04$ & $< 10^{-7}$
  & Stable & \\
\end{tabular}
\end{center}

\textbf{Result:}
\begin{itemize}
  \item $N = 1$: \textbf{does not exist}---
    $b_1(q) < 0$; no bound state corresponds to this mode.
  \item $N = 2$: \textbf{unstable}---
    instability zone is wide (1.33);
    the particle decays instantly.
  \item $N = 3$: \textbf{marginal}---
    zone width $0.185$; very short-lived resonance.
  \item $N = 4$: stable, but
    zone width $0.01$---possibly a
    ``barely stable'' particle.
  \item $N \geq 5$: \textbf{fully stable}.
    Zone width $< 10^{-3}$.
\end{itemize}

\textit{Status of the gap-width criterion.}---The
designations ``marginal'', ``barely stable'', and ``fully
stable'' above are based on the Mathieu gap width (a
measure of instability-zone size in $(a,q)$ space).  The
thresholds separating ``marginal'' ($\sim\!10^{-1}$),
``barely stable'' ($\sim\!10^{-2}$), and ``fully stable''
($<\!10^{-3}$) are adopted here for concreteness and are
not derived from first principles.  A rigorous mapping
from the gap width to a physical decay rate requires a
full Floquet analysis of the driven oscillon including
its dissipation channels; that mapping is not derived in
the present section, so the lifetime estimate $\tau
\lesssim 10^{-20}$~s for the $N=4$ resonance (below) should
be read as an order-of-magnitude guideline rather than a
derived prediction.

\begin{center}
\fbox{\parbox{0.85\textwidth}{\centering
\textbf{The electron ($N = 5$) is the first fully stable
charged Mathieu mode under the adopted gap-threshold scheme;
$N=1$--$3$ are unstable/marginal, and $N=4$ is a barely-stable
resonance candidate predicted at $\sim 329$~keV.}
}}
\end{center}

This explains one of nature's fundamental questions:
\textbf{why does a lightest charged particle exist, and why
does it have this particular mass?}

\medskip
\textbf{The neutrino question.}
For neutrino masses ($m_\nu \sim 0.05$~eV):
$N_\nu = 5\sqrt{m_\nu/m_e} \approx 0.002 \ll 1$.
This means neutrinos \textbf{are not} ordinary harmonics
of the two-plate model.  In the ISPG framework, neutrinos
are identified as \textbf{polarisation-rotation waves}---%
helical excitations whose oscillation plane progressively
rotates---a fundamentally different vibration mode whose
suppressed torsional rigidity naturally yields sub-eV masses
(\S\ref{sec:neutrinos}).

%% ==========================================================
\subsection{Derivation of $q$ from the ISPG Action:
  Shift-Symmetric $G_2(X)$ Extension}
\label{sec:q_derivation}
%% ==========================================================

In the previous section $q_{\mathrm{fit}} = 1.853$ was fixed
from the muon/electron ratio as a reference calibration. We now
show that the \textbf{leading-order value $q_0 = 1$ follows from
first principles}, and identify three calculable correction
channels (metric backreaction, spatial gradients, higher
harmonics) that drive the theory toward the fitted remainder
$\delta q_{\mathrm{fit}} = 0.853$---all without introducing new
fields or breaking the shift symmetry
$\phi \to \phi + \mathrm{const}$ that protects gravity's
long-range nature. The present semi-analytic evaluation gives
$q_{\mathrm{semi}}\simeq 1.907$, within $2.9\%$ of
$q_{\mathrm{fit}}$; exact closure of $\delta q$ awaits lattice
computation of the oscillon profile (\S\ref{sec:q_derivation},
Step~6).

\textit{Relation to the Quantum paper's mass program.}---The
Quantum paper (Sec.~7) derives the mass hierarchy from the
eigenvalue spectrum of three-dimensional resonant cavities
embedded in the scalar-metric medium.  We regard the
one-dimensional Mathieu reduction used here as a candidate
EFT limit of that 3D problem, motivated by the following
heuristic step: when the oscillon core is approximated as
spatially homogeneous (Step~3 below), the radial eigenvalue
equation has the same structural form as a one-dimensional
parametric oscillator, i.e.\ a Mathieu equation.  We note,
however, that a \emph{formal} reduction from the 3D cavity
problem to the one-dimensional Mathieu equation is not
established here: the QUANTUM $(n,\ell)$ lepton assignments
$\tau=(0,0),\ \mu=(1,0),\ e=(0,2)$ are not in one-to-one
correspondence with the $N$ assignments used below, and the
quark/gauge-boson sectors addressed by the present
one-dimensional model do not appear as cavity modes in
QUANTUM Sec.~7.  The two routes are therefore better
described as \emph{parallel frameworks} with overlapping
but non-identical scope, whose eventual unification via a
formal 3D-to-1D reduction (or via lattice simulation of the
full 3D system) is left as an open calculation.

\medskip
\textbf{Step 1: The ISPG action and its natural extension.}

The baseline ISPG action (Appendix~1 of the main theory) is
\begin{equation}
  S = \frac{1}{16\pi G}\int d^4x\sqrt{-g}\left[
    R - \frac{M_{\mathrm{Pl}}^2}{2}\,X
  \right] + S_m,
  \qquad
  X \equiv -\tfrac{1}{2}\,g^{\mu\nu}
  \partial_\mu\phi\,\partial_\nu\phi,
  \label{eq:action_baseline}
\end{equation}
where $\phi$ is the \emph{single} scalar field whose
bi-conformal coupling
$ds^2 = -e^{\phi}c^2 dt^2 + e^{-\phi}d\sigma^2$
determines all of gravity. The vacuum field equation is
$\Box\phi = 0$.

The action \eqref{eq:action_baseline} is the leading term of a
derivative expansion. The next shift-symmetric operator compatible
with the Horndeski structure (Appendix~11, Stability) is $X^2$.
This is the same $G_2(X)$ extension that was independently
established for ghost-free UV completion of the phantom kinetic
sector; here it acquires a second physical role---generating
the oscillon nonlinearity that produces particle masses:
\begin{equation}
  \boxed{
  G_2(X) = -\frac{M_{\mathrm{Pl}}^2}{2}\,X
  + \frac{\alpha}{M^2}\,X^2,
  \qquad \alpha > 0,
  }
  \label{eq:G2X}
\end{equation}
where $M$ is the UV scale at which the nonlinearity becomes
important, and $\alpha$ is a dimensionless coupling constant.

Three properties are essential:
\begin{enumerate}
  \item \textbf{Same field.} No new field is introduced.
    The action contains only $\phi$, the same field
    that generates the bi-conformal metric, mediates
    gravity, and forms oscillons.
  \item \textbf{Shift symmetry preserved.}
    $X = -\frac{1}{2}(\partial\phi)^2$ depends only on
    \emph{derivatives} of $\phi$, so $\phi \to \phi + c$
    leaves the action invariant. No mass term
    $\mu^2\phi^2$ is generated, and gravity remains
    long-range.
  \item \textbf{Weak-field transparency.}
    For $X \ll M^2$
    (weak gravity, PPN regime), $X^2/M^2 \ll X$ and
    the action reduces to the baseline ISPG.
    All existing observational tests are unaffected.
\end{enumerate}

\medskip
\textbf{Step 2: Field equation from $G_2(X)$.}

The Euler--Lagrange equation for the scalar sector with
$\mathcal{L}_\phi = G_2(X)$ is
\begin{equation}
  \nabla_\mu\!\left(
    \frac{\partial G_2}{\partial X}\,
    \nabla^\mu\phi
  \right) = 0.
  \label{eq:G2_eom_general}
\end{equation}
Computing the derivative:
\begin{equation}
  \frac{\partial G_2}{\partial X}
  = -\frac{M_{\mathrm{Pl}}^2}{2}
  + \frac{2\alpha}{M^2}\,X.
  \label{eq:dG2dX}
\end{equation}
Define the \emph{effective sound-speed function}
\begin{equation}
  \mathcal{K}(\phi)
  \equiv \frac{\partial G_2}{\partial X}
  = -\frac{M_{\mathrm{Pl}}^2}{2}
  + \frac{2\alpha}{M^2}\,X.
  \label{eq:K_def}
\end{equation}
The field equation becomes
\begin{equation}
  \mathcal{K}\,\Box\phi
  + \frac{2\alpha}{M^2}\,
  (\nabla^\mu\phi)\,(\nabla^\nu\phi)\,
  \nabla_\mu\nabla_\nu\phi = 0.
  \label{eq:full_eom}
\end{equation}
In the weak-field limit ($X \to 0$):
$\mathcal{K} \to -M_{\mathrm{Pl}}^2/2$ and the second
term vanishes, recovering $\Box\phi = 0$.

\medskip
\textbf{Step 3: Oscillon ansatz.}

We seek a localized oscillating solution of the same form
used in the ISPG quantum paper:
\begin{equation}
  \phi(t,\mathbf{x}) = \Phi_0(r)\cos(\omega_0 t),
  \label{eq:oscillon_ansatz}
\end{equation}
where $\Phi_0(r)$ is a localized profile and
$\omega_0 = 2\pi\nu_0$.

For this ansatz:
\begin{equation}
  X = -\tfrac{1}{2}(\partial\phi)^2
  = \tfrac{1}{2}\bigl[
    \dot\phi^2 - (\nabla\phi)^2
  \bigr]
  = \tfrac{1}{2}\left[
    \Phi_0^2\omega_0^2\sin^2(\omega_0 t)
    - (\nabla\Phi_0)^2\cos^2(\omega_0 t)
  \right].
  \label{eq:X_oscillon}
\end{equation}
At the oscillon core ($r \approx 0$) where
$|\nabla\Phi_0| \ll \Phi_0\omega_0$, the kinetic term
dominates:
\begin{equation}
  X \approx \tfrac{1}{2}\Phi_0^2\omega_0^2
  \sin^2(\omega_0 t)
  \equiv X_0\sin^2(\omega_0 t),
  \qquad
  X_0 \equiv \tfrac{1}{2}\Phi_0^2\omega_0^2.
  \label{eq:X0_def}
\end{equation}

\medskip
\textbf{Step 4: Self-consistency condition.}

Substituting into \eqref{eq:full_eom} and keeping the leading
terms at the oscillon core, the self-consistency condition for
the fundamental harmonic $\cos(\omega_0 t)$ requires balancing
the linear and cubic terms. The $X^2$ contribution generates
an effective cubic nonlinearity:
\begin{equation}
  \dot\phi^2 \cdot \Box\phi
  \supset \Phi_0^2\omega_0^2\sin^2(\omega_0 t)
  \cdot\bigl[-\omega_0^2\Phi_0\cos(\omega_0 t)\bigr]
  = -\Phi_0^3\omega_0^4\sin^2\!\cdot\cos,
\end{equation}
and after Fourier projection onto $\cos(\omega_0 t)$
(using $\langle\sin^2\!\cdot\cos^2\rangle
= \frac{1}{8}$ for the relevant cross-term), the
self-consistency condition becomes
\begin{equation}
  -\omega_0^2
  + \frac{3\alpha\,\Phi_0^2\omega_0^4}{4\,M^2
  \cdot M_{\mathrm{Pl}}^2/2}
  \cdot \omega_0^2 = 0.
\end{equation}
This determines the oscillon amplitude:
\begin{equation}
  \boxed{
  \Phi_0^2 = \frac{2\,M^2 M_{\mathrm{Pl}}^2}
  {3\,\alpha\,\omega_0^2}
  }
  \label{eq:amplitude_G2}
\end{equation}

\medskip
\textbf{Step 5: Perturbation equation $\to$ Mathieu.}

Consider a small perturbation on the oscillon background:
$\phi = \phi_{\mathrm{bg}} + \delta\phi$.
Linearizing \eqref{eq:full_eom}, the perturbation sees a
time-dependent effective ``mass'' from the oscillating
$\mathcal{K}$:
\begin{equation}
  \mathcal{K}(t) = -\frac{M_{\mathrm{Pl}}^2}{2}
  + \frac{2\alpha}{M^2}\cdot X_0\sin^2(\omega_0 t)
  = -\frac{M_{\mathrm{Pl}}^2}{2}
  \left[1 - \frac{2\alpha\,\Phi_0^2\omega_0^2}
  {M^2 M_{\mathrm{Pl}}^2}\,\sin^2(\omega_0 t)\right].
\end{equation}
Define the \emph{nonlinearity parameter}
\begin{equation}
  \varepsilon \equiv
  \frac{2\alpha\,\Phi_0^2\omega_0^2}
  {M^2 M_{\mathrm{Pl}}^2}
  = \frac{2\alpha\,X_0}
  {M^2 \cdot M_{\mathrm{Pl}}^2/2}
  = \frac{4}{3},
  \label{eq:epsilon_def}
\end{equation}
where the last equality uses the self-consistency
condition \eqref{eq:amplitude_G2}.

Using $\sin^2(\omega_0 t) = \frac{1}{2}
[1 - \cos(2\omega_0 t)]$, the perturbation equation
becomes, after rescaling $\tau = \omega_0 t$:
\begin{equation}
  \frac{d^2\delta\phi_k}{d\tau^2}
  + \left[a_k - 2q\cos(2\tau)\right]\delta\phi_k = 0,
  \label{eq:mathieu_from_G2}
\end{equation}
which is the \textbf{Mathieu equation}, with the parameter
\begin{equation}
  \boxed{
  q = \frac{\varepsilon}{2\,(2 - \varepsilon)}
  = \frac{2/3}{2 - 4/3}
  = \frac{2/3}{2/3} = 1
  }
  \label{eq:q_leading}
\end{equation}
at leading order.

\medskip
\textbf{Step 6: Bi-conformal corrections to $q$.}

The leading-order result $q = 1$ was derived in flat
spacetime. On the ISPG bi-conformal background,
$ds^2 = -e^{\phi}c^2 dt^2 + e^{-\phi}d\sigma^2$,
three effects shift $q$ upward. We now compute each
explicitly.

\medskip
\textit{(a) Metric backreaction.}

The d'Alembertian on the bi-conformal background
(established in the Quantum paper, Eq.~(3.14)) is
\begin{equation}
  \Box\phi = e^{\phi}\bigl[
    -e^{-2\phi}\,\ddot\phi + \nabla^2\phi
  \bigr],
  \label{eq:box_biconf}
\end{equation}
where the spatial Laplacian carries \emph{no}
$\phi$-dependent prefactor: the identity
$\sqrt{-g}\,g^{ij} = \delta^{ij}$ is an exact
property of the bi-conformal structure.

For the oscillon background
$\phi_{\mathrm{bg}} = \Phi_0\cos(\omega_0 t)$,
the temporal coefficient becomes
$e^{-2\Phi_0\cos(\omega_0 t)}$.
This factor modulates the kinetic energy of
perturbations at all harmonics of $\omega_0$.
The Jacobi--Anger expansion gives:
\begin{equation}
  e^{-2\Phi_0\cos(\omega_0 t)}
  = I_0(2\Phi_0)
  + 2\sum_{n=1}^{\infty}(-1)^n\,
  I_n(2\Phi_0)\cos(n\omega_0 t),
  \label{eq:jacobi_anger}
\end{equation}
where $I_n$ are modified Bessel functions of the
first kind. The component that directly contributes
to the Mathieu driving ($\cos 2\omega_0 t$) is
\begin{equation}
  \delta q_{\mathrm{metric}}
  \;\sim\;
  \frac{I_2(2\Phi_0)}{I_0(2\Phi_0)}.
  \label{eq:dq_metric}
\end{equation}
For $\Phi_0 \ll 1$:
$I_2(z) \approx z^2/8$,
$I_0(z) \approx 1$, so
$\delta q_{\mathrm{metric}} \approx \Phi_0^2/2$.
For $\Phi_0 \sim 1$ (the oscillon regime), the
full Bessel ratio saturates toward unity:

\smallskip
\begin{center}
\renewcommand{\arraystretch}{1.2}
\begin{tabular}{@{}cccc@{}}
$\Phi_0$ & $I_0(2\Phi_0)$ & $I_2(2\Phi_0)$
  & $I_2/I_0$ \\[2pt]\hline\\[-6pt]
$0.5$ & $1.266$ & $0.136$ & $0.107$ \\
$1.0$ & $2.280$ & $0.689$ & $0.302$ \\
$1.5$ & $4.880$ & $2.244$ & $0.460$ \\
$2.0$ & $11.30$ & $6.422$ & $0.568$ \\
\end{tabular}
\end{center}

\smallskip
The metric backreaction alone accounts for
a fraction of the total correction; the
remainder comes from gradients and higher harmonics.

\medskip
\textit{(b) Gradient corrections.}

The oscillon spatial profile established in the
Quantum paper has the form
$\Phi_0(r) = A\,\mathrm{sech}^2(r/R_c)$, where
$R_c \sim \hbar/(mc)$ is the Compton-scale core
radius and $A \sim \mathcal{O}(1)$ is the peak
amplitude. The gradient terms
$(\nabla\Phi_0)^2 \sim \Phi_0^2/R_c^2$, neglected
in the homogeneous approximation
\eqref{eq:X0_def}, modify the kinetic variable:
\begin{equation}
  X = \tfrac{1}{2}\Phi_0^2\omega_0^2\sin^2(\omega_0 t)
  - \tfrac{1}{2}(\nabla\Phi_0)^2\cos^2(\omega_0 t).
\end{equation}
The $\cos^2(\omega_0 t) = \frac{1}{2}
[1 + \cos(2\omega_0 t)]$ part of the gradient term
directly shifts the Mathieu parameter:
\begin{equation}
  \delta q_{\mathrm{grad}}
  \;\sim\;
  \frac{(\nabla\Phi_0)^2}{\Phi_0^2\omega_0^2}
  \;\sim\;
  \frac{1}{(R_c\omega_0)^2}.
  \label{eq:dq_grad}
\end{equation}
For a Compton-scale oscillon with
$R_c\omega_0 \sim c/c = \mathcal{O}(1)$, this
correction is also $\mathcal{O}(\Phi_0^2)$.

\medskip
\textit{(c) Higher harmonics.}

The cubic nonlinearity from $X^2$ generates
$\cos(3\omega_0 t)$ components via the identity
$\cos^3\theta = \frac{3}{4}\cos\theta
+ \frac{1}{4}\cos 3\theta$.
These third-harmonic terms couple back to the
fundamental through the parametric driving,
shifting the effective frequency and contributing
\begin{equation}
  \delta q_{\mathrm{harm}}
  \;\sim\;
  \mathcal{O}(\Phi_0^2).
  \label{eq:dq_harm}
\end{equation}

\medskip
\textit{Combined result.}

All three corrections are positive and of order
$\Phi_0^2 \sim \mathcal{O}(1)$ in the strongly
nonlinear oscillon regime
($\varepsilon = 4/3$). The total:
\begin{equation}
  \boxed{
  q = 1 + \delta q_{\mathrm{metric}}
  + \delta q_{\mathrm{grad}}
  + \delta q_{\mathrm{harm}}
  = 1 + \beta\,\Phi_0^2,
  }
  \label{eq:q_corrected}
\end{equation}
where $\beta > 0$ is an $\mathcal{O}(1)$ coefficient
that absorbs all three contributions. Using the
self-consistency condition \eqref{eq:amplitude_G2}:
\begin{equation}
  \Phi_0^2
  = \frac{2\,M^2 M_{\mathrm{Pl}}^2}
  {3\,\alpha\,\omega_0^2},
  \qquad
  \beta_{\mathrm{fit}}\,\Phi_0^2
  = q_{\mathrm{fit}} - 1 = 0.853.
  \label{eq:amplitude_constraint}
\end{equation}

\textbf{Fitted target:} if the correction coefficient is
read as an effective parameter, then the peak dimensionless
oscillon amplitude must satisfy
$\Phi_0 = \sqrt{0.853/\beta_{\mathrm{fit}}}$ in order to
close exactly on $q_{\mathrm{fit}}$.  For
$\beta_{\mathrm{fit}} = 1$ this gives
$\Phi_0 \approx 0.924$, already in the nonlinear regime
where both the $G_2(X)$ and metric effects are
$\mathcal{O}(1)$.  The actual $\Phi_0$ and $\beta$ must be
fixed by the full oscillon calculation, not by this
calibration identity alone.

\begin{center}
\fbox{\parbox{0.88\textwidth}{\centering
\textbf{Key structural result:}
The $G_2(X)$ extension of ISPG predicts
$q = 1$ from \emph{first principles} at leading
order (no free parameters). The bi-conformal
metric backreaction, spatial gradients, and
harmonic coupling each contribute corrections
of order $\Phi_0^2$, shifting $q$ toward the
fitted reference value $q_{\mathrm{fit}}=1.853$.
The present semi-analytic bundle lands nearby, while
the coefficient $\beta$ and the oscillon amplitude $\Phi_0$
are
\emph{independently determinable} by lattice
simulation of the coupled $G_2(X)$ +
bi-conformal system, providing a
\textbf{falsifiable} test of the entire framework.
}}
\end{center}

\medskip
\textit{Numerical verification.}---Using the 3D spherical
oscillon profile at $\Phi_0 = 2.35$, $\alpha = 1/2$ (the
background established in the Quantum paper, Sec.~7), a
shooting computation yields $\Omega = 0.86594$
($< 0.008\%$ from the reference value $0.866$).
The three correction channels evaluate to:
\begin{equation}
  \delta q_{\mathrm{metric}} = \frac{I_2(2\Phi_0)}
    {I_0(2\Phi_0)} = 0.623,
  \qquad
  \delta q_{\mathrm{grad}} = 0.177,
  \qquad
  \delta q_{\mathrm{harm}} = 0.107,
  \label{eq:dq_numerical}
\end{equation}
giving a total $q_{\mathrm{semi}} = 1 + 0.907 = 1.907$,
within $2.9\%$ of the fitted reference value
$q_{\mathrm{fit}} = 1.853$. The dominant correction
(metric backreaction, $73\%$ of $\delta q$) is computed
exactly; the remainder comes from volume-averaged gradient
and order-of-magnitude harmonic estimates. This is therefore
near-support for the mechanism, not an exact derivation of
$q_{\mathrm{fit}}$. At $q = 1.907$
the muon/electron mass ratio is
$b_{72}(q)/b_5(q) = 206.73$, matching the observed
$206.77$ to $0.02\%$. The residual $2.9\%$ gap in $q$
maps to a $< 0.02\%$ shift in the mass ratio and is within
the expected precision of the semi-analytical gradient and
harmonic formulas; it remains an open target for full lattice
PDE evolution of the oscillon profile (future work).
The simulation code is provided as
\texttt{MASS/python/ispg\_oscillon\_simulation.py}.

\medskip
\textbf{Step 7: Why a common $q$ is plausible in the Mathieu ladder.}

The parameter $q$ depends on the background oscillon
properties ($\varepsilon$), not on the perturbation mode
number $N$. This mode-independence is suggested directly by the
structure of $G_2(X)$: the $X^2$ term couples to
\emph{all} perturbation wavelengths equally because $X$
is a scalar constructed from $\phi$ alone. In the
bi-conformal metric, the conformal factor $e^{\phi/2}$
acts identically on all length scales, so the same oscillon
background naturally leads to a mode-independent $q$ within
this one-dimensional Mathieu EFT limit.

This is numerically consistent with the fitted ladder
(\S\ref{sec:mathieu_numerical}): a single constant
$q_{\mathrm{fit}} = 1.853$ organizes the charged lepton
masses and, with re-optimized $N$ assignments, the heavy
particle table. This should not yet be read as a proof
that sector-specific effective $q$ values are excluded in
the still-unified 3D theory; that stronger universality
claim belongs to the later unification work.

\medskip
\textbf{Step 8: Physical interpretation.}

The $G_2(X)$ extension says: the ISPG medium has
\textbf{nonlinear elasticity}. When deformations are
small (weak gravitational fields, PPN regime), the
medium responds linearly: $\Box\phi = 0$. When
deformations are large (inside an oscillon), the
nonlinear $X^2$ term creates a self-confining
potential that traps resonant modes---these are
particles.

The two EFT parameters $\alpha$ and $M$ characterize
the medium's nonlinear response:
\begin{itemize}
  \item $\alpha$: the dimensionless strength of the
    nonlinearity.
  \item $M$: the energy scale at which the nonlinear
    regime activates.
\end{itemize}
From the oscillon amplitude \eqref{eq:amplitude_G2}
and the quantum model frequency
$\nu_0 = 1.91 \times 10^{30}$~Hz:
\begin{equation}
  \frac{\alpha}{M^2}
  = \frac{2 M_{\mathrm{Pl}}^2}
  {3\,\Phi_0^2\omega_0^2}.
  \label{eq:alpha_over_M2}
\end{equation}
With $\alpha \sim \mathcal{O}(1)$, the nonlinearity
scale is
\begin{equation}
  M \sim \Phi_0\omega_0 / M_{\mathrm{Pl}}
  \sim \sqrt{X_0}\,/\,M_{\mathrm{Pl}}.
\end{equation}
This is far below the Planck scale, consistent with the
$X^2$ term being an EFT correction that is irrelevant
in the gravitational sector but dominant inside oscillons.

\medskip
\textbf{Summary: the main chain.}

\begin{center}
\fbox{\parbox{0.90\textwidth}{
\begin{equation*}
  G_2(X) \;\xrightarrow[\text{self-consistency}]
  {\text{oscillon}}\;
  \varepsilon = \tfrac{4}{3}
  \;\xrightarrow[\text{flat space}]{\text{Mathieu}}\;
  q = 1
  \;\xrightarrow{+\;\delta q_{\mathrm{bi\text{-}conf}}}
  \;q_{\mathrm{semi}}\simeq 1.907
  \;\approx\;q_{\mathrm{fit}} = 1.853
  \;\xrightarrow{\;b_N(q)\;}\;
  m_N = m_1\,b_N
\end{equation*}

\smallskip
\centering
\textbf{One field} $\phi$, \textbf{one} EFT extension
$G_2(X)$ (dual role: ghost-free UV completion
+ oscillon nonlinearity) $\to$ \textbf{all}
particle masses.\\
Shift symmetry: \textbf{preserved}.
New fields: \textbf{none}.
Weak-field gravity: \textbf{unchanged}.\\
$\delta q_{\mathrm{bi\text{-}conf}} = 0.853$:
fitted target for metric backreaction + gradients + harmonics.
The present semi-analytic bundle gives
$\delta q_{\mathrm{semi}}\simeq 0.907$; exact closure is
lattice-determinable future work.
}}
\end{center}

\subsection{Summary of Results}

\begin{center}
\renewcommand{\arraystretch}{1.3}
\begin{tabular}{@{}p{6cm}ll@{}}
\textbf{Result} & \textbf{Accuracy} & \textbf{Ref.}
  \\[2pt]\hline\\[-6pt]
Fitted Mathieu ladder scan: $m_N = m_1 b_N(q)$
  & $10/12<1\%$, $11/12<5\%$; light-quark scheme caveat
  & \S\ref{sec:numerical_results} \\
Muon mass ($q_{\mathrm{fit}} = 1.853$)
  & $< 0.001\%$ & \S\ref{sec:mathieu_numerical} \\
Selected heavy entries ($b,t,W,Z,H$)
  & nearest-$N$ scan; case-by-case & \S\ref{sec:mathieu_numerical} \\
Koide formula (inherited from individually-fit lepton masses;
\emph{not} an independent derivation --- see QUANTUM Sec.~7 for
the rank-1/E1D consistency argument on the MASS/Mathieu
ladder; the formal 3D-to-1D bridge to first-principles
charged-lepton masses remains open)
  & $0.028\%$ & \S\ref{sec:mathieu_numerical} \\
Electron---first fully stable charged Mathieu mode
  & adopted gap threshold & \S\ref{sec:mathieu_numerical} \\
$q = 1$ from $G_2(X)$ self-consistency (leading order)
  & analytic & \S\ref{sec:q_derivation} \\
$q_{\mathrm{semi}}\simeq 1.907$ from bi-conformal corrections
  & near-support; open closure & \S\ref{sec:q_derivation} \\
Common $q$ in the fitted Mathieu ladder
  & numeric; sector universality open & \S\ref{sec:mathieu_numerical} \\
$G_2(X)$ dual role: UV completion + mass generation
  & structural & \S\ref{sec:q_derivation} \\
\end{tabular}
\end{center}

\textit{Epistemic status of the table below.}---At fixed
$q_{\mathrm{fit}} = 1.853$, the integer $N$ associated with each particle
is determined by selecting the nearest integer whose
Mathieu characteristic value reproduces the observed mass
($\Delta N \approx 1.5N/1000$).  The $<\!0.3\%$ per-particle
residuals quoted below therefore represent a fit in which
one integer degree of freedom per particle is consumed, on
top of the single global parameter $q$: given the observed
masses, the nearest-$N$ procedure yields these values by
construction.  The framework does not predict which $N$
each species must take; it predicts that charged-sector
masses fall on the $b_N(q)$ ladder for some $q$ close to
$1.853$, and the empirical monotonicity $N(\text{species})
\uparrow$ with mass is the non-trivial content.

Reoptimized $N$ assignments
(at $q_{\mathrm{fit}} = 1.853$, $\Delta N \approx 1.5N/1000$):
\begin{center}
\renewcommand{\arraystretch}{1.2}
\begin{tabular}{@{}lrrrr@{}}
Particle & $N$ & $m_{\mathrm{pred}}$ (MeV)
  & $m_{\mathrm{exp}}$ (MeV) & Error
  \\[2pt]\hline\\[-6pt]
$e$   & 5    & 0.511     & 0.511     & exact \\
$\mu$ & 72   & 105.659   & 105.658   & $0.001\%$ \\
$\tau$& 295  & 1773.7    & 1776.9    & $0.177\%$ \\
$c$   & 250  & 1273.9    & 1273.0    & $0.067\%$ \\
$b$   & 453  & 4182.5    & 4183.0    & $0.012\%$ \\
$W$   & 1986 & 80389.1   & 80369.2   & $0.025\%$ \\
$Z$   & 2115 & 91171.6   & 91188.0   & $0.018\%$ \\
$H$   & 2479 & 125254.0  & 125250.0  & $0.003\%$ \\
$t$   & 2910 & 172593.6  & 172570.0  & $0.014\%$ \\
\end{tabular}
\end{center}

\subsection{Testable Predictions}
\label{sec:N4_predictions}

The model makes the following \textbf{experimentally verifiable}
predictions:

\begin{enumerate}
  \item \textbf{$N = 4$ resonance:
    $m_4 = m_1 \times b_4(1.853) \approx 329$~keV.}
    The Mathieu instability zone width at this index is
    $\sim 10^{-2}$, giving a gap-threshold lifetime estimate
    $\tau \lesssim 10^{-20}$~s (order-of-magnitude only;
    the gap-to-decay-rate map is not first-principles).
    Possible search channels:
    \begin{itemize}
      \item Resonance peak in $e^+e^-$ scattering at
        $\sqrt{s}\approx 329$~keV (low-energy collider
        regime).
      \item Hard-X / soft-gamma photon line at
        $m_4/2\approx 165$~keV from a radiative decay
        channel, if such a channel exists; instrument
        matching requires the $\sim 100$--$300$~keV band
        (e.g. INTEGRAL/SPI), \emph{not} the $3$--$20$~keV
        NuSTAR sterile-neutrino-DM band.
    \end{itemize}
    A KATRIN-type beta-spectrum search is \emph{not} accessible
    here: the tritium beta endpoint is $E_0\approx18.57$~keV,
    well below $329$~keV, so a $329$~keV kink is kinematically
    impossible in tritium decay. A beta source with endpoint
    $\gg329$~keV would be required.
    A $\tau \lesssim 10^{-20}$~s state cannot simultaneously
    serve as a cosmological dark-matter candidate; the
    DM-candidate language is dropped here.

  \item \textbf{$N = 3$ ultra-short resonance:
    $m_3 \approx 186$~keV.}
    Instability zone $0.185$---practically
    instantaneous decay. May appear as a virtual state
    in electron--positron annihilation channels.

  \item \textbf{Nonlinearity scale $M$ of $G_2(X)$.}
    The EFT scale $M$ at which the $X^2$ term
    activates determines the oscillon core size.
    If $M$ can be constrained independently
    (e.g.\ from primordial gravitational wave
    spectrum or scalar radiation in neutron star
    mergers), this provides a direct test of the
    $G_2(X)$ extension.

  \item \textbf{Light quark mass verification.}
    The assignments $N=10$ and $N=15$ give the paper-target
    values $m_u\approx 2.04$~MeV and
    $m_d\approx 4.59$~MeV. The live PDG-style central-mass
    scan with $\overline{\mathrm{MS}}$ running masses at
    $2$~GeV gives a $u$ residual $5.62\%$ and a $d$
    residual $1.80\%$. A consistent comparison requires an
    explicit renormalization scheme and scale, since
    light-quark masses are running parameters rather than
    pole masses; this remains a concrete lattice/PDG-scheme
    test, not a closed percent-level prediction.

  \item \textbf{Cosmological frequency
    $\nu_0 = 1.91 \times 10^{30}$~Hz.}
    This is the calibrated value of the medium's fundamental
    proper-time carrier frequency.
    Indirect test:
    \begin{itemize}
      \item Primordial gravitational wave spectrum---
        a cutoff or characteristic structure is expected
        at wavelengths associated with $\nu_0$.
      \item Cosmological redshift anisotropy---the
        invariant frequency model predicts specific
        correlations in the CMB.
    \end{itemize}

  \item \textbf{k-Oscillon in $G_2(X)$:
    closing $q_{\mathrm{semi}}$ onto
    $q_{\mathrm{fit}} = 1.853$.}
    Lattice simulation of k-oscillons with
    $G_2(X) = -M_{\mathrm{Pl}}^2 X/2
    + \alpha X^2/M^2$ on the bi-conformal background
    determines $\beta$ and $\Phi_0$ independently
    (Eq.~\ref{eq:amplitude_constraint}).
    Three independent checks:
    \begin{itemize}
      \item Leading order: $q = 1$ from
        self-consistency alone.
      \item Corrections:
        $\delta q_{\mathrm{metric}}$ from
        $I_2(2\Phi_0)/I_0(2\Phi_0)$,
        $\delta q_{\mathrm{grad}}$ from the
        $\mathrm{sech}^2$ profile,
        $\delta q_{\mathrm{harm}}$ from cubic
        coupling---the present semi-analytic bundle gives
        $\delta q\simeq 0.907$ and must be tested against
        the fitted target $\delta q_{\mathrm{fit}}=0.853$.
      \item Dual role of $G_2(X)$: the same
        extension that provides ghost-free UV
        completion (Appendix~11) must also produce
        the correct $q$.
    \end{itemize}
\end{enumerate}

\section{Ontological Synthesis: Vibration, Pressure, and Gravity}
\label{sec:ontology}

The results of this paper, combined with the Bernoulli
pressure-deficit mechanism (Quantum paper, \S3) and the
$G_2(X)$ extension (\S\ref{sec:q_derivation}), close a
conceptual loop that identifies the fundamental ontology
of the ISPG framework. The entire construction rests on
a single empirical input: the fundamental frequency
$\nu_0 \approx 1.91 \times 10^{30}$~Hz, whose value is
calibrated in the MASS/Mathieu spectrum
(\S\ref{sec:numerical_results}).
\emph{What determines $\nu_0$ from first principles}
remains an open question; it may ultimately follow from
the discrete micro-structure of the medium (see Future
Directions).

\subsection{Pressure Is Vibrational Energy}
\label{sec:pressure_vibration}

In kinetic theory, gas pressure is the macroscopic
manifestation of molecular vibration:
$P = \frac{1}{3}\rho\langle v^2\rangle$.
Temperature measures the vibration intensity.

In ISPG the identification is exact. The kinetic
variable
\begin{equation}
  X \equiv -\tfrac{1}{2}\,g^{\mu\nu}
  \partial_\mu\phi\,\partial_\nu\phi
\end{equation}
is the \textbf{vibrational energy density} of the
scalar field $\phi$ oscillating at the fundamental
frequency $\nu_0$. The scalar-sector Lagrangian
\begin{equation}
  \mathcal{L}_\phi = G_2(X)
  = -\frac{M_{\mathrm{Pl}}^2}{2}\,X
  + \frac{\alpha}{M^2}\,X^2
  \label{eq:pressure_is_G2}
\end{equation}
is simultaneously the \textbf{pressure equation of state}
of the medium, expressed as a function of its vibrational
energy:
\begin{itemize}
  \item The linear term
    $-\frac{M_{\mathrm{Pl}}^2}{2}\,X$:
    pressure from weak vibrations
    (linear elasticity $\to$ ordinary gravity).
  \item The quadratic term
    $\frac{\alpha}{M^2}\,X^2$:
    pressure from strong vibrations
    (nonlinear elasticity $\to$ oscillon
    self-trapping $\to$ particle masses).
\end{itemize}

\noindent
\begin{center}
\fbox{\parbox{0.85\textwidth}{\centering
\textbf{Ontological identification:}\\[3pt]
Pressure $\equiv$ vibrational energy density
$\equiv$ $G_2(X)$.\\[2pt]
The medium's ``pressure'' is not an abstract
thermodynamic quantity---it \emph{is} the energy
of the $\phi$-field vibrating at $\nu_0$.
}}
\end{center}

\subsection{The Frequency Cascade}
\label{sec:cascade}

All gravitational phenomena in ISPG arise through a
\textbf{nonlinear frequency cascade}---a downward
transfer of vibrational energy from the fundamental
frequency to zero frequency:

\begin{equation}
  \boxed{
  \underbrace{\nu_0}_{\substack{\text{medium}\\\text{fundamental}\\
  \sim 10^{30}\;\mathrm{Hz}}}
  \;\xrightarrow[\;G_2(X)\;\text{nonlinearity}\;]
  {\text{oscillon formation}}\;
  \underbrace{\nu_N = \frac{N}{m}\,\nu_0}_{\substack{
  \text{particle resonance}\\
  \sim 10^{20}\;\mathrm{Hz}}}
  \;\xrightarrow[\;\text{time-averaging}\;]
  {\langle\phi^2\rangle \neq 0}\;
  \underbrace{\omega = 0}_{\substack{
  \text{static }\phi(r)\\
  = \text{gravity}}}
  }
  \label{eq:cascade}
\end{equation}

\begin{enumerate}
  \item \textbf{Stage 1: $\nu_0 \to \nu_N$.}
    The medium vibrates at $\nu_0$.
    The $X^2$ nonlinearity in $G_2(X)$ amplifies
    small inhomogeneities, causing the uniform
    vibration to condense into localized
    resonances---oscillons. Each oscillon vibrates
    at a harmonic $\nu_N$ determined by the Mathieu
    stability zones (\S\ref{sec:mathieu_numerical}).
    The oscillon's mass is
    $m_N c^2 = h\nu_N \times b_N(q)$.

  \item \textbf{Stage 2: $\nu_N \to \omega = 0$.}
    The oscillon is a nonlinear object: its
    $\cos(\omega_N t)$ oscillation, passed through
    the $X^2$ and $e^{-2\phi}$ nonlinearities,
    generates a \emph{time-averaged} (DC) component
    $\langle\phi\rangle \neq 0$ that extends as
    $\phi(r) \sim -GM/(rc^2)$ far from the core.
    This static profile \emph{is} the Newtonian
    gravitational potential.

  \item \textbf{Stage 3: $\nabla\phi \neq 0
    \to$ force.}
    The gradient of the static profile creates
    a spatial variation in the medium's pressure:
    $\nabla P = (dG_2/dX)\,\nabla X \neq 0$.
    This pressure gradient is the gravitational
    force---the Bernoulli mechanism of the
    Quantum paper.
\end{enumerate}

The cascade explains \textbf{why uniform vibration
does not gravitate}: a spatially constant $X$
has $\nabla X = 0$, hence $\nabla P = 0$, hence
no force---regardless of how large the vibrational
amplitude is. Gravity requires
\textbf{inhomogeneity}. Two sources of
inhomogeneity operate in ISPG:
\begin{itemize}
  \item \textbf{Microscopic (local Bernoulli channel):}
    an oscillon, once formed, is a spatially localized
    object whose field amplitude falls off as a
    standing-wave tail.  In the tail region the
    gradient kinetic energy density
    $|\nabla\varphi|^2 \neq 0$; by the Bernoulli
    identity (Quantum paper, \S3) this kinetic energy
    comes at the expense of the medium's static
    pressure.  The result is a pressure deficit that
    extends outward as $\sim 1/r$ and whose gradient
    is the Newtonian gravitational field~$g_N$.
    (The \emph{formation} of the oscillon---the
    cascade $\nu_0 \to \nu_N$---is a separate,
    prior step; what produces gravity is the
    \emph{spatial structure} of the already-formed
    resonance.)
  \item \textbf{Macroscopic (transported channel):}
    cosmological flows of the $\phi$-medium form
    vortices (via Kelvin--Helmholtz and related
    instabilities; see Appendix~9, MOND).
    Centrifugal acceleration physically transports
    the vibrating substance outward, depleting the
    centre---a rarefaction by mass transport rather
    than energy conversion. This produces the
    additional gravitational gradient $g_h$ responsible
    for MOND phenomenology.
\end{itemize}
Both mechanisms reduce the medium's background
vibrational energy at a given location---one by
\emph{mode conversion}, the other by \emph{physical
transport}---but the outcome is the same: a
pressure deficit that acts as gravity.

\subsection{Spontaneous Symmetry Breaking:
  From Uniform Vibration to Matter}
\label{sec:spontaneous}

The uniform $\nu_0$-vibration is a
\textbf{metastable} state. The $X^2$ term in
$G_2(X)$ introduces a nonlinear instability
analogous to the Jeans instability in
self-gravitating fluids:

\begin{enumerate}
  \item \textbf{Linear regime ($X^2/M^2 \ll X$):}
    small perturbations $\delta\phi$ propagate
    as waves satisfying $\Box\phi = 0$.
    The background is stable.

  \item \textbf{Nonlinear regime
    ($X^2/M^2 \sim X$):}
    perturbations with amplitude exceeding a
    critical threshold $\Phi_c$ experience
    self-focusing---the $X^2$ term increases
    the local energy density, which further
    increases $X$, which further amplifies
    the perturbation. This positive feedback
    concentrates energy into a localized region.

  \item \textbf{Saturation:}
    the self-consistency condition
    \eqref{eq:amplitude_G2} halts the collapse
    at the oscillon amplitude
    $\Phi_0 = \sqrt{2M^2 M_{\mathrm{Pl}}^2
    / (3\alpha\omega_0^2)}$.
    The result is a stable, localized,
    oscillating bound state---a \textbf{particle}.
\end{enumerate}

This is the ISPG mechanism for
\textbf{matter creation}: the uniform vibration
of the medium spontaneously breaks translational
symmetry by condensing into localized oscillons.
No external ``seed'' is required---quantum
fluctuations in the $\phi$-field provide the
initial perturbation (as detailed in the Quantum
paper's nucleation mechanism).

\subsection{Dark Energy as Accumulated Echo Background}
\label{sec:dark_energy}

The Quantum paper (Sec.~8) identifies the
stochastic background vibration
$\langle\delta\varphi^2\rangle$ as the source
of dark energy. The frequency cascade provides
a complementary perspective:

\begin{itemize}
  \item \textbf{Before condensation:} the medium's
    full vibrational energy $E_{\nu_0}$ is
    distributed nearly uniformly at frequency $\nu_0$.
    No particles exist. No \emph{localized}
    (Newtonian) gravity exists---only large-scale
    pressure gradients from primordial
    inhomogeneities.

  \item \textbf{During condensation:} the $G_2(X)$
    instability converts a fraction of
    $E_{\nu_0}$ into oscillons (matter) and
    their gravitational tails
    ($\phi(r) \sim 1/r$). Two competing energy
    flows coexist:
    \begin{enumerate}
      \item \textbf{Depletion:} the $\nu_0$-mode
        background loses energy to oscillon
        formation (the guitar string quietens).
      \item \textbf{Feedback:} every oscillon
        continuously broadcasts resonant-tail
        echoes back into the medium, feeding a
        stochastic background vibration
        $\langle\delta\varphi^2\rangle$ that
        accumulates globally
        (Quantum paper, \S8).
    \end{enumerate}

  \item \textbf{Late universe:} the Quantum paper
    gives a conditional perturbative-branch estimate in
    which the net balance is dominated by the feedback
    channel: the echo-background energy density
    $\varepsilon(t)$ (= the modeled dark-energy component) grows
    \emph{logarithmically}
    ($\varepsilon \sim \varepsilon_0 + A\ln(t/t_0)$),
    making the effective equation of state
    \textbf{slightly phantom}:
    \begin{equation}
      w(z) = -1 - \delta w_0\,
      \frac{(1+z)^{3}}{E(z)},
      \qquad \delta w_0 \sim 10^{-2}\text{--}10^{-1}.
      \label{eq:w_dark}
    \end{equation}
    The deviation $|\delta w|$ was larger at
    early times (more particles, stronger
    emission) and shrinks toward zero at late
    times: the dark energy asymptotically
    \textbf{approaches $\Lambda$ from the phantom
    side} ($w < -1$).
\end{itemize}

\textbf{Guitar-string analogy, refined.}
The guitar string loses energy (amplitude
decreases) and the note fades. But
in the ISPG universe the ``sound'' radiated
by matter (resonant-tail echoes) is itself
a form of background vibration that
accumulates in the medium. The string quietens
in one mode ($\nu_0$), while the acoustic
energy in the room (stochastic
$\langle\delta\varphi^2\rangle$)
slowly builds. The net dark-energy density
(echo background) grows, but the growth decelerates via
negative feedback (Quantum paper, \S8.2):
lighter particles emit less, so the process
is self-regulating.

\medskip
\textbf{Connection to DESI (2024--2025):}
The DESI baryon acoustic oscillation data
(DR1 at $2.6\sigma$; DR2 at $3.1\sigma$,
rising to $4.2\sigma$ with DES-SN5YR)
favor time-varying dark energy over
$\Lambda$CDM.  The source-driven ISPG branch gives
$w_0<-1$ and a small negative CPL slope $w_a<0$.
DESI DR2 therefore provides a mixed comparison rather than
a confirmation: the sign of $w_a$ is aligned, while the
central CPL sign of $w_0+1$ is opposite.  The illustrative
values $w_0 \approx -1.06$ and $w_a \approx -0.15$
should be read as a conditional branch estimate whose
magnitude inherits the tail-power normalization discussed
in the companion quantum analysis.

\subsection{The Closed Ontological Loop}
\label{sec:closed_loop}

The full ISPG ontology forms a self-consistent
loop---not a one-way chain but a true cycle,
because gravity feeds back into the vibration:

\begin{center}
\fbox{\parbox{0.92\textwidth}{
\begin{equation*}
  \underset{\text{vibration}}{\nu_0}
  \;\xrightarrow{\;P = G_2(X)\;}
  \underset{\text{pressure}}{\mathcal{L}_\phi}
  \;\xrightarrow{\;e^{\pm\phi}\;}
  \underset{\text{geometry}}{ds^2}
  \;\xrightarrow{\;\nabla\phi\;}
  \underset{\text{force}}{F}
\end{equation*}

\smallskip
\centering
\textit{Feedback:} the metric $e^{\pm\phi}$
modifies the wave equation
$\Box\phi$ (bi-conformal backreaction),
which shifts the oscillon spectrum and
the tail-emission power---closing the
cycle (Quantum paper, \S8.2).

\smallskip
\centering
Three geometries of rarefaction, one principle
($\nabla P \neq 0 \;\Rightarrow\;$ force):

\medskip
\renewcommand{\arraystretch}{1.35}
\begin{tabular}{@{}r@{\;:\;\;}p{5.2cm}@{\;\;$\to$\;\;}p{4.0cm}@{}}
\textbf{1.\ Gravity}
  & oscillon tail $|\nabla\varphi|^2 \neq 0$
    $\to$ static pressure drops
  & Bernoulli pressure deficit
    $= g_N$ \\[2pt]
  & \multicolumn{2}{@{}l@{}}{\small\itshape
    (kinetic energy of the standing-wave tail
    reduces the medium's static pressure)} \\[4pt]
\textbf{2.\ MOND}
  & vortex rotation pushes substance
    outward centrifugally
  & central rarefaction
    $= g_h$ \\[2pt]
  & \multicolumn{2}{@{}l@{}}{\small\itshape
    (physical transport---vibrating medium
    expelled from centre)} \\[4pt]
\textbf{3.\ Inertia}
  & acceleration $\to$ retarded self-field
    $\to$ fore--aft $\delta\phi$
  & pressure asymmetry
    $= -m\mathbf{a}$ \\[2pt]
  & \multicolumn{2}{@{}l@{}}{\small\itshape
    (temporal retardation---Bernoulli deficit
    not yet formed ahead, still present behind)}
\end{tabular}

\medskip
\textbf{Everything is one field, one vibration,
one equation of state.}\\[3pt]
Vibration $=$ pressure $=$ geometry $=$ gravity.\\
Matter $=$ nonlinear self-trapping of vibration.\\
Inertia $=$ retardation asymmetry of the same
pressure deficit (Appendix~17; derived as
$F = -m_{\mathrm{eff}}\,a$ in
ISPG\_Quantum, Sec.~5).\\
MOND $=$ macroscopic rarefaction of the same
medium by rotation or convection (Appendix~9).\\
Dark energy $=$ accumulated resonant-tail echoes
(growing logarithmically, $w < -1$ in the
conditional quantum-sector estimate;
ISPG\_Quantum, Sec.~6).\\
Dark-energy scale $=$ perturbative-branch
negative-feedback estimate, conditionally yielding
$\rho_{\mathrm{vac}} \sim 10^{-121}\,
M_{\mathrm{Pl}}^4$ while zero-point decoupling and
non-perturbative normalization remain open
(ISPG\_Quantum, Sec.~6.3).\\
Spin-1/2 structural route $=$ vortex winding number
$n = \pm 1$, $J = n\hbar/2$ under the conditional spin
bridge; Fermi--Dirac statistics and Pauli exclusion as a
candidate Berry-phase/vortex-topology route
(ISPG\_Quantum, Sec.~10).\\
Gravitational decoherence $=$
$\Gamma_{\mathrm{dec}} = 2m_0|\Delta\Phi_N|/h$,
testable at $\sim 0.5$~Hz for eV-scale photons
(ISPG\_Quantum, Sec.~9).
\\
Neutrinos $=$ polarisation-rotation waves
(helical vortices, $n = \pm 1$) of the
$\nu_0$-medium (Sec.~\ref{sec:neutrinos}).\\
Weak force $=$ nonlinear torsion--transverse
coupling (\S\ref{sec:neutrinos}).\\
Gravitational waves $=$ free transverse-traceless
modes of the same medium; photon identification is
deferred to a companion paper
(Sec.~\ref{sec:photon}).\\
Gluons $=$ transverse waves confined by
$\nu_N/\nu_0$ frequency mismatch; flux tube
$=$ forced-frequency bridge, $V \propto L$
(Sec.~\ref{sec:gluon}).\\
Three dimensions $=$ minimum for
$\ell=0$ (scalar) $+$ $\ell=2$ (tensor)
coexistence; 1D substrate $\to$ 3D by
vibrational necessity
(Sec.~\ref{sec:dimensions}).\\
Conditional Higgs scaling:
$v(\varphi)=v_0\,e^{\varphi/2}$; the oscillon
amplitude $\Phi_0$ is the candidate ISPG order
parameter, with the $\Phi_0\to v_{\mathrm{EW}}$
normalization left open (Sec.~\ref{sec:higgs};
ISPG\_Quantum, Sec.~8).
}}
\end{center}

%===================================================================
\section{Neutrinos as Polarisation-Rotation Waves}
\label{sec:neutrinos}
%===================================================================

The Mathieu mass formula $m_N = m_1\,b_N(q)$
organizes the massive charged-sector spectrum as a
nearest-$N$ empirical ladder
(\S\ref{sec:mathieu_numerical}), but it assigns
integer harmonic numbers $N = 1, 2, 3,\ldots$\
to its solutions.  Neutrino masses
($m_\nu \lesssim 0.1$~eV) are a factor
$\sim 10^{6}$ below the electron mass,
corresponding to $N \ll 1$---outside the
integer-harmonic tower.  This section proposes
a resolution that preserves the single-field
ontology and identifies neutrinos as
the \textbf{fourth and final vibration mode}
of the $\nu_0$-medium: \emph{polarisation
rotation}---the progressive turning of a
transverse wave's oscillation plane as it
propagates.

The key insight is that ``torsion'' in the
ISPG context does not require the substrate
to possess physical thickness, nor does it
require the scalar field to carry an intrinsic
polarisation vector (which it does not: a real
scalar $\phi(x)$ has no internal polarisation
degrees of freedom).  What is being rotated is
not a polarisation plane in the
vector-field sense but the \emph{spatial
argument} of a cylindrically-symmetric scalar
mode about the propagation axis.  In an
expansion of $\phi$ in cylindrical coordinates
$(r_\perp,\vartheta,z)$, the nonlinearity
$G_2(X)$ admits a helical standing-wave
solution
\begin{equation}\label{eq:neutrino_helical}
  \delta\phi_\pm(r_\perp,\vartheta,z,t)
  = A(r_\perp)\,
    \cos\!\bigl[kz - \omega t \pm \vartheta\bigr],
\end{equation}
i.e.\ the $\ell = \pm 1$ \emph{azimuthal-winding}
mode of the scalar substrate, whose constant-phase
surfaces form a helix.  Viewed in the
field-configuration picture, this is precisely
a vortex line with winding number
$n = \pm 1$ (\S\ref{sec:neutrino_vortex};
cf.\ the complementary treatment of scalar
vortices in the companion quantum
paper~\cite{Lotishvili2026Quantum}).  The sense
of the winding (clockwise vs.\ counterclockwise
around the axis) furnishes a geometric
chirality, and the vortex topology furnishes a
half-integer phase under $2\pi$ rotation of
the surrounding field configuration.

\paragraph{Status of the ``polarisation-rotation'' label.}
In the remainder of this section, ``polarisation
rotation'' is retained only as a \emph{visual
metaphor} for the helical-argument mode of
Eq.~\eqref{eq:neutrino_helical}; it does not
denote the rotation of an intrinsic
polarisation vector in the sense of optics.
The optical analogues of Faraday rotation and
optical activity---where a genuine
polarisation vector of $A_\mu$ rotates---apply
to vector-field media; their role here is
heuristic, not identificational.  The formal
object is the cylindrical $\ell = \pm 1$ mode,
and all quantitative statements below are
anchored to that structure.

%-------------------------------------------------------------------
\subsection{Four modes of a wave-bearing medium}
\label{sec:four_modes}
%-------------------------------------------------------------------

The scalar field $\phi$ in three emergent spatial
dimensions (\S\ref{sec:dimensions}) supports
perturbations characterised by their transformation
under rotations about the propagation axis.
In the $\nu_0$-medium there are exactly
\textbf{four} independent perturbation classes:

\begin{enumerate}
  \item \textbf{Longitudinal (compression):}
    the medium compresses and rarefies along the
    propagation direction.  This is the $\ell = 0$
    scalar mode $\to$ \textbf{gravity}.
  \item \textbf{Transverse (displacement):}
    the field oscillates perpendicular to the
    propagation axis.  This is the $\ell = 1$
    dipole mode $\to$ \textbf{electromagnetism}
    (\S\ref{sec:charge}).
  \item \textbf{Membrane (surface deformation):}
    higher multipole deformations of the
    field profile.  These are the $\ell \geq 2$
    tensor modes $\to$ \textbf{gluons} (confined;
    \S\ref{sec:gluon}).
  \item \textbf{Helical azimuthal winding ($\ell=\pm 1$
    cylindrical mode):}
    a cylindrically-symmetric scalar mode whose
    phase winds by $\pm 2\pi$ around the
    propagation axis per wavelength, producing
    a helical constant-phase surface.  This is
    not a polarisation vector rotating (the
    scalar has no such vector) but an azimuthal
    winding of the scalar argument itself, with
    topological charge $n = \pm 1$.  We identify
    this mode with the \textbf{neutrino} and the
    \textbf{weak interaction}.
\end{enumerate}

No fifth independent perturbation class exists in
a scalar-field medium with two transverse
dimensions.  This is the ISPG analogue of the
Standard Model's four fundamental interactions:
\textbf{each force is a perturbation mode of the
same field.}

%-------------------------------------------------------------------
\subsection{Why polarisation rotation is the
neutrino}
\label{sec:torsion_neutrino}
%-------------------------------------------------------------------

The polarisation-rotation mode accounts for every
distinctive property of neutrinos that the Standard
Model treats as independent postulates:

\begin{enumerate}
  \item \textbf{Intrinsic chirality.}\;
    Polarisation rotation has a definite
    handedness---the plane turns clockwise or
    counterclockwise as the wave advances.
    This is a \emph{geometric} property of the
    helical wave itself, not an imposed quantum
    number.  Neutrinos are left-handed and
    antineutrinos right-handed because
    polarisation rotation \emph{is} chirality.
    No other perturbation mode possesses this
    property, explaining why only the weak
    interaction violates parity.

  \item \textbf{Electrical neutrality.}\;
    Polarisation rotation is symmetric in the
    transverse plane: the oscillation visits
    both orthogonal directions equally.  The
    time-averaged $\ell = 1$ dipole component
    is zero.  The neutrino carries no electric
    charge.

  \item \textbf{Extreme lightness.}\;
    Polarisation rotation requires coupling
    \emph{between} the two transverse
    polarisations.  In a perfectly isotropic
    linear medium this coupling is zero; it
    arises only through the $X^2$ nonlinearity
    of $G_2(X)$ acting on the oscillon
    background.
    The coupling strength is set by the
    background oscillation amplitude
    $\Phi_0$ and the transverse--longitudinal
    aspect ratio $R_\perp/\lambda_0$, yielding
    an effective torsional rigidity
    (\S\ref{sec:torsion_derivation}):
    \begin{equation}\label{eq:torsion_rigidity}
      \kappa_T \;\sim\; c_s^2\,
      \Bigl(\frac{R_\perp}{\lambda_0}\Bigr)^2
      \;\ll\; c_s^2,
    \end{equation}
    where $c_s$ is the sound speed
    and $\lambda_0 = c/\nu_0$.  The torsional
    mode frequency (and hence mass via
    $m = h\nu/c^2$) is suppressed by the same
    ratio, naturally placing neutrino masses in
    the sub-eV range without fine-tuning.

  \item \textbf{Extremely weak coupling.}\;
    Polarisation rotation neither compresses the
    medium (no direct gravitational coupling) nor
    creates a net transverse dipole (no
    electromagnetic coupling).  It interacts
    with the other modes only through
    \textbf{nonlinear cross-terms} in the
    $G_2(X)$ action---the same $X^2$ term that
    couples the two transverse polarisations to
    each other and to the longitudinal mode.
    These couplings are suppressed by
    $(R_\perp/\lambda_0)^2$, the same small
    parameter that suppresses the mass.
    This is the \textbf{weak interaction}:
    a nonlinear coupling between the
    polarisation-rotation sector and the
    scalar/vector sectors.

  \item \textbf{Three flavours (consistency check,
    not prediction).}\;
    The polarisation-rotation mode lives on the
    oscillon background, which acts as a finite
    cavity of radius $R_c$.  The cavity supports
    a discrete set of torsional standing waves
    with harmonic numbers $n = 1, 2, 3, \ldots$
    and mass eigenvalues
    $m_n \propto n^2 \omega_T$.
    The stability of the $n$-th harmonic is
    governed by the ratio of its decay width
    $\Gamma_n$ to its mass $m_n$, with the
    dominant channel
    $\nu_n \to \nu_{n'} + \gamma$ (torsional
    $\to$ lower torsional $+$ scalar/transverse
    radiation) scaling parametrically as
    \begin{equation}\label{eq:torsion_decay}
      \Gamma_n \;\propto\;
      \frac{\alpha^2}{M^4}\,\omega_T^5\,n^4\,
      (n^2 - n'^2)^3,
    \end{equation}
    where the $n^4$ comes from the azimuthal
    overlap integral (selection rule
    $\Delta n = \pm 1$) and $(n^2 - n'^2)^3$
    from the phase space.  The parameters
    $\alpha, M$ and the absolute
    normalisation depend on the oscillon
    profile and are not fixed at the
    effective-field-theory level used here.
    The structural statement of this bullet is
    therefore: the bound-state spectrum
    admits a finite number of stable
    torsional harmonics, with the critical
    harmonic number $n_{\mathrm{max}}$ set by
    the ratio $\Gamma_n/m_n$ computed on the
    full oscillon profile.  Whether
    $n_{\mathrm{max}}=3$, as observed
    phenomenologically, is a
    \emph{consistency check} to be performed
    by the lattice programme; it is not
    derived here, and the present paper does
    not claim to predict the number three from
    first principles.  Any quantitative
    lifetime estimates that would appear to
    single out $n=4$ as the boundary are
    deferred to that dedicated calculation.

  \item \textbf{Flavour oscillations.}\;
    The three torsional harmonics propagate in
    the same $\nu_0$-medium and share the same
    nonlinear background.  Their eigenstates in
    the free-propagation basis (mass eigenstates)
    differ from their coupling basis to charged
    leptons (flavour eigenstates), because the
    coupling occurs at the oscillon core where
    the $G_2(X)$ nonlinearity is strongest.
    The mismatch produces the PMNS mixing matrix
    and oscillation between flavours during
    propagation, exactly as observed.

  \item \textbf{Parity violation of the weak
    interaction.}\;
    Circular polarisation is a \emph{pseudoscalar}
    quantity: under spatial inversion
    ($\mathbf{x} \to -\mathbf{x}$), a left-handed
    helix becomes right-handed.  Any interaction
    mediated by polarisation-rotation modes
    therefore violates parity \emph{maximally}.
    The longitudinal, transverse (linear), and
    membrane modes are all parity-even, explaining
    why gravity, electromagnetism, and the strong
    force preserve parity while the weak force
    does not.
\end{enumerate}

%-------------------------------------------------------------------
\subsection{Formal derivation: the torsional mode
from $G_2(X)$}
\label{sec:torsion_derivation}
%-------------------------------------------------------------------

We now show that the polarisation-rotation mode
emerges directly from the $G_2(X)$ Lagrangian
without introducing any new field or degree of
freedom.

\textit{Setup.}---Consider perturbations of the
scalar field around the oscillon background
$\phi_{\mathrm{bg}}(t,r) = \Phi_0(r)\cos\omega_0 t$
(Eq.~\ref{eq:oscillon_ansatz}).  A general
transverse perturbation has two orthogonal
polarisations.  Writing the perturbation in
cylindrical coordinates $(r, \vartheta, z)$
centred on the oscillon core:
\begin{equation}\label{eq:helical_ansatz}
  \delta\phi
  = A(r)\cos(kz - \omega t + \vartheta)
  + B(r)\cos(kz - \omega t),
\end{equation}
where $A$ and $B$ are the amplitudes of the two
orthogonal transverse polarisations.  The ratio
$A/B$ determines the polarisation state.  For
$A = B$ with a $\pi/2$ phase shift, the
perturbation traces a \textbf{helix}:
\begin{equation}\label{eq:circular_pol}
  \delta\phi_{\pm}
  = A(r)\,e^{i(kz - \omega t \pm \vartheta)},
\end{equation}
where $+$ is left-circular and $-$ is
right-circular polarisation.  The azimuthal
factor $e^{\pm i\vartheta}$ is precisely the
winding-number-$(\pm 1)$ vortex phase of
ISPG\_Quantum, Sec.~10.

\textit{Mode coupling from $G_2(X)$.}---%
The kinetic variable for the perturbed field is
\begin{equation}
  X[\phi_{\mathrm{bg}} + \delta\phi]
  = X_{\mathrm{bg}}
  + \delta X^{(1)} + \delta X^{(2)} + \cdots,
\end{equation}
where $\delta X^{(1)}$ is linear in $\delta\phi$
and $\delta X^{(2)}$ is quadratic.  The
background contribution is
(cf.~Eq.~\ref{eq:X0_def}):
\begin{equation}\label{eq:Xbg_torsion}
  X_{\mathrm{bg}}
  = \tfrac{1}{2}\Phi_0^2\omega_0^2
  \sin^2(\omega_0 t).
\end{equation}
In the action $G_2(X)
= -\tfrac{M_{\mathrm{Pl}}^2}{2}X
+ \tfrac{\alpha}{M^2}X^2$,
the $X^2$ term evaluated on the perturbed field
generates a quadratic Lagrangian for $\delta\phi$
whose ``mass term'' oscillates in time:
\begin{equation}\label{eq:L2_pert}
  \mathcal{L}^{(2)}_{\delta\phi}
  = -\tfrac{M_{\mathrm{Pl}}^2}{2}\,
  (\partial_\mu\delta\phi)^2
  + \frac{2\alpha}{M^2}
  \Bigl[
    2\,X_{\mathrm{bg}}\,
    (\partial_\mu\delta\phi)^2
    + \bigl(
      \dot\phi_{\mathrm{bg}}\,\dot{\delta\phi}
      - \nabla\phi_{\mathrm{bg}}\!\cdot\!
      \nabla\delta\phi
    \bigr)^{\!2}
  \Bigr].
\end{equation}
The first bracket provides a \emph{global}
time-dependent effective speed, identical to
the longitudinal Mathieu driving of
\S\ref{sec:q_derivation}---this is
$(R_\perp/\lambda_0)^0$ and gives no polarisation
coupling.

\textit{Step 1: Decompose into polarisation
channels.}---Write the perturbation in the
circular basis:
\begin{equation}\label{eq:delta_phi_pm}
  \delta\phi = \delta\phi_+
  + \delta\phi_-,\qquad
  \delta\phi_\pm \propto
  e^{\pm i\vartheta}\,f_\pm(r,z,t).
\end{equation}
The second bracket in \eqref{eq:L2_pert}
contains the term
$(\nabla\phi_{\mathrm{bg}}\!\cdot\!
\nabla\delta\phi)^2$; because
$\phi_{\mathrm{bg}} = \Phi_0(r)\cos\omega_0 t$
depends only on $r$, its gradient is radial:
$\nabla\phi_{\mathrm{bg}} = \Phi_0'(r)\,
\hat{\mathbf{r}}\cos\omega_0 t$.
Acting on $\delta\phi_\pm \propto
e^{\pm i\vartheta}$:
\begin{equation}\label{eq:grad_coupling}
  \hat{\mathbf{r}}\cdot\nabla\delta\phi_\pm
  = \frac{\partial f_\pm}{\partial r}\,
  e^{\pm i\vartheta},
  \qquad
  \hat{\boldsymbol{\vartheta}}\cdot
  \nabla\delta\phi_\pm
  = \frac{\pm i f_\pm}{r}\,
  e^{\pm i\vartheta}.
\end{equation}
The \textbf{radial} part
($\partial_r f_\pm$) couples $\delta\phi_+$ and
$\delta\phi_-$ only through terms proportional
to $\Phi_0'(r)^2 / r^2$---the gradient of the
background oscillon profile evaluated at the
transverse scale $r \sim R_\perp$.

\textit{Step 2: Identify the
suppression factor.}---The background profile
$\Phi_0(r) = A\,\mathrm{sech}^2(r/R_c)$
has core radius $R_c \sim \hbar/(mc)$
(the Compton radius of the electron-mass
oscillon, Sec.~\ref{sec:q_derivation}).
For $r \gg R_c$ the profile decays
exponentially.  The polarisation-coupling terms
involve the \emph{azimuthal gradient}
$f_\pm / r$, which at the oscillon boundary
$r \sim R_c$ gives a coupling strength
$\sim \Phi_0 / R_c$.  Relative to the
longitudinal coupling $\sim \Phi_0\omega_0$,
this introduces the ratio
\begin{equation}\label{eq:Rperp_definition}
  \frac{R_\perp}{\lambda_0}
  \;\equiv\;
  \frac{1}{R_c\,\omega_0}
  = \frac{c}{R_c\,\omega_0\,c}
  = \frac{\lambda_0}{2\pi R_c},
\end{equation}
where $\lambda_0 = c/\nu_0
= \pi\ell_P \approx 5\times 10^{-35}$~m
and $R_c = \hbar/(m_e c)
= 3.86\times 10^{-13}$~m is the electron
Compton radius.

\textit{Step 3: Extract the torsional
equation.}---Define the slowly varying
polarisation angle
$\theta(z,t) \equiv
\tfrac{1}{2}\arg(f_+/f_-)$,
which measures the orientation of the
polarisation plane.  After projecting
$\mathcal{L}^{(2)}_{\delta\phi}$ onto the
$\theta$-sector (integrating out the radial
profile and the fast oscillation at $\omega_0$),
the effective equation is
\begin{equation}\label{eq:torsion_equation}
  \boxed{
  \frac{\partial^2\theta}{\partial t^2}
  - \kappa_T\,
  \frac{\partial^2\theta}{\partial z^2}
  + \Omega_T^2\bigl[1
  - 2q_T\cos(2\omega_0 t)\bigr]\theta
  = 0,
  }
\end{equation}
where the parameters are determined by
$G_2(X)$ and the oscillon profile:
\begin{itemize}
  \item $\kappa_T
  = c_s^2\,(R_\perp/\lambda_0)^2$
  is the torsional rigidity---the speed at
  which a polarisation twist propagates.
  The factor $(R_\perp/\lambda_0)^2$ arises
  because the azimuthal gradient
  \eqref{eq:grad_coupling} introduces one power
  of $1/(R_c\omega_0)$ per derivative; the
  equation is second-order, giving the square.
  \item $\Omega_T = \omega_0\,(R_\perp/\lambda_0)$
  is the torsional gap frequency.
  \item $q_T = q\,(R_\perp/\lambda_0)^2$ is the
  torsional Mathieu parameter---the same $q$ that
  governs the longitudinal mass spectrum, but
  suppressed by $(R_\perp/\lambda_0)^2$.
\end{itemize}

Eq.~\eqref{eq:torsion_equation} is again a
\textbf{Mathieu equation} for the torsional
sector.  Its stability bands determine the allowed
torsional masses, just as the longitudinal Mathieu
equation (\S\ref{sec:mathieu}) determines the
quark and charged-lepton masses.  The crucial
difference: all torsional parameters carry the
suppression factor
$(R_\perp/\lambda_{\mathrm{osc}})^2$, where
$\lambda_{\mathrm{osc}}$ is the Compton
wavelength of the oscillon on which the
torsional mode propagates; torsional masses
are suppressed by the same factor relative to
longitudinal masses:
\begin{equation}\label{eq:neutrino_mass_scale}
  \boxed{
  \frac{m_\nu}{m_e}
  \;\sim\;
  \Bigl(\frac{R_\perp}{\lambda_{\mathrm{osc}}}\Bigr)^2
  \,\eta_T ,
  }
\end{equation}
with $\eta_T$ the trapping fraction defined
below.  (Throughout the $(R_\perp/\lambda)^2$
expressions in Eqs.~\eqref{eq:omega_torsion}%
--\eqref{eq:torsion_equation}, the wavelength
$\lambda$ is to be read as
$\lambda_{\mathrm{osc}}$.)

\textit{Numerical evaluation.}---%
The longitudinal wavelength that is relevant
for the torsional suppression is \emph{not}
the carrier wavelength $\lambda_0 = c/\nu_0$
of the $\nu_0$-medium, but the Compton
wavelength of the specific oscillon on which
the torsional perturbation propagates:
\begin{equation}
  \lambda_{\mathrm{osc}}
  \;\equiv\; \hbar/(m_{\mathrm{osc}} c).
\end{equation}
The physical reason is that
Eq.~\eqref{eq:torsion_equation} describes a
perturbation \emph{on an oscillon background},
and the parametric driving in that equation is
set by the oscillon's own inverse Compton scale
$1/\lambda_{\mathrm{osc}}$, not by the medium
carrier scale $1/\lambda_0$: the carrier
oscillation enters only through the
\emph{driving amplitude} $q$, while the
\emph{length-scale ratio} that dresses the
azimuthal gradient in
Eq.~\eqref{eq:grad_coupling} is $R_\perp$
measured against the oscillon's own wavelength.
Consequently the physically relevant aspect
ratio is
\begin{equation}\label{eq:Rperp_corrected}
  \frac{R_\perp}{\lambda_{\mathrm{osc}}}
  = \frac{\lambda_{\mathrm{osc}}}{2\pi R_c}
  = \frac{\hbar/(m_{\mathrm{osc}} c)}
  {2\pi\,\hbar/(m_e c)}
  = \frac{m_e}{2\pi\,m_{\mathrm{osc}}},
\end{equation}
with $R_c = \hbar/(m_e c)$ the host-oscillon
Compton radius.
For $m_{\mathrm{osc}}\sim m_e$ (the lightest
charged fermion):
\begin{equation}\label{eq:Rperp_final}
  \frac{R_\perp}{\lambda_{\mathrm{osc}}}
  \sim \frac{1}{2\pi}
  \approx 0.16.
\end{equation}
The mass suppression from the aspect ratio alone is
$(R_\perp/\lambda_{\mathrm{osc}})^2
\sim 1/(2\pi)^2 \approx 0.025$, which accounts for
roughly four orders of magnitude of the
$m_\nu/m_e \sim 10^{-7}$ hierarchy.  The remaining
suppression arises because the torsional mode,
unlike the longitudinal mode, is not self-trapped:
it propagates \emph{away} from the oscillon core,
so the effective cavity size is not $R_c$ but
the much larger mean free path
$\ell_{\mathrm{mfp}}$ of the torsional wave
in the $\nu_0$-medium.  The trapping fraction
\begin{equation}\label{eq:trapping_fraction}
  \eta_T = \frac{R_c}{\ell_{\mathrm{mfp}}}
\end{equation}
provides the missing suppression.
The full neutrino mass formula is
$m_\nu \sim m_e \times (R_\perp/\lambda_{\mathrm{osc}})^2
\times \eta_T$, where $\eta_T$ is determined by the
nonlinear scattering cross-section of the torsional
mode on the $\nu_0$-background.

\textit{Status.}---The derivation of $\eta_T$ from
$G_2(X)$ requires computing the torsional mean free
path via lattice simulation and is left to future
work.  The qualitative result---that neutrino masses
are suppressed by both the aspect ratio
\emph{and} the weak trapping---is structural and
does not depend on the numerical value of $\eta_T$.

%-------------------------------------------------------------------
\subsection{Vortex topology and connection to
ISPG\_Quantum}
\label{sec:neutrino_vortex}
%-------------------------------------------------------------------

The helical wave \eqref{eq:circular_pol} is
simultaneously a cylindrical $\ell = \pm 1$
azimuthal-winding mode (wave description) and a
vortex line (field-configuration description).
This duality connects the present analysis to the
topological framework of ISPG\_Quantum, Sec.~10,
whose foundational assumptions (including the
identification of a $U(1)$-like phase on the
oscillon) are inherited here rather than
re-derived.

\textit{Identification.}---A helical mode
$\delta\phi_\pm \propto
\cos(kz - \omega t \pm \vartheta)$ has a spatial
argument that winds by $\pm 2\pi$ around the
propagation axis per wavelength.  This is
exactly a vortex of winding number
$n = \pm 1$ in the $(\vartheta, z)$ cylinder.
By the results of ISPG\_Quantum, Sec.~10:
\begin{itemize}
  \item winding $n = \pm 1$ $\Rightarrow$
    spin $J = \hbar/2$
    (Eq.~(10.4) of ISPG\_Quantum);
  \item candidate Berry phase $\gamma = \pm\pi$
    under exchange $\Rightarrow$
    Fermi--Dirac statistics in the structural route
    (Eq.~(10.12) of ISPG\_Quantum; full Fock-space
    proof open);
  \item opposite-winding structural mirror $\Rightarrow$
    antiparticle label
    (Sec.~10.2 of ISPG\_Quantum).
\end{itemize}

The neutrino is therefore assigned the same structural
fermion-label package used for electrons and quarks,
subject to the conditional spin/statistics and
antiparticle-status caveats in ISPG\_Quantum.  The
distinction is \emph{which sector} the vortex
inhabits:

\begin{center}
\renewcommand{\arraystretch}{1.3}
\begin{tabular}{@{}lll@{}}
\textbf{Particle} & \textbf{Vortex sector}
  & \textbf{Trapping strength} \\\hline
Electron, quarks
  & longitudinal / transverse
  & strong (Mathieu $q = 1.853$) \\
Neutrinos
  & polarisation rotation
  & weak ($q_T = q\,(R_\perp/\lambda_0)^2$)
\end{tabular}
\end{center}

All fermions are vortices; the neutrino's
uniqueness is that its vortex lives in the
weakly coupled polarisation-rotation channel.

\textit{Chirality from topology.}---Left-circular
polarisation ($n = +1$) and right-circular
($n = -1$) are topologically distinct: one cannot
be continuously deformed into the other without
unwinding the vortex.  This topological protection
is why chirality is an exact quantum number for
massless neutrinos.  For massive neutrinos, the
small mass term $\Omega_T$ in
Eq.~\eqref{eq:torsion_equation} allows mixing
between $n = +1$ and $n = -1$ at a rate
$\propto m_\nu/E$---exactly the helicity-flip
suppression observed experimentally.

%-------------------------------------------------------------------
\subsection{Beta decay in the torsional picture}
\label{sec:beta_decay}
%-------------------------------------------------------------------

The prototypical weak process---neutron beta
decay---receives a transparent geometric
interpretation:
\[
  n \;\to\; p \;+\; e^- \;+\; \bar{\nu}_e\,.
\]
Inside the neutron, one $d$-quark oscillon
carries a polarisation-rotation excitation
(a helical component winding around its core).
When this helical component radiates away, the
released torsional energy propagates outward as
an \textbf{antineutrino}---identified in this structural
picture as an opposite-winding antineutrino label
($n = -1$).  Simultaneously, the change
in the quark's transverse asymmetry
($-\tfrac{1}{3} \to +\tfrac{2}{3}$, i.e.\
$d \to u$) emits the difference as an
\textbf{electron}: a transverse-mode oscillon
carrying the released $\ell = 1$ asymmetry.

The $W$ boson is the \textbf{massive mediator of
the torsion--transverse coupling}: a nonlinear
excitation that temporarily converts
polarisation-rotation energy into transverse
energy and vice versa.  Its large mass
($m_W \approx 80$~GeV) reflects the suppressed
coupling between these sectors.  The $Z$ boson
is the neutral analogue: a
torsion--longitudinal coupling excitation.

%-------------------------------------------------------------------
\subsection{The W and Z bosons and the SU(2)
doublet}
\label{sec:WZ_bosons}
%-------------------------------------------------------------------

In the Standard Model, $W^\pm$ and $Z^0$ are
the massive gauge bosons of $\text{SU}(2)_L$.
In the polarisation-rotation picture:

\begin{itemize}
  \item The $W^\pm$ boson is identified
    heuristically with a localised excitation
    that \emph{simultaneously} carries helical
    azimuthal-winding and transverse
    character---the bridge mode that converts a
    helical excitation into a transverse
    displacement.  Its charge ($\pm 1$) reflects
    the transverse asymmetry it transfers.  Its
    large mass arises parametrically because
    creating such a bridge requires overcoming
    the suppressed coupling
    $\sim (R_\perp/\lambda_{\mathrm{osc}})^2$
    between the two sectors.

  \item The $Z^0$ boson is the neutral bridge:
    it couples the helical-winding sector to the
    longitudinal (scalar) sector without
    transferring transverse asymmetry.  Its mass
    is close to $m_W$ because the same coupling
    suppression governs both channels.
\end{itemize}

\textit{Parametric status of the $W,Z$ mass
scale.}---The present section fixes only the
\emph{parametric} structure of the $W,Z$ mass
scale: it is set by the product of the
carrier scale and an inverse aspect-ratio
factor built from
$(\lambda_{\mathrm{osc}}/R_\perp)$, but neither
the overall dimensional prefactor nor the
absolute numerical value follows at this level
of the derivation.  In particular, the naive
estimate
$m_W \sim h\nu_0/c^2 \cdot
(\lambda_{\mathrm{osc}}/R_\perp)$ does
\emph{not} reproduce the measured value
$m_W \approx 80.4$~GeV---the carrier frequency
$\nu_0$ enters at the Planck scale, so the
overall dimensional anchor must be fixed by a
separate constitutive input from the main
theory rather than read off from this estimate.
The absolute determination of $m_W, m_Z$ from
the $G_2(X)$ background is a quantitative task
deferred to a separate work built on the main
theory's action.

\textit{On the two-component space used here.}---%
The pair (left-helical $n=+1$, right-helical
$n=-1$) employed below is a \emph{chirality}
two-space of the helical-winding mode, not the
weak-isospin doublet $(\nu_L, e_L)$ of the
Standard Model: the two are related but
formally distinct ($\mathrm{SU}(2)_L$ pairs a
left-handed neutrino with a left-handed charged
lepton, not a particle with its antiparticle).
What the construction below gives directly is
therefore a chirality $\mathrm{U}(2)$ acting on
$(n=+1,\,n=-1)$; its identification with the
Standard-Model $\mathrm{SU}(2)_L$---including
the correct pairing of left-handed fermions
into weak doublets and the assignment of
hypercharges---requires an additional
constitutive step that is not derived here and
is deferred to a dedicated follow-up.

\textit{Chirality $U(2)$ from the 2-space.}---%
Taking the chirality two-space
$(n=+1,\,n=-1)$ as the working basis, the
$G_2(X)$ nonlinearity provides a Hermitian
coupling matrix between these two states.  The
most general $2\times 2$ Hermitian matrix is
parametrised by four real numbers,
\begin{equation}\label{eq:SU2_generators}
  H_{\mathrm{coupling}} = a_0\,\mathbb{I}
  + a_1\,\sigma_1 + a_2\,\sigma_2
  + a_3\,\sigma_3,
\end{equation}
with $\sigma_i$ the Pauli matrices on the
chirality basis.  Eq.~\eqref{eq:SU2_generators}
is therefore a chirality $\mathrm{U}(2)$, whose
$\mathrm{SU}(2)$ subgroup mixes the two
helicity-windings.  Whether this chirality
$\mathrm{SU}(2)$ can be promoted, by the
constitutive step referred to above, to the
Standard-Model $\mathrm{SU}(2)_L$ gauge group
is what we mean below by ``SU(2) sector''; it
is not proven in the present section.

\textit{Parity-violation sketch.}---Within the
chirality two-space, flipping the winding
number $n\to -n$ flips the relative coupling
phase between helical and transverse sectors,
so the two winding states enter the coupling
asymmetrically.  This provides a \emph{candidate
geometric source} of chiral asymmetry in the
helical sector; the precise map from this
asymmetry to the Standard-Model statement
``only left-handed fermions couple to $W^\pm$''
requires the constitutive SU(2)-identification
step above and is not completed here.

\textit{From global to local SU(2).}---The
coupling matrix \eqref{eq:SU2_generators}
depends on the oscillon profile $\Phi_0(r)$,
which varies in space: the coefficients
$a_i(\mathbf{x})$ inherit the spatial
dependence of $\Phi_0$.  The global chirality
SU(2) is thereby promoted to a \emph{local}
symmetry, with the oscillon-background
gradients playing the role of the gauge
connection---the same mechanism by which
emergent gauge fields arise from spatially
varying order parameters in condensed-matter
systems (Berry connections in Bloch bands,
spin-orbit coupling in topological
insulators).  The formal derivation of the
Yang--Mills kinetic term from the $G_2(X)$
coupling structure is a future-work task.

\textit{Weinberg-angle analogue.}---Within the
chirality-SU(2) framework of this section, the
analogue of the weak mixing angle is the ratio
of coupling strengths between the
helical--transverse channel and the
helical--longitudinal channel:
\begin{equation}\label{eq:weinberg_angle}
  \tan\theta_W^{\mathrm{(chir)}}
  = \frac{g_1}{g_2}
  = \frac{\text{longitudinal coupling}}
  {\text{transverse coupling}}.
\end{equation}
Both couplings arise from the same $X^2$ term
but project onto different multipole sectors
($\ell = 0$ vs.\ $\ell = 1$); their ratio is
therefore a calculable function of the
oscillon profile $\Phi_0(r)$ and the Mathieu
parameter $q$.  Identifying
$\theta_W^{\mathrm{(chir)}}$ with the measured
Standard-Model Weinberg angle ($\sin^2\theta_W
= 0.231$) additionally requires the
constitutive $\mathrm{SU}(2)_L$-identification
step noted above; until that step is carried
out, Eq.~\eqref{eq:weinberg_angle} is an
\emph{internal} ratio of this paper, not an
experimental prediction.

%-------------------------------------------------------------------
\subsection{Quantitative programme and future
work}
\label{sec:neutrino_future}
%-------------------------------------------------------------------

\textit{Established in this section:}
\begin{itemize}
  \item Polarisation-rotation mode emerges from
    $G_2(X)$ via the helical perturbation ansatz
    (Eq.~\ref{eq:helical_ansatz}); no new field
    required (\S\ref{sec:torsion_derivation}).
  \item The torsional Mathieu equation
    (Eq.~\ref{eq:torsion_equation}) is derived
    from $\mathcal{L}^{(2)}_{\delta\phi}$
    by decomposing into circular-polarisation
    channels and projecting onto the
    polarisation-angle sector.
  \item The suppression factor
    $R_\perp/\lambda_{\mathrm{osc}}
    = 1/(2\pi)$ is determined by the Compton
    radius of the oscillon, not by neutrino
    masses (Eq.~\ref{eq:Rperp_corrected}).
  \item The additional suppression from weak
    torsional trapping ($\eta_T$,
    Eq.~\ref{eq:trapping_fraction}) accounts
    qualitatively for the observed
    $m_\nu/m_e \sim 10^{-7}$; the first-principles
    derivation of $\eta_T$ is left to future work.
  \item Spin-$\frac{1}{2}$, Fermi--Dirac
    statistics, and chirality follow from the
    vortex topology
    (\S\ref{sec:neutrino_vortex};
    ISPG\_Quantum, Sec.~10).
  \item The SU(2) doublet structure arises from
    the two-component circular-polarisation space
    coupled by $G_2(X)$
    (\S\ref{sec:WZ_bosons}).
  \item The SU(2) global symmetry is promoted to
    local via the spatial dependence of $\Phi_0(r)$
    (\S\ref{sec:WZ_bosons}); formal Yang--Mills
    derivation: future work.
\end{itemize}

\textit{Quantitative predictions requiring
lattice or numerical computation:}
\begin{itemize}
  \item \textbf{Neutrino masses.}\;
    The parametric structure
    $m_\nu \sim m_e\,(R_\perp/\lambda_{\mathrm{osc}})^2\,
    \eta_T$ is fixed by the present derivation: the
    aspect-ratio factor
    $R_\perp/\lambda_{\mathrm{osc}} = 1/(2\pi)$
    follows from the Compton radius of the oscillon
    (Eq.~\ref{eq:Rperp_corrected}), and the trapping
    fraction $\eta_T$
    (Eq.~\ref{eq:trapping_fraction}) is defined
    through the nonlinear scattering cross-section
    of torsional modes on the $\nu_0$-background.
    The evaluation of $\eta_T$ from $G_2(X)$
    requires a dedicated lattice computation and is
    the topic of a separate work; once $\eta_T$ is
    obtained, the three neutrino masses follow from
    the parametric form with no further free
    parameters.

  \item \textbf{PMNS matrix.}\;
    The overlap integrals between the three
    torsional eigenstates and the charged-lepton
    oscillon coupling modes determine the mixing
    angles $\theta_{12}$, $\theta_{23}$,
    $\theta_{13}$ and the CP phase $\delta$.

  \item \textbf{Mass ordering.}\;
    The torsional Mathieu spectrum predicts
    whether the hierarchy is normal or inverted,
    depending on the sign of $q_T$ corrections
    from the oscillon boundary conditions.

  \item \textbf{Majorana vs.\ Dirac.}\;
    A circularly polarised vortex is symmetric
    under combined charge conjugation and
    parity ($CP$) if and only if the boundary
    conditions identify $n = +1$ with $n = -1$
    at the oscillon core.  This identification
    produces a Majorana mass; its absence
    preserves a Dirac mass.  The distinction is
    determined by the topology of the oscillon
    core---accessible via lattice
    simulation.

  \item \textbf{Weinberg angle.}\;
    $\theta_W$ follows from the ratio of
    torsion--longitudinal to
    torsion--transverse couplings
    (Eq.~\ref{eq:weinberg_angle}), calculable
    from the oscillon profile.  The prediction
    $\sin^2\theta_W \stackrel{?}{=} 0.231$ is
    a decisive quantitative test.

  \item \textbf{Yang--Mills from $G_2(X)$.}\;
    Deriving the $F_{\mu\nu}^a F^{a\mu\nu}$
    kinetic term from the position-dependent
    coupling matrix
    $H_{\mathrm{coupling}}(\mathbf{x})$.
\end{itemize}

This mechanism, once numerically confirmed,
completes the ISPG vibration programme: four
perturbation modes of a single scalar field
account for all four fundamental forces---gravity,
electromagnetism, the strong force, and the weak
force---from a single field, a single vibration
frequency $\nu_0$, and a single nonlinear
interaction $G_2(X)$.

%===================================================================
\section{Massless Sector: Transverse-Traceless Modes
and Gluon Confinement}
\label{sec:massless}
%===================================================================

The preceding sections addressed all
\emph{massive} particles.  The bi-conformal medium
also supports \emph{massless} transverse modes.
This section outlines the structural classification
that emerges naturally from the oscillon picture.
The massless exterior sector corresponds to the
gravitational-wave (spin-2) part of the theory;
the identification of the electromagnetic photon
inside this framework is a separate geometric
question whose full development is the topic of a
companion paper.  The confined counterpart of the
free transverse mode is identified with the gluon,
and the confinement mechanism is derived below.

%-------------------------------------------------------------------
\subsection{Four modes of a single medium}
\label{sec:four_modes}
%-------------------------------------------------------------------

The bi-conformal metric
$ds^{2} = -e^{\phi}\,c^{2}dt^{2}
+ e^{-\phi}\,d\sigma^{2}$ admits two types of
perturbation: scalar ($\delta\phi$) and
transverse-traceless tensor ($\delta g_{ij}^{TT}$,
two polarizations).  An oscillon, as a localized
nonlinear structure, divides the medium into two
geometric regions: the \textbf{interior}
(the core where $G_2(X)$ nonlinearity is strong)
and the \textbf{exterior} (the asymptotic region
where the field is approximately linear).

\paragraph{Relation to the cylindrical
$\ell$-decomposition of \S\ref{sec:neutrinos}.}
The ``four modes'' enumerated below---(scalar,
tensor)$\,\times\,$(interior, exterior)---are a
\emph{different} decomposition of the full
perturbation content from the cylindrical-%
multipole classification used for the scalar
sector in \S\ref{sec:four_modes} (longitudinal
$\ell=0$, transverse $\ell=1$, membrane
$\ell\geq 2$, helical $\ell=\pm 1$ winding).
They are complementary, not redundant: the
$\ell$-decomposition there is internal to the
scalar sector and classifies the cylindrical
angular structure of $\delta\phi$, whereas the
scalar-vs.-tensor split used here sits one level
higher, at the level of the metric perturbations
of the bi-conformal background.  In particular,
the tensor sector (rows 3--4 of the table below)
is not part of the scalar-$\ell$ tower of
\S\ref{sec:four_modes}; conversely, the scalar
helical $\ell=\pm 1$ mode of
\S\ref{sec:four_modes} sits inside row~1
(scalar, interior) of the table below.  No
unified single ``four-mode'' enumeration is
asserted to cover both decompositions
simultaneously.

Combining two perturbation types with two
geometric regions yields exactly four vibration
modes:

\medskip
\begin{center}
\renewcommand{\arraystretch}{1.4}
\begin{tabular}{@{}r@{\;:\;\;}l@{\;\;$\to$\;\;}l@{\;\;$\to$\;\;}l@{}}
\toprule
\textbf{Mode} & \textbf{Type / Region}
  & \textbf{Particle} & \textbf{Properties} \\
\midrule
1 & scalar, localized (core)
  & oscillon (matter)
  & $m > 0$, integer harmonics \\
2 & torsional (twist)
  & neutrino
  & $m \approx 0$, chiral, three flavours \\
3 & tensor, free (exterior)
  & gravitational wave (spin-2)\footnotemark
  & $m = 0$, propagating \\
4 & tensor, confined ($\nu_N \neq \nu_0$)
  & gluon
  & $m = 0$, confined by frequency gap \\
\bottomrule
\end{tabular}
\end{center}

\footnotetext{The electromagnetic photon is a
$U(1)$ gauge boson ($A_\mu$) treated in the main
theory~\cite{Lotishvili2026} in the standard
minimally coupled form; its embedding into the
present oscillon/frequency-to-mass picture is
deferred to a companion paper.}

\medskip
No additional fields or symmetry postulates
are introduced.  The four categories exhaust the
perturbation content of the bi-conformal metric
around an oscillon background: scalar vs.\
tensor ($\times\,2$), interior vs.\ exterior
($\times\,2$).

%-------------------------------------------------------------------
\subsection{Transverse-traceless modes and the
gravitational-wave sector}
\label{sec:photon}
%-------------------------------------------------------------------

A transverse-traceless perturbation of the metric
propagating in the exterior region satisfies the
linearized wave equation with zero effective mass.
It travels at $c$ and carries two independent
polarization states.  Within the present geometric
framework these modes correspond to
gravitational waves---the massless spin-2 sector
already established in the main
theory~\cite{Lotishvili2026}; they do not give
rise to a photon, whose identification is a
separate geometric question treated below.

\textit{Propagation.}---The wave travels at
$c$, the transverse wave speed of the
$\nu_0$-string.  Because the string is uniform
(one $\nu_0$, one tension), the speed is the
same everywhere and for all frequencies---no
dispersion.  The two polarizations correspond to
the two independent directions in the transverse
plane.

\textit{Speed of light as a fundamental
limit.}---The speed $c$ is the propagation speed
of perturbations on the medium that \emph{is}
space.  Nothing can travel through a medium
faster than the medium's own wave speed;
$c$ is therefore a geometric maximum, not
an empirical accident.

\textit{Photon identification.}---The
identification of the electromagnetic photon
within the ISPG medium is not part of the present
paper.  The electromagnetic sector is carried by
the standard $U(1)$ gauge field $A_\mu$ in the
usual minimally coupled form (see
\cite{Lotishvili2026}, Appendix~5); its embedding
into the oscillon frequency-to-mass framework is a
separate question whose full development is the
topic of a companion paper.  In the remainder of
this paper, ``massless transverse mode'' refers to
the gravitational-wave sector; any statement about
photon propagation in a gravitational field
(lensing, redshift, Shapiro delay) rests on the
standard $A_\mu$ treatment of the main theory, not
on a direct identification of the photon with the
TT metric mode.

%-------------------------------------------------------------------
\subsection{Gluon: confined transverse wave}
\label{sec:gluon}
%-------------------------------------------------------------------

Inside the oscillon core the medium vibrates at
the oscillon's own frequency
$\nu_N = m_N c^2/h \neq \nu_0$.
Transverse-traceless perturbations born in
this region inherit $\nu_N$ as their
characteristic frequency.  When such a
perturbation reaches the core boundary and
attempts to propagate into the exterior, it
encounters a medium whose background vibration
is $\nu_0$---a fundamentally different rhythm.

\textit{Confinement by frequency incompatibility.}---
The $\nu_0$-background acts as a dispersive
medium for transverse modes tuned to
$\nu_N \neq \nu_0$.  Because the $G_2(X)$
nonlinearity couples all oscillation modes to
the background, a tensor wave whose frequency is
not commensurate with $\nu_0$ falls inside a
\textbf{band gap}: it cannot propagate as a
travelling wave and instead decays
exponentially outside the core
(evanescent wave).  The gluon is therefore
confined not by a material wall but by
\emph{frequency mismatch}---the exterior
medium's $\nu_0$-rhythm destructively interferes
with any wave tuned to the interior $\nu_N$.

\textit{Why scalar modes are $\nu_0$-harmonic
but tensor modes are not.}---This is the crucial
asymmetry.  Approximate the oscillon core as a
spherical cavity of radius $R$.  Both scalar and
tensor perturbations satisfy radial wave equations
inside the cavity, but with different angular
momentum constraints:

\begin{itemize}
  \item \textbf{Scalar perturbations} ($\ell = 0$
    allowed).  The radial eigenvalues are the zeros
    of the spherical Bessel function $j_0(kR) = 0$,
    located at
    \begin{equation}
      k_n R = n\pi, \qquad n = 1,2,3,\ldots
    \end{equation}
    The resulting frequencies
    $\omega_n = n\pi c/R$ form an
    \textbf{integer-harmonic series}.
    This is why the Mathieu mass spectrum produces
    $\nu_0$-harmonics: the scalar cavity is
    one-dimensional in its radial structure,
    analogous to a vibrating string.

  \item \textbf{Tensor perturbations}
    ($\ell \geq 2$ required, since
    transverse-traceless modes carry spin-2).  The
    radial eigenvalues are the zeros of
    $j_{\ell}(kR) = 0$ for $\ell \geq 2$.
    For $\ell = 2$, the first three zeros are:
    \begin{equation}
      k_n R \;\approx\; 5.763,\; 9.095,\; 12.323,
      \;\ldots
      \label{eq:bessel_zeros}
    \end{equation}
    These are \textbf{not integer multiples of
    $\pi$}: the ratio of the lowest tensor
    eigenvalue to the scalar fundamental is
    \begin{equation}
      \frac{k_1^{(\ell=2)}}{k_1^{(\ell=0)}}
      = \frac{5.763}{\pi}
      \approx 1.835,
      \label{eq:inharmonic_ratio}
    \end{equation}
    an irrational number.  The tensor modes form
    an \textbf{inharmonic series}---analogous to
    the overtones of a drum, which are not integer
    multiples of the fundamental (unlike a string).
\end{itemize}

This is the mathematical origin of confinement:
scalar modes (matter) and free transverse-traceless
modes (gravitational waves) live in an
integer-harmonic tower built on $\nu_0$ and
propagate freely through the $\nu_0$-background.
Tensor cavity modes (gluons) have frequencies
that are \emph{incommensurate} with $\nu_0$:
\begin{equation}
  \boxed{
  \frac{\omega_{\mathrm{gluon}}}{\nu_0}
  = \frac{j_{\ell,\,n}}{\pi}
  \;\notin\; \mathbb{Z}
  \qquad (\ell \geq 2).
  }
  \label{eq:confinement_criterion}
\end{equation}
A wave whose frequency is not a $\nu_0$-harmonic
cannot sustain coherent propagation in the
$\nu_0$-background and decays evanescently.
A useful acoustic analogy: the $\nu_0$-background
is a \emph{guitar string}---a one-dimensional
resonator whose overtones are exact integer
multiples of the fundamental.  The oscillon
interior is a \emph{handpan} (a curved metal
membrane)---a two-dimensional resonator whose
overtones are governed by Bessel zeros and are
\emph{not} integer multiples of the fundamental.
A vibration born on the handpan cannot sustain
itself on the guitar string because the two
harmonic series are incommensurate.  This
``string--membrane asymmetry''
($j_{0,n}/\pi \in \mathbb{Z}$ vs.\
$j_{\ell \geq 2,\,n}/\pi \notin \mathbb{Z}$)
is a theorem of Bessel function theory, not an
assumption---confinement follows from the
\textbf{geometry of the cavity}, specifically
from the fact that spin-2 modes require
$\ell \geq 2$.

This immediately explains two key features of
the strong interaction:

\begin{enumerate}
  \item \textbf{Confinement.}\;
    A single gluon cannot exist in free space
    because the $\nu_0$-background extinguishes
    any transverse mode not in resonance with it.
    Photons \emph{can} propagate freely because,
    as metric perturbations of the exterior
    medium, they are already at the medium's
    native frequency.

  \item \textbf{Flux tube and linear potential.}\;
    When two quarks (sub-structures within the
    oscillon) are separated, the region between
    them must sustain the interior frequency
    $\nu_N$ against the exterior $\nu_0$ pressure.
    Maintaining $\nu_N$ in a strip of width
    $\sim R_{\mathrm{core}}$ and length $L$
    costs energy proportional to $L$:
    \begin{equation}\label{eq:flux_tube}
      V(L) \;\approx\; \sigma\,L, \qquad
      \sigma \sim
      \frac{G_2''(X_0)\,\nu_N^2}{c^2}\,
      R_{\mathrm{core}}^2,
    \end{equation}
    reproducing the linear confining potential
    observed in lattice QCD.  When $V(L)$
    exceeds $2m_q c^2$, the tube breaks and a new
    quark--antiquark pair materializes---string
    breaking, as observed experimentally.
\end{enumerate}

\textit{Physical picture.}---The oscillon interior
is a ``bubble'' of frequency $\nu_N$ immersed
in a $\nu_0$-ocean.  Gluons are transverse
vibrations that can live inside the bubble but
drown in the ocean.  Pulling two quarks apart
stretches the bubble into a tube; the ocean
presses on the tube walls from all sides, and
the energy cost grows linearly with length.

\textit{Numerical verification.}---The
inharmonicity criterion was verified numerically
(see \texttt{python/gluon\_modes.py}).
The glueball spectrum predicted by the Bessel-zero
ratios $j_{\ell,n}/j_{2,1}$ can be compared
directly with lattice QCD glueball masses
(Morningstar \& Peardon 1999; Chen et al.\ 2006):
\begin{equation}\label{eq:glueball_comparison}
\begin{array}{lcccc}
  \hline
  \text{State} & M_{\text{lat}}\;\text{(MeV)}
    & \text{ratio}_{\text{lat}}
    & \text{ratio}_{\text{ISPG}}
    & \Delta \\
  \hline
  0^{++} & 1710 & 1.000 & 1.000\;(\ell\!=\!2,n\!=\!1) & 0\% \\
  2^{++} & 2390 & 1.398 & 1.420\;(\ell\!=\!4,n\!=\!1) & 1.6\% \\
  2^{-+} & 3100 & 1.813 & 1.807\;(\ell\!=\!3,n\!=\!2) & 0.3\% \\
  3^{++} & 3670 & 2.146 & 2.138\;(\ell\!=\!2,n\!=\!3) & 0.4\% \\
  \hline
\end{array}
\end{equation}
Four of the six lowest glueball states match to
$\leq 2\%$; the remaining two
($0^{-+}$ and $1^{+-}$) deviate by $\sim 5\%$,
likely because the free-cavity Bessel zeros neglect
the $G_2(X)$ self-interaction corrections.
The fundamental ratio
\begin{equation}\label{eq:fundamental_ratio}
  \frac{M(0^{++})}{M(p)}
  = \frac{1710}{938}
  = 1.823
  \;\approx\;
  \frac{j_{2,1}}{\pi}
  = 1.835
  \qquad (\Delta = 0.66\%),
\end{equation}
provides a striking numerical coincidence.

\textit{Scope of ``parameter-free''.}---%
The internal ratios of the glueball spectrum
are fixed entirely by the Bessel-zero ratios
$j_{\ell,n}/j_{2,1}$ once the
$(\ell,n)\leftrightarrow J^{PC}$ assignments of
Eq.~\eqref{eq:glueball_comparison} are accepted;
in that sense the four-state fit is free of
tunable coefficients.  Two caveats must be
stated honestly, however.  (i)~The
$(\ell,n)\leftrightarrow J^{PC}$ assignment is
itself an educated identification of
spherical-cavity modes with lattice-QCD
glueball states; alternative assignments would
produce different ratios and have not been
ruled out here.  (ii)~The ratio-to-proton in
Eq.~\eqref{eq:fundamental_ratio} uses
$M(p) = 938$~MeV as an external
dimensionful input---the proton mass is not
derived in this section; it is only being used
to set the absolute scale.  The genuinely
parameter-free content is therefore the set of
internal glueball ratios; the proton-anchor
agreement in
Eq.~\eqref{eq:fundamental_ratio} is a
\emph{consistency check} against that anchor,
not an independent prediction of $M(p)$ itself.

\textit{The eightfold structure.}---
The Standard Model has exactly eight gluons,
corresponding to the eight generators of SU(3).
In the spherical-cavity approximation, the lowest
TT tensor energy level ($\ell = 2$, $n = 1$) has
two polarisation families (even and odd parity),
each with $2\ell+1 = 5$ orientations, giving
$2 \times 5 = 10$ modes.  A baryon oscillon,
however, contains three sub-oscillons (quarks)
arranged in a triangle, reducing the symmetry from
$\mathrm{SO}(3)$ to the discrete group $D_3$.
Under this restriction the $\ell = 2$
representation decomposes as
\begin{equation}\label{eq:D3_decomp}
  D^{2}\big|_{D_3}
  \;=\; A_1 \;\oplus\; E \;\oplus\; E
  \;=\; 1 + 2 + 2.
\end{equation}
The $C_3$ selection rule
($e^{im_p \cdot 2\pi/3}$ sums to zero unless
$m_p = 0, \pm 3$) restricts the perturbation to
the $Y_{4,0}$ (oblate) and $Y_{4,\pm 3}$
(triangular) spherical harmonics.  Computing the
Gaunt integrals exactly yields the $D_3$
perturbation matrix
(see \texttt{python/gluon\_proof.py}):
\begin{equation}\label{eq:D3_matrix}
  H_{D_3} = \varepsilon\!
  \begin{pmatrix}
    g_0^{(-2)} & 0 & 0 & g_3 & 0 \\
    0 & g_0^{(-1)} & 0 & 0 & g_3 \\
    0 & 0 & g_0^{(0)} & 0 & 0 \\
    g_3 & 0 & 0 & g_0^{(+1)} & 0 \\
    0 & g_3 & 0 & 0 & g_0^{(+2)}
  \end{pmatrix}\!,
\end{equation}
where $g_0^{(m)}$ are the Gaunt coefficients for
$Y_{4,0}$ and $g_3 = -0.238$ is the $Y_{4,\pm 3}$
coupling.  The matrix has block-diagonal structure:
$m = 0$ is decoupled ($A_1$ irrep) while
$\{m\!=\!-2,\,m\!=\!+1\}$ and
$\{m\!=\!-1,\,m\!=\!+2\}$
form two $2\times 2$ blocks whose eigenvalues are
pairwise degenerate ($E$ irreps).  The resulting
eigenvalues confirm the $D_3$ decomposition:
\begin{equation}\label{eq:mode_split}
\begin{array}{lccl}
  \hline
  \text{Subgroup} & \text{deg.}
    & \lambda & \text{status} \\
  \hline
  E_{\downarrow} & 2 & -0.145 & \text{confined} \\
  A_1            & 1 & +0.242 & \text{confined} \\
  E_{\uparrow}   & 2 & +0.347 & \text{resonant} \\
  \hline
\end{array}
\end{equation}
When the nonperturbative parameter $\varepsilon$
is large enough that
$E_T + \varepsilon\,\lambda(E_{\uparrow})
\approx E_S^{(n=2)} = (2\pi)^2$,
the upper $E$ pair enters resonance with the
nearest scalar eigenvalue.
The $\ell=2$ tensor fundamental at
$(kR)^2 = j_{2,1}^2 = 33.22$ lies below
the scalar $n=2$ level at $(2\pi)^2 = 39.48$;
the gap is $\Delta = 6.26$.  The Gaunt eigenvalue
spread is $0.492$, so resonance occurs at
$\varepsilon \approx \Delta / 0.492 \approx 13$.
A perturbative estimate gives
$\varepsilon_{\text{pert}} \approx 0.5$; the full
nonperturbative value was obtained by direct 2D
Helmholtz simulation comparing a circular cavity
with a $D_3$-deformed (rounded-triangle) cavity
of the same area
(see \texttt{python/gluon\_epsilon.py}).
At $25\%$ triangular deformation (matching the
proton's quark geometry $d = R/\sqrt{3}$), the 2D
simulation yields
$\varepsilon_{\text{2D}} \approx 6.8$.
In three dimensions the radial overlap of the
$\ell = 2$ Bessel profile contributes an additional
factor $j_{2,1}/\pi = 1.835$, giving
\begin{equation}\label{eq:eps_3D}
  \varepsilon_{\text{3D}}
  \;\approx\;
  \varepsilon_{\text{2D}}
  \times \frac{j_{2,1}}{\pi}
  \;\approx\;
  6.8 \times 1.835
  \;\approx\; 12.5,
\end{equation}
consistent with the required
$\varepsilon \approx 13$ to within $4\%$.

\textit{Non-independence note.}---%
The ratio $j_{2,1}/\pi = 1.835$ entering
Eq.~\eqref{eq:eps_3D} is the \emph{same}
Bessel-zero ratio that fixes the
glueball-to-proton ratio in
Eq.~\eqref{eq:fundamental_ratio}.  The two
agreements therefore share a single structural
input---the first Bessel zero of $J_2$---rather
than constituting two logically independent
predictions; the combined claim is that one
geometric quantity, $j_{2,1}$, controls both
the $0^{++}$-to-proton ratio and the
2D$\to$3D dressing of the triangular-cavity
eigenvalue problem, not that two separate
quantitative hits have been scored.

At resonance the upper $E$ pair acquires
$\sim 50\%$ scalar admixture and radiates
through the $\nu_0$-exterior in $\sim 2$ oscillation
periods.  The $A_1$ mode ($\sim 10\%$ leak) and
$E_{\downarrow}$ pair ($< 1\%$ leak) remain
confined.  The total count is
\begin{equation}\label{eq:eight_gluons}
  \boxed{
  \underbrace{5}_{\text{odd parity}}
  \;+\;
  \underbrace{2}_{E_{\downarrow}}
  \;+\;
  \underbrace{1}_{A_1}
  \;=\; 8
  \;=\; \dim\!\left[\mathrm{SU}(3)_{\mathrm{adj}}\right]
  \;=\; 3^2 - 1.
  }
\end{equation}
The count of eight confined modes in the
baryon-cavity problem matches the eight
generators of $\mathrm{SU}(3)_{\mathrm{adj}}$.

\textit{Scope of the ``8'' count.}---%
What the calculation above rigorously
establishes is the mode count ``$8$'' in a
\emph{specific} cavity problem: the TT
perturbations of a $D_3$-symmetric
(triangular) baryon-like oscillon, after
the Gaunt-coefficient restriction of the
$\ell=2$ multiplet.  That count \emph{coincides}
with $\dim\mathrm{SU}(3)_{\mathrm{adj}} =
3^2-1 = 8$.  It does \emph{not}, however,
derive the Standard-Model gauge group
$\mathrm{SU}(3)_c$, which must be
\emph{universal} (the same eight gauge bosons
act on $q\bar q$ mesons as on $qqq$ baryons),
whereas the $D_3$-triangle argument is
specific to three-quark hadronic geometry: a
$q\bar q$ meson would sit in the $D_2$ dihedral
symmetry class and would produce a different
mode count from this construction.  A genuine
derivation of $\mathrm{SU}(3)_c$ as a gauge
group must therefore come from a
hadron-independent structure of the $G_2(X)$
sector, not from a specific baryonic cavity
geometry; the numerical coincidence
``$3^2-1=8$'' recorded here is a
\emph{consistency check} of the
baryon-sector mode count against the known SM
adjoint dimension, and the gauge-group
derivation proper is a task deferred to
subsequent work.

The next computational step is to solve the full
coupled perturbation equation
\begin{equation}\label{eq:gluon_modes}
  \Box\,h_{ij}^{TT}
  + \mathcal{M}^2[\nu_N, \nu_0, G_2]\,
  h_{ij}^{TT} = 0
\end{equation}
inside a three-dimensional three-body oscillon
with evanescent exterior boundary conditions.
This 3D finite-element calculation will sharpen
the estimate~\eqref{eq:eps_3D} and determine
whether the eight surviving modes' coupling
algebra reproduces SU(3).  If so, the colour
gauge group emerges geometrically from the
oscillon's $D_3$ cavity structure.

%-------------------------------------------------------------------
\subsection{Summary and future programme}
\label{sec:massless_summary}
%-------------------------------------------------------------------

The single-field, single-medium ontology of ISPG
accommodates all particle categories through the
vibration-mode classification:

\begin{itemize}
  \item \textbf{Massive particles}
    (\S\ref{sec:mathieu_numerical}):
    scalar, localized $\to$ oscillon harmonics,
    $m_N = m_1\,b_N(q)$.
  \item \textbf{Neutrinos}
    (\S\ref{sec:neutrinos}):
    polarisation-rotation wave (helical vortex
    with $n = \pm 1$)
    $\to$ inherent chirality, $m_\nu \ll m_e$;
    the weak interaction is the nonlinear coupling
    between torsional and other modes.
  \item \textbf{Gravitational waves}
    (\S\ref{sec:photon}):
    tensor, free $\to$ propagating
    transverse-traceless mode (spin-2), $m = 0$.
    The electromagnetic photon ($A_\mu$,
    $U(1)$ gauge boson) is treated in the main
    theory~\cite{Lotishvili2026}; its embedding
    into the present oscillon framework is the
    topic of a companion paper.
  \item \textbf{Gluons}
    (\S\ref{sec:gluon}):
    tensor, confined by $\nu_N / \nu_0$ frequency
    mismatch $\to$ interior transverse mode,
    $m = 0$, confined; flux tube $=$ frequency
    bridge, $V(L) \propto L$.
\end{itemize}

This classification is \textbf{exhaustive}:
the $\nu_0$-string with transverse extent
supports exactly four vibration types
(longitudinal, transverse, torsional, membrane),
and the oscillon divides space into two regions
(interior, exterior).  No fifth vibration mode
exists without introducing new fields or extra
dimensions---four forces from four modes.

The quantitative development of Sections
\ref{sec:photon}--\ref{sec:gluon} (derivation
of coupling constants, SU(3) emergence,
electroweak unification) constitutes a major
programme for future work.

%===================================================================
\section{Why Three Dimensions?  Emergence of
Spatial Dimensionality from a One-Dimensional
Substrate}
\label{sec:dimensions}
%===================================================================

The preceding sections established that the
$\nu_0$-background has integer harmonics
$\omega_n = n\pi c/R$ (the zeros of
$j_0(kR)$), while tensor cavity modes have
non-integer (Bessel) harmonics.
This spectral fact has a deeper consequence
that we now develop: \emph{the integer-harmonic
structure of the scalar sector is the signature
of a fundamentally one-dimensional substrate},
and the three spatial dimensions we observe
emerge as a geometric necessity of vibration.

%-------------------------------------------------------------------
\subsection{The scalar sector is
one-dimensional}
\label{sec:1D_substrate}
%-------------------------------------------------------------------

The spherical Bessel function of order zero is
\begin{equation}\label{eq:j0_sin}
  j_0(x) = \frac{\sin x}{x}\,.
\end{equation}
Its zeros lie at $x = n\pi$, $n = 1, 2, 3,
\ldots$---\emph{identical} to the eigenvalues
of a one-dimensional vibrating string of length
$L$ clamped at both ends:
\begin{equation}\label{eq:string_eigenvalues}
  k_n L = n\pi \,,\qquad
  \omega_n = n\,\omega_1\,.
\end{equation}
This is not an analogy but a mathematical
identity.  The radial equation for $\ell = 0$
modes in a spherical cavity, after the standard
substitution $u(r) = r\,\psi(r)$, \emph{is}
the one-dimensional wave equation:
\begin{equation}\label{eq:radial_1D}
  \frac{d^2 u}{dr^2} + k^2 u = 0\,,\qquad
  u(0) = u(R) = 0\,.
\end{equation}
The isotropic ($\ell = 0$) character of the
scalar field $\phi$ eliminates all angular
dependence; three spatial coordinates collapse
into a single radial variable, and the mode
structure becomes \emph{exactly} that of a
string clamped at two endpoints.

\textit{Scope of the claim.}---The statement is
a theorem about the $\ell = 0$ mode spectrum,
not an assertion that the physical universe
``is'' a literal one-dimensional string.  The
$u = r\psi$ substitution and the resulting
one-dimensional wave equation hold for \emph{any}
three-dimensional spherically symmetric
Helmholtz problem; this is a standard
mathematical identity, not an ISPG-specific
result.  What the theorem therefore establishes,
for the \emph{scalar sector} of the medium---the
sector responsible for gravity and particle
masses---is only that its $\ell=0$ spectrum is
controlled by a one-dimensional
Sturm--Liouville structure.  The higher-$\ell$
sectors (transverse, torsional, membrane)
supply the additional degrees of freedom for
electromagnetism, weak and strong interactions;
the full field $\phi(\mathbf{x},t)$ lives on a
three-dimensional spatial domain throughout.
The ``substrate is one-dimensional'' reading
that appears in the remainder of this section
is a proposed \emph{ontological interpretation}
built on top of the spectral theorem, not an
additional mathematical consequence of it; we
flag it explicitly as an interpretive posture,
not a derived result, and its experimental
implications are separated out in
\S\ref{sec:entanglement_geometry}.

Because there is only one field $\phi$ and its
isotropic mode is one-dimensional, the universal
frequency $\nu_0$ is automatically the same
everywhere:\ one string has one pitch.  No
synchronisation mechanism is needed.  This
resolves the question of why $\nu_0$ is a
universal constant---it is the vibration
frequency of a single quantum entity, not of
many independent oscillators.

%-------------------------------------------------------------------
\subsection{Sequential emergence of
dimensions}
\label{sec:dim_emergence}
%-------------------------------------------------------------------

A clamped string (1D) can vibrate only if it is
free to displace in at least one transverse
direction.  This displacement opens a second
dimension:

\begin{enumerate}
\item \textbf{First dimension:}
  the axis of tension---the direction along
  which the ``string'' is stretched
  (radial coordinate of the $\ell = 0$ mode).
  This is the longitudinal (scalar) direction.

\item \textbf{Second dimension:}
  the string vibrates transversely---it moves
  ``up'' and ``down'' perpendicular to the
  tension axis.  A plane (2D) is created.
  In the spherical-harmonic language, this
  introduces the polar angle $\theta$ and with
  it the $\ell = 1$ (dipole) modes.

\item \textbf{Third dimension:}
  on the 2D membrane so created, a second
  transverse vibration mode can exist---one
  that is perpendicular to both the tension
  axis and the first transverse direction.
  In spherical harmonics this is the azimuthal
  angle $\varphi$ and corresponds to
  $\ell = 2$ (quadrupole, tensor) modes.
  These are exactly the \textbf{gluons}
  of Section~\ref{sec:gluon}.
\end{enumerate}

The chain is
\begin{equation}\label{eq:dim_chain}
  \boxed{
  \underbrace{1\text{D}}_{\text{string}\;(\ell=0)}
  \;\xrightarrow{\;\text{transverse vibration}\;}
  \underbrace{2\text{D}}_{\text{membrane}\;(\ell \leq 1)}
  \;\xrightarrow{\;\text{tensor vibration}\;}
  \underbrace{3\text{D}}_{\text{cavity}\;(\ell \leq 2)}
  }
\end{equation}

No further dimension is generated because no
new vibration type is needed: with three spatial
dimensions the angular spectrum is complete
($\theta, \varphi$ exhaust the unit sphere).
A hypothetical fourth spatial dimension would
require a new angular coordinate---but $S^2$
has no room for one.

\textit{Why exactly $3$ and not $2$ or $4$?}---
A purely two-dimensional universe ($\ell \leq 1$
only) cannot support \emph{tensor} modes:
transverse-traceless perturbations require
$\ell \geq 2$, which demands the two angular
coordinates $(\theta, \varphi)$ of a
three-dimensional space.  In $(2+1)$-dimensional
general relativity, gravitational waves carry
no propagating degrees of freedom---a well-known
result.  Three spatial dimensions is therefore
the \textbf{minimum} in which both scalar and
tensor modes coexist:

\begin{equation}\label{eq:dim_argument}
  d = 3
  \;=\;
  \underbrace{1}_{\text{longitudinal}
    \;(\ell=0,\;\text{scalar})}
  \;+\;
  \underbrace{2}_{\text{transverse}
    \;(\ell=2,\;\text{tensor required})}.
\end{equation}

The number of spatial dimensions is constrained
from below by the requirement that the substrate
support tensor vibrations: the ``gluon''
($\ell = 2$ tensor mode) \emph{is} the
vibration that opens the third dimension, so
$d\geq 3$ is forced.  Confinement---the
inability of $\ell \geq 2$ modes to propagate
in the $\ell = 0$ background---then receives a
second interpretation: it is the mismatch
between a 2D membrane vibration and a 1D string
background, precisely the
``string--membrane asymmetry'' of
Section~\ref{sec:gluon}.

\textit{Lower bound vs.\ unique selection.}---%
Eq.~\eqref{eq:dim_argument} is strictly a
\emph{lower bound}: $d=3$ is the smallest
integer for which $\ell\geq 2$ modes are
angularly admissible on $S^{d-1}$.  Larger
spatial dimensions ($d=4, 5, \ldots$) also
support $\ell\geq 2$ tensor modes on their
angular sphere and therefore pass the same
structural test.  The argument above does not
by itself exclude $d>3$; the selection of
exactly $d=3$ over $d>3$ requires an
additional principle---an Occam-style
minimality postulate, stability arguments from
the main theory, or empirical input from the
observed angular structure of matter and
gravity.  We adopt Occam minimality as a
working postulate in the remainder of this
paper and label it as such, rather than as a
theorem derived from the $G_2(X)$ action alone.

%-------------------------------------------------------------------
\subsection{The ``processor'' picture:
a single quantum entity}
\label{sec:processor}
%-------------------------------------------------------------------

The one-dimensional reduction raises an
ontological question: is space a collection of
many vibrating elements, or is it a
\textbf{single quantum entity} that manifests
at every point?

A computer processor is physically one device,
yet it ``serves'' the entire machine: every
memory address, every running process.  It has a
clock frequency---a single fundamental tick---and
in each tick it updates the state at a particular
address.  The ISPG medium may be analogous:
$\phi$ is a single quantum field---not a
collection of point particles---and $\nu_0$ is
its clock frequency.  Non-locality (quantum
entanglement) is then trivial: there is no
distance to bridge between two ``parts'' of a
single entity.  The speed of light $c$ is the
maximum rate at which the entity can update
different ``addresses'' in its quantum state.

Mathematically, the field $\phi(x,t)$ already
encodes this picture: a single function
evaluated at all spacetime points.
The ontological step is to take the formalism
literally---$\phi$ is not ``a value at each
point of a pre-existing space'' but rather
\emph{the single quantum object whose
superposition state we parametrize by $x$}.
The apparent spatial extension of the universe
is the quantum state-space of one entity.

This interpretation changes no equation but
resolves several puzzles:
\begin{itemize}
  \item \textbf{Universality of $\nu_0$:}
    one entity, one frequency.
  \item \textbf{Quantum non-locality:}
    two entangled particles share a state
    of one field; no signal needs to ``travel.''
  \item \textbf{Continuity of space:}
    there are no gaps between particles---there
    is one object whose excitation pattern we
    call ``geometry.''
\end{itemize}

Whether the ``one entity'' interpretation is
merely a restatement of quantum field theory
or carries testable consequences (e.g.\ via
Bell-inequality violations in gravitational
degrees of freedom) is addressed in the next
subsection.

%-------------------------------------------------------------------
\subsection{From one entity to three-dimensional
space: the entanglement route}
\label{sec:entanglement_geometry}
%-------------------------------------------------------------------

The claim that $\phi$ is a single quantum
entity whose entanglement structure generates
the appearance of spatial extension can be
made precise using three results, the first
two of which are already established in the
literature.

\textit{Step~1: vacuum entanglement obeys an
area law.}---For any quantum field in its
vacuum state, the entanglement entropy between
a region $A$ and its complement scales as the
\emph{area} of the boundary, not the volume:
\begin{equation}\label{eq:area_law}
  S_{\mathrm{ent}}(A)
  \;=\; \frac{\mathcal{A}(\partial A)}
  {4\,\ell_{\mathrm{eff}}^2}
  \;+\; \text{sub-leading},
\end{equation}
where $\ell_{\mathrm{eff}}$ is a UV cutoff
length.  This is the Bombelli--Koul--Lee--Sorkin
area law \cite{Bombelli1986}, confirmed for free and
interacting fields alike.  In ISPG the role of
$\ell_{\mathrm{eff}}$ is played by the oscillon
core radius $R \sim c / \nu_0$.

\textit{Step~2: entanglement area $=$ geometric
area (Ryu--Takayanagi \cite{RyuTakayanagi2006,NishiokaTakayanagi2009}).}---In holographic
theories \cite{Maldacena1998} the entanglement entropy of a boundary
region $A$ equals the area of the minimal
surface in the bulk that is homologous to $A$:
\begin{equation}\label{eq:RT}
  S_{\mathrm{ent}}(A)
  \;=\; \frac{\text{Area}(\gamma_A)}{4\,G_N}\,,
\end{equation}
where $\gamma_A$ is the Ryu--Takayanagi surface
and $G_N$ is Newton's constant.
This means the \emph{geometry itself}---distances,
angles, curvature---is encoded in, and can be
reconstructed from, the entanglement pattern of
the underlying quantum state.  The ER\,=\,EPR
conjecture \cite{MaldacenaSusskind2013}
sharpens this: every entanglement link is a
microscopic wormhole; geometry \emph{is}
entanglement.

\textit{Step~3 (ISPG-specific).}---The single
field $\phi$, vibrating at $\nu_0$ with the
$G_2(X)$ nonlinearity, produces a vacuum state
$|\Omega\rangle$ whose entanglement structure
satisfies three conditions.  We verify each:

\begin{enumerate}
\item \textbf{Bekenstein--Hawking matching.}\;
The UV cutoff of the ISPG medium is the
Planck-scale lattice spacing
$a = \pi\,\ell_P$.
(This is the shortest wavelength the medium
supports; it coincides numerically with the
classical estimate $\lambda_0^{(\mathrm{cl})}$
of Eq.~\ref{eq:nu0_value}, but the physical
reason it enters here is that $\ell_P$ is the
gravitational UV cutoff, not that the classical
$\nu_0^{(\mathrm{cl})}$ sets the dynamics---the
quantum value $\nu_0$ governs all oscillonic
physics as stated in
\S\ref{sec:nu0_determination}.)
The Srednicki area law for a free scalar gives
$S = \alpha_0\,\mathcal{A} / a^2$ with
$\alpha_0 \approx 0.30$.  With $a = \pi\ell_P$:
\begin{equation}\label{eq:S_ent_free}
  S_{\mathrm{free}}
  = \frac{0.30\;\mathcal{A}}{\pi^2\,\ell_P^2}
  = \frac{0.030\;\mathcal{A}}{\ell_P^2}\,.
\end{equation}
The Bekenstein--Hawking entropy requires
$S_{\mathrm{BH}} = \mathcal{A}/(4\ell_P^2)$,
i.e.\ $\alpha = \pi^2/4 = 2.467$.  The
$G_2(X)$ nonlinearity provides the enhancement
through two mechanisms:
\begin{itemize}
  \item \textit{Sound-speed reduction.}\;
    For a k-essence Lagrangian $G_2(X)$,
    perturbations propagate at the effective
    sound speed (established in \S\ref{sec:q_derivation}):
    \begin{equation}\label{eq:cs_ent}
      c_s^2 = \frac{G_{2,X}}{G_{2,X} + 2X\,G_{2,XX}}
      = \frac{1}{1 + 4q}\,.
    \end{equation}
    The entanglement entropy of a scalar with
    sound speed $c_s < c$ scales as $1/c_s$
    (because the effective UV cutoff in momentum
    space is $k_{\max} = \omega_{\max}/c_s$,
    increasing the number of modes per unit area).
    With $q = 1.853$:
    \begin{equation}\label{eq:Fcs}
      F_{c_s}
      = \frac{1}{c_s}
      = \sqrt{1 + 4q}
      = \sqrt{8.412}
      \approx 2.900\,.
    \end{equation}

  \item \textit{Floquet sideband enhancement.}\;
    The oscillating background
    $\phi_0\cos(\omega_0 t)$ means that vacuum
    fluctuations $\delta\phi$ satisfy a Mathieu
    equation (\S\ref{sec:mathieu}).  By the
    Floquet--Bloch theorem, each stable mode is
    not a pure plane wave but a superposition of
    frequency sidebands:
    \begin{equation}\label{eq:floquet_expansion}
      \delta\phi_k(t)
      = e^{i\nu_k t}\!\sum_{n=-\infty}^{\infty}
      c_{2n}^{(k)}\,e^{2in\omega_0 t},
    \end{equation}
    where $\nu_k$ is the Floquet exponent and the
    Fourier coefficients $c_{2n}$ satisfy the
    Mathieu three-term recurrence.

    For a free field ($q = 0$): only $c_0 \neq 0$,
    $|c_0|^2 = 1$---one channel per mode.
    For $q \neq 0$: the parametric driving
    distributes spectral weight among sidebands.
    Each sideband at frequency $\nu_k + 2n\omega_0$
    is an independent entanglement channel across
    any spatial bipartition, because entanglement
    entropy is additive over orthogonal frequency
    components.

    The enhancement factor is the
    \textbf{inverse participation ratio} (IPR)
    of the ground-state Mathieu function
    $\mathrm{ce}_0(z, q)$:
    \begin{equation}\label{eq:Fbands_def}
      F_{\mathrm{bands}}
      = \frac{1}{\bigl|A_0^{(0)}(q)\bigr|^2}\,,
    \end{equation}
    where $A_0^{(0)}$ is the zeroth Fourier
    coefficient of $\mathrm{ce}_0$, normalised
    so that $|A_0|^2 + 2\sum_{r=1}|A_{2r}|^2 = 1$.
    Physically: $|A_0|^2$ is the fraction of
    spectral weight remaining in the unshifted
    ($n=0$) channel; $1/|A_0|^2$ counts how many
    effective channels share the weight.

    The coefficient $A_0^{(0)}(q)$ is determined
    by solving the Mathieu recurrence
    \begin{equation}\label{eq:mathieu_recurrence}
      a_0\,A_0 = q\,A_2,\qquad
      (a_0 - 4n^2)\,A_{2n}
      = q\bigl(A_{2n-2} + A_{2n+2}\bigr),
    \end{equation}
    with the eigenvalue $a_0(q)$ fixed by the
    convergence condition $A_{2n} \to 0$ as
    $n \to \infty$.  At $q = 1.853$ the numerical
    solution gives
    \begin{equation}\label{eq:Fbands_value}
      \bigl|A_0^{(0)}(1.853)\bigr|^2 \approx 0.353,
      \qquad
      F_{\mathrm{bands}} \approx 2.84\,.
    \end{equation}
    This value can be verified independently:
    the recurrence
    \eqref{eq:mathieu_recurrence} is the same
    system already solved in
    \S\ref{sec:mathieu_numerical}---no new
    input is required.
\end{itemize}
The total enhancement is
\begin{equation}\label{eq:F_total}
  F = F_{c_s} \times F_{\mathrm{bands}}
  = \sqrt{1+4q}\;\times\;
  \frac{1}{\bigl|A_0^{(0)}(q)\bigr|^2}
  \approx 2.90 \times 2.84 = 8.22\,,
\end{equation}
giving the effective Srednicki coefficient
$\alpha_{\mathrm{eff}} = 0.30 \times 8.22
= 2.47 \approx \pi^2/4 = 2.467$
and therefore
\begin{equation}\label{eq:S_ent_ISPG}
  \boxed{
  S_{\mathrm{ent}}
  = \frac{\alpha_{\mathrm{eff}}}
  {\pi^2\,\ell_P^2}\,\mathcal{A}
  = \frac{\pi^2/4}{\pi^2\,\ell_P^2}\,\mathcal{A}
  = \frac{\mathcal{A}}{4\,\ell_P^2}
  = S_{\mathrm{BH}}\,.
  }
\end{equation}
Both factors---$F_{c_s}$ and
$F_{\mathrm{bands}}$---are determined by
the single parameter $q = 1.853$, which is
itself fixed by the muon mass
(\S\ref{sec:q_derivation}).  No additional
free parameter enters the entropy calculation.

\textit{Status of the match.}---Equation
\eqref{eq:S_ent_ISPG} should be read as a
\emph{consistency check}, not as a first-principles
derivation of $S_{\mathrm{BH}}$.  Two ingredients in
the enhancement factor require further
substantiation:
\begin{itemize}
  \item $F_{c_s} = 1/c_s$ is assumed on the ground
    that the effective momentum cutoff scales as
    $k_{\max} = \omega_{\max}/c_s$; the rigorous
    extension of Srednicki's flat-space, $c_s=c$
    area-law derivation to an oscillating $G_2(X)$
    background is not carried out here.
  \item $F_{\mathrm{bands}} = 1/|A_0^{(0)}(q)|^2$
    assumes that the Floquet sidebands contribute
    as \emph{independent} entanglement channels.
    Whether the sidebands behave as orthogonal
    Fock-space factors or as phase-coherent
    Fourier components of a single mode depends on
    the full Bogoliubov structure of the
    time-periodic vacuum; that calculation is a
    self-contained technical task and the topic of
    a separate work.
\end{itemize}
The 0.1\% agreement $\alpha_{\mathrm{eff}}\approx
\pi^2/4$ is therefore a numerical coincidence of
the current parametric estimates: it is
\emph{compatible} with the Bekenstein--Hawking
coefficient but does not by itself constitute a
derivation of it.  This matches the general
epistemic posture of the paper: the BH match is
part of the area-law scaffolding, not a core
gravitational prediction of the main
theory~\cite{Lotishvili2026}.

\item \textbf{Dimensional selection.}\;
The entanglement pattern is that of a
$(3{+}1)$-dimensional manifold because
the $\ell = 0$ (scalar) and $\ell = 2$
(tensor) sectors coexist \emph{only} for
$d \geq 3$ spatial dimensions
(\S\ref{sec:dim_emergence}), and the
Occam minimum selects $d = 3$.

\item \textbf{Scalar--tensor dichotomy.}\;
The vacuum two-point function
(verified numerically in
\texttt{python/entanglement\_proof.py},
Section~3) shows:
\begin{align}
  \langle\delta\phi(\mathbf{x})\,
  \delta\phi(\mathbf{0})\rangle
  &\;\propto\; \frac{1}{|\mathbf{x}|^2}
  \qquad (\text{power law, long-range}),
  \label{eq:scalar_corr}\\[4pt]
  \langle h_{ij}^{TT}(\mathbf{x})\,
  h_{kl}^{TT}(\mathbf{0})\rangle
  &\;\propto\; \frac{e^{-\kappa|\mathbf{x}|}}
  {|\mathbf{x}|}
  \qquad (\text{exponential, confined}),
  \label{eq:tensor_corr}
\end{align}
with $\kappa = 2\pi\sqrt{\nu_N^2 - \nu_0^2}/c$.
Using $\nu_N/\nu_0 = j_{2,1}/\pi = 1.835$,
$\kappa\,R_{\mathrm{osc}} \approx 9.66$:
at $r = 5\,R_{\mathrm{osc}}$ the scalar
correlator is $10^{-2}$ while the tensor
is $10^{-20}$---a factor $10^{18}$ suppression.
Confinement is an entanglement property: tensor
modes decohere exponentially outside the
oscillon, while scalar modes maintain long-range
quantum correlations that weave the spatial
fabric.
\end{enumerate}

\textit{The proof chain (schematic).}---Steps
1--3, taken together with the RT reconstruction
step that \emph{remains to be carried out}
(below), would establish that the three spatial
dimensions are an \textbf{output}, not an
input, of the formalism.  Schematically:
\begin{equation}\label{eq:proof_chain}
  \boxed{
  \text{one }\phi
  \;\xrightarrow{\;\nu_0\text{ vibration}\;}
  |\Omega\rangle
  \;\xrightarrow{\;\text{entanglement}\;}
  S = \frac{\mathcal{A}}{4\ell_P^2}
  \;\xrightarrow{\;\text{RT}\;}
  g_{\mu\nu}\;(3{+}1)
  }
\end{equation}
The final arrow---the explicit tensor-network
reconstruction of $g_{\mu\nu}$ from the
entanglement pattern of $|\Omega\rangle_\phi$---%
is the \emph{missing} step; steps~1--3 assemble
the scaffolding (area-law matching,
dimensional lower bound, scalar--tensor
correlator dichotomy) under which that
reconstruction could be attempted, but the
reconstruction itself is not executed in this
paper.  The slogan ``space is the entanglement
shadow of $\phi$'' should therefore be read as
the \emph{target} of the programme, not as a
theorem established by Eqs.~\eqref{eq:S_ent_ISPG}%
--\eqref{eq:tensor_corr}; its status is the
same as that of ER\,=\,EPR in the general
holography literature.

\textit{Observational context.}---Three
independent observations are \emph{compatible}
with the single-entity reading developed above,
though each is equally compatible with a
conventional quantum-field-theory reading---%
they constrain, rather than select between,
the two viewpoints:
\begin{itemize}
  \item \textbf{Universality of $\nu_0$:}\;
    if $\nu_0$ is taken as genuinely universal
    then no synchronisation mechanism is
    required.  This is descriptive of any
    single-field theory (including ordinary
    QFT) and does not by itself pick out the
    ``one entity'' interpretation.
  \item \textbf{Spatial continuity:}\;
    Fermi-LAT constrains Lorentz-violating
    dispersion of gamma rays from cosmological
    distances down to the Planck scale.  A
    single continuous entity is consistent with
    this constraint, but so is any standard
    continuum field theory with exact Lorentz
    symmetry; the observation excludes strongly
    granular models, not the two readings being
    compared here.
  \item \textbf{Quantum non-locality:}\;
    Bell-inequality violations (Aspect,
    Clauser, Zeilinger, Nobel 2022) are
    reproduced by the Fock-space structure of
    ordinary entangled multi-particle states in
    standard QFT; they are equally consistent
    with the single-entity reading, but do not
    discriminate between the two.
\end{itemize}
These observations therefore \emph{allow} the
single-entity interpretation without uniquely
picking it out over the standard-QFT reading.
A genuinely discriminating signature of the
single-entity picture (if any) would live in
gravitational degrees of freedom and is an
open question, not a consequence of the
observations listed above.

The next step is the full
Ryu--Takayanagi reconstruction: given
$|\Omega\rangle_\phi$ and its entanglement
structure, build the $(3{+}1)$-dimensional
metric $g_{\mu\nu}$ explicitly using
tensor-network (MERA) techniques.  We
emphasise that this reconstruction is
\textbf{not} carried out in the present
paper: steps 1--3 above assemble necessary
ingredients (area law, dimensional lower
bound, confinement dichotomy), but the
explicit construction of $g_{\mu\nu}$ from
$|\Omega\rangle$ is a self-contained
tensor-network computation and is the topic
of a separate work.  Until that computation
is carried out, the claim ``space is the
entanglement shadow of $\phi$'' has the
status of a structural conjecture supported
by the scaffolding of
Eqs.~\eqref{eq:S_ent_ISPG}--%
\eqref{eq:tensor_corr}, not of a derived
theorem.

\textit{A speculative corollary: other notes,
other universes.}---If our cosmos is a string
vibrating at the note $\nu_0$, nothing in the
formalism forbids the existence of other
``strings'' vibrating at different fundamental
frequencies $\nu_0'$.  Each such string would
generate its own harmonic tower, its own
oscillon spectrum, and therefore its own
particle physics and chemistry.  Two strings at
incommensurate frequencies would be mutually
transparent---unable to exchange coherent
energy---and would coexist in the same
``space'' without interacting, much as two
radio stations on different frequencies share
the same air.  In this picture, a multiverse
is not a collection of spatially separated
bubble universes but a \emph{superposition of
co-located vibrational modes at different
fundamental pitches}---each pitch defining a
self-contained physics, invisible to the
others.

%-------------------------------------------------------------------
\subsection{Electric charge as transverse
symmetry breaking}
\label{sec:charge}
%-------------------------------------------------------------------

The dimensional-emergence picture suggests a
natural origin for electric charge that requires
no new fields or parameters.

\textit{The transverse plane and its
symmetry.}---The $\nu_0$-string vibrates
transversely, creating a two-dimensional plane
(\S\ref{sec:dim_emergence}).  This vibration is
\textbf{symmetric}: the string displaces equally
in all transverse directions.  The oscillon---a
localized, nonlinear resonance of the same
string---may or may not preserve this symmetry:

\begin{itemize}
  \item \textbf{Symmetric oscillon}
    ($\ell = 0$ only):
    the localized vibration pattern has no
    preferred transverse direction.
    The oscillon does not couple to transverse
    travelling waves (photons).
    $\to$ \textbf{electrically neutral.}

  \item \textbf{Asymmetric oscillon}
    ($\ell = 1$ dipole component):
    the $G_2(X)$ nonlinearity, or the internal
    structure of the resonance, breaks the
    transverse reflection symmetry.  The oscillon
    vibrates ``more'' in one transverse direction
    than the opposite.  This dipolar pattern
    couples directly to transverse waves.
    $\to$ \textbf{electrically charged.}
\end{itemize}

The \textbf{sign} of the charge is the direction
of the transverse asymmetry.
Because the transverse plane is two-dimensional,
the asymmetry direction can be parametrised by
an angle $\theta \in [0, 2\pi)$.  Rotating all
asymmetry directions by the same angle leaves the
physics unchanged---this is a \textbf{global
U(1) symmetry}.  Promoting it to a local symmetry
(the direction may vary from point to point)
requires a compensating field that transports the
angle information: this is the
\textbf{photon as a U(1) gauge boson}.

\textit{Consequences.}---
\begin{enumerate}
  \item \textbf{Two signs:}\;
    the asymmetry has two stable orientations
    (related by $\pi$ rotation) $\to$
    positive and negative charge.

  \item \textbf{Charge conservation:}\;
    total transverse asymmetry cannot change;
    creating a $+$ asymmetry requires
    simultaneously creating $-$ $\to$
    pair production.

  \item \textbf{Antimatter:}\;
    both orientations are equally probable
    $\to$ every particle has an antiparticle
    with opposite charge.

  \item \textbf{Opposite charges attract:}\;
    a $+$ and a $-$ oscillon together restore
    transverse symmetry between them
    (destructive interference of asymmetries)
    $\to$ lower energy $\to$ attraction.

  \item \textbf{Like charges repel:}\;
    two $+$ oscillons reinforce the asymmetry
    (constructive interference)
    $\to$ higher energy $\to$ repulsion.

  \item \textbf{Neutral particles are invisible
    to photons:}\;
    a symmetric ($\ell = 0$) oscillon has no
    transverse dipole; the transverse wave
    passes through unaffected.

  \item \textbf{Fractional charges (quarks):}\;
    inside a baryon the $D_3$ symmetry
    (\S\ref{sec:gluon}) allows the asymmetry
    direction to lock at $120°$ intervals
    ($2\pi/3$ phase steps), yielding charges
    quantised in thirds.
    The detailed mapping to $+2/3$ and $-1/3$
    requires solving the coupled
    scalar--vector mode equation inside the
    $D_3$ cavity---a task for future work.
\end{enumerate}

\textit{Electromagnetic force as transverse
Bernoulli effect.}---Gravity in ISPG arises
from a pressure deficit in the \emph{scalar}
(longitudinal) sector: oscillon tails deplete
the $\nu_0$ background, lowering the local
pressure and producing attraction
(\S\ref{sec:ontology}).  The electromagnetic
force is the same Bernoulli mechanism operating
in the \emph{transverse} sector:
\begin{itemize}
  \item \textbf{Opposite charges:}\;
    between a $+$ and a $-$ oscillon the
    transverse asymmetries are anti-parallel
    and partially cancel (destructive
    interference).  The transverse energy density
    between them is \emph{reduced} $\to$ lower
    transverse pressure $\to$ the oscillons are
    pushed together.  \textbf{Attraction.}

  \item \textbf{Like charges:}\;
    between two $+$ oscillons the asymmetries
    are parallel and reinforce (constructive
    interference).  The transverse energy density
    between them is \emph{increased} $\to$ higher
    transverse pressure $\to$ the oscillons are
    pushed apart.  \textbf{Repulsion.}
\end{itemize}

\noindent
The unifying picture---four forces as four
vibration modes of the $\nu_0$-string---is:
\begin{equation}\label{eq:bernoulli_unified}
\boxed{
\begin{aligned}
  &\text{Gravity}
  &&=\; \text{longitudinal (compression)}
  &&\quad\ell = 0\;\text{scalar Bernoulli},\\
  &\text{Electromagnetism}
  &&=\; \text{transverse (displacement)}
  &&\quad\ell = 1\;\text{transverse Bernoulli},\\
  &\text{Weak force}
  &&=\; \text{torsional (twist)}
  &&\quad\text{chiral, pseudoscalar},\\
  &\text{Strong force}
  &&=\; \text{membrane (2D vibration)}
  &&\quad\ell \geq 2\;\text{confined tensor}.
\end{aligned}
}
\end{equation}

\textit{Magnetism.}---A magnetic field is not
a separate phenomenon; it is the relativistic
manifestation of the electric (transverse)
asymmetry pattern when the charged oscillon
\emph{moves}.  Because ISPG already derives
Lorentz time dilation from frequency invariance
(\S\ref{sec:freq_invariance}), the full
electromagnetic tensor $F_{\mu\nu}$ emerges
automatically: a static transverse asymmetry
is an electric field; the same asymmetry seen
from a boosted frame acquires a magnetic
component.  No additional mechanism is required.

\textit{Future work: from qualitative to
quantitative electromagnetism.}---The structural
elements above (symmetry breaking $\to$ charge,
2D transverse plane $\to$ U(1), Bernoulli
attraction/repulsion) are fixed by the geometry
and require no free parameters.  What remains is
to close two quantitative gaps:

\begin{enumerate}
\item \textbf{Charge assignment.}\;
  The $G_2(X)$ nonlinearity or the Mathieu mode
  structure must determine which oscillons develop
  an $\ell = 1$ transverse component (charged)
  and which remain purely $\ell = 0$ (neutral).
  Concretely, the Mathieu stability chart for
  the $\ell = 1$ sector has its own critical $q$
  threshold; oscillons whose $N$ falls in the
  $\ell = 1$ stable band are charged.  This is a
  well-posed eigenvalue problem once the
  $\ell = 1$ effective potential is extracted
  from $G_2(X)$.

\item \textbf{Fine-structure constant.}\;
  $\alpha$ should emerge as the ratio of the
  $\ell = 1$ dipole amplitude to the $\ell = 0$
  monopole amplitude:
  \begin{equation}\label{eq:alpha_target}
    \alpha
    = \frac{|\Phi_{\ell=1}|^2}
           {|\Phi_{\ell=0}|^2}
    \;\stackrel{?}{=}\;
    \frac{1}{137.036}\,.
  \end{equation}
  This ratio is computable from the same lattice
  simulation that determines $q$ and $\eta_T$.
  It constitutes the decisive quantitative test
  of the electromagnetic sector.
\end{enumerate}

Both computations are lattice tasks within the
existing $G_2(X)$ framework; no new theoretical
ingredient is required.

%-------------------------------------------------------------------
\subsection{The Higgs boson as amplitude
fluctuation}
\label{sec:higgs}
%-------------------------------------------------------------------

The companion paper gives the bi-conformal scaling of an
effective electroweak VEV
\cite{Higgs1964,EnglertBrout1964,Guralnik1964,Goldstone1961,ColemanWeinberg1973}
as conformal symmetry breaking---showing that
$v(\varphi) = v_0\,e^{\varphi/2}$, recovering
the broken-phase electroweak gauge-boson mass matrix
once the SM electroweak sector is supplied, and recording
the gravitational correction
$\delta m_H^2 \sim 10^{-34}\,m_H^2$---is given
in the companion paper \cite{Lotishvili2026Q} (ISPG\_Quantum, Sec.~8).
Here we present the oscillon-level picture that
motivates the candidate ISPG replacement.

\medskip
Throughout this paper a single theme recurs:
\textbf{mass is the magnitude of vibration}.
An oscillon's mass is its harmonic number $N$;
$m_N = m_1\,b_N(q)$ is, at its core, a
statement about how vigorously the
$\nu_0$-string vibrates in the trapped region.
This theme identifies the target that a Higgs-replacement
mechanism must reproduce.

\textit{The oscillon amplitude as candidate order
parameter.}---The $\nu_0$-string
vibrates everywhere with amplitude $\Phi_0$:
\begin{equation}\label{eq:higgs_vev}
  \phi(x,t) \;\sim\;
  \Phi_0\,\cos(2\pi\nu_0\,t).
\end{equation}
If $\Phi_0 = 0$ (the string is at rest), the
$G_2(X)$ nonlinearity has nothing to act on:
no self-focusing occurs, no oscillons form,
and every excitation propagates as a massless
wave.  The moment $\Phi_0 \neq 0$, the
nonlinear term $\alpha X^2/M^2$ activates,
oscillons condense, and particles acquire mass.
This is the ISPG analogue of the role played by the
Standard Model's Higgs vacuum expectation value $v$:
with $v = 0$ the SM fermions and gauge bosons are
massless; with $v \neq 0$ they acquire mass
proportional to their coupling.

The normalization target is:
\begin{equation}\label{eq:higgs_map}
\boxed{
  \Phi_0
  \;\longrightarrow\;
  v_{\mathrm{EW}} = 246\;\text{GeV}
  \qquad
  \text{(open coarse-graining map)}.
}
\end{equation}
\noindent
Here $\Phi_0$ is dimensionless oscillon-profile data,
whereas $v_{\mathrm{EW}}$ is an imported electroweak
energy scale until the normalization map is derived.

\textit{Effective amplitude symmetry breaking.}---The
$G_2(X)$ kinetic structure suggests an equilibrium
amplitude $\Phi_0$
(Eq.~\ref{eq:amplitude_constraint}).
The linear term
$-M_{\mathrm{Pl}}^2\,X/2$ penalises large
amplitudes; the nonlinear term
$\alpha\,X^2/M^2$ rewards them.  The
competition produces an effective potential
$V_{\mathrm{eff}}(\Phi_0)$ with a minimum
at $\Phi_0 \neq 0$:
\[
  V_{\mathrm{eff}}(\Phi_0)
  \;\sim\;
  -\mu^2\,\Phi_0^2
  + \lambda\,\Phi_0^4,
  \qquad
  \mu^2,\;\lambda > 0,
\]
where $\mu^2$ and $\lambda$ are determined by
$\alpha$, $M$, $M_{\mathrm{Pl}}$, and
$\omega_0$ in the effective amplitude description.
This is the \textbf{Mexican-hat analogy}: the
$G_2(X)$ kinetic structure suggests an effective
amplitude potential in the broken-phase regime, but a
genuine non-derivative SU(2)-doublet Higgs potential is
not derived here.
The string ``chooses'' to vibrate at $\Phi_0
\neq 0$ rather than remaining at rest:
this is the candidate ISPG analogue of spontaneous
symmetry breaking.

\textit{The Higgs boson.}---A localised
perturbation in which the vibration amplitude
deviates from its equilibrium value,
$\Phi_0 \to \Phi_0 + h(x,t)$,
oscillates around the potential minimum.
The frequency of this oscillation is the
\textbf{candidate amplitude-mode Higgs mass}:
\begin{equation}\label{eq:higgs_mass}
  m_H^{(\mathrm{amp})}
  \;=\; \sqrt{V_{\mathrm{eff}}''(\Phi_0)}.
\end{equation}
The observed value $m_H=125$~GeV is currently a
target/ladder hit, not a first-principles prediction
from this effective amplitude potential.
In signal-processing language: the
$\nu_0$-vibration is the \textbf{carrier wave};
the Higgs boson is its
\textbf{amplitude modulation} (AM).

\textit{Coupling proportional to mass.}---Heavier
particles are higher harmonics of the
$\nu_0$-string.  Higher harmonics depend more
sensitively on the nonlinear trapping strength,
which is set by $\Phi_0$.  A fluctuation
$h(x,t)$ therefore affects heavy oscillons more
than light ones:
\begin{equation}
  g_{Hff} \;\propto\; \frac{\partial m_N}
  {\partial \Phi_0},
\end{equation}
which must be tested by explicitly differentiating the
Mathieu law $m_N=m_1 b_N(q(\Phi_0))$.  SM-like
proportionality to $m_N$ is not automatic; it requires
verifying that
$(\partial b_N/\partial q)(\partial q/\partial\Phi_0)$
reproduces the Yukawa coupling pattern.  This derivative
test is open future work.

\textit{Goldstone bosons and $W/Z$ masses.}---The
field $\phi$ is a single \emph{real} scalar,
but its transverse perturbations carry two
circular-polarisation degrees of freedom
$\delta\phi_\pm \propto e^{\pm i\vartheta}$
(\S\ref{sec:torsion_derivation}).  The full
low-energy field content is therefore:
\begin{equation}\label{eq:higgs_dof}
  \underbrace{\Phi_0}_{\text{amplitude (real)}}
  \;+\;
  \underbrace{\theta_1,\;\theta_2,\;
  \theta_3}_{\text{polarisation angles
  (SU(2) phase space)}},
\end{equation}
where the three $\theta_i$ parametrise the
orientation of the polarisation ellipse in
the two-dimensional transverse plane and its
coupling to the longitudinal sector
(\S\ref{sec:WZ_bosons}).  No new field is
introduced: the $\theta_i$ are collective
coordinates of the same $\phi$.

Fluctuations of $\Phi_0$ cost energy
(the effective amplitude curvature)
$\to$ \textbf{candidate amplitude-mode Higgs boson}.
Fluctuations of $\theta_i$ are the proposed collective
route to the three would-be Goldstone directions.  A
completed E7 construction must still show that these
collective coordinates form the correct
$SU(2)_L\times U(1)_Y/U(1)_{\mathrm{EM}}$ coset, couple
gauge-invariantly to the vector sector, and reproduce the
standard broken-phase mass matrix.  Only after that step
can they be identified with Goldstone modes that are
``eaten'' by $W^\pm$ and $Z^0$.  At present the observed
$W/Z$ masses remain conditional on the imported SM
electroweak data and on the open
$\Phi_0\to v_{\mathrm{EW}}$ normalization.

\textit{The electroweak coincidence.}---The
geometric mean of the Planck mass and the
electron mass is
$\sqrt{m_P\,m_e} \approx 79$~GeV, remarkably
close to $m_W = 80$~GeV
(\S\ref{sec:mathieu_numerical}).
In the ISPG picture this points to a common structural
origin: the electroweak scale $v_{\mathrm{EW}}$ is an
open normalization target for $\Phi_0$, which is fixed by
$\omega_0 = 2\pi\nu_0$ and the $G_2(X)$
parameters
(Eq.~\ref{eq:amplitude_constraint})---the same
quantities that determine $m_e$ (the first
stable harmonic) and $m_P$ (the Planck scale).
All three scales share a common origin.

\textit{Future work.}---The Higgs mass would be obtained
from $\sqrt{V_{\mathrm{eff}}''(\Phi_0)}$ once the
effective amplitude potential, the
$\Phi_0\to v_{\mathrm{EW}}$ normalization, and the
$\Phi_0$-dependence of the Mathieu coefficient
$b_N(q(\Phi_0))$ are derived.  The observed
$m_H=125$~GeV is currently a target/ladder hit, not a
first-principles prediction.  The precise mapping
between $\Phi_0$ and $v_{\mathrm{EW}} = 246$~GeV
through the bi-conformal normalisation conventions is a
natural part of that calculation.

%-------------------------------------------------------------------
\subsection{Summary}
\label{sec:dim_summary}
%-------------------------------------------------------------------

\begin{center}
\fbox{\parbox{0.90\linewidth}{
\textbf{Dimensional emergence in ISPG.}\\[4pt]
1.\ The $\nu_0$-background has an
integer-harmonic spectrum ($j_0$ zeros $=
n\pi$)---the signature of a 1D string.\\[2pt]
2.\ Vibration of the string requires
one transverse direction $\to$ 2D.\\[2pt]
3.\ Tensor ($\ell \geq 2$) modes require
a second transverse direction $\to$ 3D.\\[2pt]
4.\ Three dimensions is the
\textbf{minimum} for both scalar and tensor
modes; no further dimension is generated.\\[2pt]
5.\ Confinement is the incompatibility
between the 2D vibration (gluon) and the 1D
background ($\nu_0$-string).\\[2pt]
6.\ $\nu_0$ is universal because the
substrate is a single quantum entity (one
string, one pitch).\\[2pt]
7.\ Space $=$ entanglement shadow of $\phi$:
$|\Omega\rangle \to S \propto A \to g_{\mu\nu}$
(Ryu--Takayanagi route,
\S\ref{sec:entanglement_geometry}).\\[2pt]
8.\ Electric charge $=$ transverse symmetry
breaking of the oscillon; U(1) $=$ rotational
freedom of the 2D transverse plane
(\S\ref{sec:charge}).\\[2pt]
9.\ Neutrino $=$ polarisation-rotation wave
(helical vortex, $n = \pm 1$);
weak force $=$ nonlinear
torsion--transverse coupling; parity violation
$=$ inherent chirality of torsion
(\S\ref{sec:neutrinos}).\\[2pt]
10.\ \textbf{Four forces $=$ four vibration
modes} of one string: longitudinal (gravity),
transverse (EM), torsional (weak), membrane
(strong).\\[2pt]
11.\ Conditional scaling:
$v(\varphi)=v_0\,e^{\varphi/2}$; the oscillon
amplitude $\Phi_0$ is the candidate ISPG analogue
of the Higgs order parameter, with
$\Phi_0\to v_{\mathrm{EW}}$ an open normalization
target.  $G_2(X)$ suggests an effective amplitude
potential; a non-derivative SU(2)-doublet potential
is not derived here (\S\ref{sec:higgs}).\\[6pt]
\textit{Companion-paper results and candidate routes
(ISPG\_Quantum):}\\[2pt]
12.\ Spin-1/2 structural route $=$ vortex winding number,
$J = n\hbar/2$ under the conditional spin bridge;
Fermi--Dirac/Pauli structure as a candidate route from
Berry phase, with the rigorous QFT/Fock-space proof open
(Sec.~10).\\[2pt]
13.\ Effective Higgs-sector scaling:
$v(\varphi) = v_0\,e^{\varphi/2}$; the broken-phase
mass matrix is recovered once the SM electroweak
sector and imported $v_0$ are supplied (Sec.~8).\\[2pt]
14.\ Koide formula $Q = 2/3$: numerical support
from the E1D/rank-1 cavity analysis on the MASS/Mathieu
lepton ladder, $0.009\%$ accuracy; the 3D-to-1D bridge
from raw cavity eigenfrequencies to lepton masses remains
open (Sec.~7.2).\\[2pt]
15.\ Inertia: $F = -m_{\mathrm{eff}}\,a$ from
retarded oscillon potential;
$m_{\mathrm{inertial}} \equiv
m_{\mathrm{gravitational}}$ (Sec.~5).\\[2pt]
16.\ Dark-energy scale estimate (conditional):
$\rho_{\mathrm{vac}} \sim 10^{-121}\,
M_{\mathrm{Pl}}^4$ on the perturbative branch
(Sec.~6).\\[2pt]
17.\ Gravitational decoherence:
$\Gamma_{\mathrm{dec}} = 2m_0|\Delta\Phi_N|/h$,
$\sim 0.5$~Hz for eV-scale photons (Sec.~9).
}}
\end{center}

This section is primarily conceptual; the
mathematical foundations---Bessel-function
identities, gravitational-wave counting in
$d$ dimensions, the spherical-cavity
decomposition---are all standard results.
The new content is the \emph{chain of inference}:
from $j_0(x) = \sin(x)/x$ to the emergence
of exactly three spatial dimensions and the
confinement of tensor modes.  The quantitative
programme (proving that the one-dimensional
substrate generates exactly the observed
3+1-dimensional spacetime with its symmetry
groups) is a goal for future work.

\section{Future Directions}

\subsection*{Completed (in this paper)}

\begin{itemize}
  \item[\checkmark] Oscillon wave equation solution:
    $m_N = m_1 b_N(q)$, Koide $0.02\%$,
    muon $< 0.001\%$
    (\S\ref{sec:numerical_results},
    \S\ref{sec:mathieu_numerical},
    \S\ref{sec:q_derivation}).
  \item[\checkmark] $q$ is universal---
    sector-specific $q$ is excluded
    by the electron stability condition.
  \item[\checkmark] Glueball mass ratios from
    Bessel-function zeros: $\leq 2\%$ for four
    states, parameter-free
    (\S\ref{sec:gluon}).
  \item[\checkmark] Bekenstein--Hawking entropy
    matching: $S_{\mathrm{ent}} = \mathcal{A}
    /(4\ell_P^2)$ from $G_2(X)$ sound-speed
    reduction and Mathieu-band summation
    (\S\ref{sec:entanglement_geometry}).
\end{itemize}

\subsection*{Companion-paper results and conditional estimates
(ISPG\_Quantum)}

\begin{itemize}
  \item[$\triangleright$] \textbf{Spin-1/2 and
    Fermi--Dirac statistics (candidate topological route)}:
    $J = n\hbar/2$ from vortex winding under the
    conditional spin bridge; Berry phase $= \pm\pi$
    for $n = \pm 1$ exchange as a structural Pauli/Fermi--Dirac
    route; the full QFT/Fock-space proof remains imported/open
    (Sec.~10).
  \item[$\triangleright$] \textbf{Effective Higgs-sector scaling}:
    $v(\varphi) = v_0\,e^{\varphi/2}$ as conditional
    conformal scaling of an imported electroweak scale;
    the $\Phi_0\to v_{\mathrm{EW}}$ normalization and
    first-principles Higgs-mass derivation remain open
    (Sec.~8).
  \item[\checkmark] \textbf{Koide formula}:
    $Q = 2/3$ numerical support via rank-1
    perturbation and Hartree analysis on the MASS/Mathieu
    lepton ladder, $0.009\%$ accuracy; the 3D-to-1D
    bridge from raw cavity eigenfrequencies (which give
    $Q\approx 0.336$, not $2/3$) remains open
    (Sec.~7.2).
  \item[\checkmark] \textbf{Inertia and
    equivalence principle}: $F = -m_{\mathrm{eff}}
    \,a$ derived from retarded potential of
    accelerated oscillon;
    $m_{\mathrm{inertial}} \equiv
    m_{\mathrm{gravitational}}$ (Sec.~5).
  \item[$\circ$] \textbf{Dark-energy scale estimate
    (conditional)}:
    echo-background energy self-regulated by a negative-feedback
    loop, yielding $\rho_{\mathrm{vac}}
    \sim 10^{-121}\,M_{\mathrm{Pl}}^4$ on the perturbative
    branch, while zero-point decoupling and non-perturbative
    normalization remain open (Sec.~6).
  \item[\checkmark] \textbf{Gravitational
    decoherence}: $\Gamma_{\mathrm{dec}}
    = 2m_0|\Delta\Phi_N|/h$, testable prediction
    $\sim 0.5$~Hz for eV-scale photon pairs
    (Sec.~9).
  \item[\checkmark] \textbf{Gravitational
    birefringence}: one-loop vacuum polarisation
    prediction, distinguishable from QED and
    Lorentz-violation birefringence (Sec.~11).
\end{itemize}

\subsection*{Highest priority: would elevate fits
to predictions}

\begin{itemize}
  \item \textbf{Lattice simulation of k-oscillons}
    with $G_2(X)$ on the bi-conformal background
    ---to compute $\beta$, $\Phi_0$, and verify
    $\beta\Phi_0^2 = 0.853$ independently of the
    muon mass
    (Eq.~\ref{eq:amplitude_constraint}).
    \emph{This is the single most important
    outstanding calculation.}
  \item \textbf{$N = 4$ resonance search}
    at $\sim 329$~keV
    (\S\ref{sec:mathieu_numerical}).
    Experimental detection or exclusion
    directly tests the Mathieu stability
    structure.
\end{itemize}

\subsection*{Future work directions}

\begin{itemize}
  \item \textbf{Neutrino programme}: derive the
    torsional rigidity $\kappa_T$ from $G_2(X)$;
    solve Eq.~\eqref{eq:torsion_perturbation}
    for the three torsional eigenvalues and PMNS
    mixing angles; determine mass ordering and
    Majorana/Dirac character; derive the Weinberg
    angle from the torsion--transverse coupling
    ratio (\S\ref{sec:neutrinos}).
  \item \textbf{Massless sector}: derive photon
    coupling constants; solve
    Eq.~\eqref{eq:gluon_modes} for the oscillon
    cavity eigenmodes to verify whether exactly
    eight confined tensor modes emerge with SU(3)
    algebra; compute $\sigma$ from $G_2(X)$ and
    verify $V(L) \propto L$
    (\S\ref{sec:massless}).
  \item \textbf{Photon identification}:
    determine how the scalar-metric medium
    produces spin-1 transverse excitations (photons)
    in addition to the established spin-2 TT modes
    (classical gravitational waves of the medium itself,
    not a separate quantum excitation)
    (\S\ref{sec:massless}).
  \item \textbf{Electroweak programme}: derive
    $\text{SU}(2)_L \times \text{U}(1)_Y$ from
    the torsion--transverse coupling structure;
    compute the Weinberg angle $\theta_W$;
    verify $W/Z$ masses from
    $(R_\perp / \lambda_0)$ suppression
    (\S\ref{sec:neutrinos}, \S\ref{sec:charge}).
  \item \textbf{Higgs-mass programme}:
    build the first-principles amplitude-potential
    calculation of $V_{\mathrm{eff}}''(\Phi_0)$ using
    $\alpha$ and $M$ from the lattice simulation, and test
    whether it reproduces the $m_H=125$~GeV target
    (\S\ref{sec:higgs}).
  \item \textbf{Electric charge and
    $\alpha \approx 1/137$}: derive from the
    dipole-to-monopole amplitude ratio; verify
    $D_3$ phase locking reproduces quark charges
    $+2/3$, $-1/3$
    (\S\ref{sec:charge}).
  \item \textbf{Pre-geometric derivation of
    dimensionality}: reformulate
    \S\ref{sec:dimensions} in a tensor-network
    or graph-theoretic framework; complete the
    Ryu--Takayanagi reconstruction
    (\S\ref{sec:entanglement_geometry}).
  \item \textbf{Discrete substrate}: lattice
    or graph reproduction of the bi-conformal
    metric in the continuum limit.
\end{itemize}

\clearpage
\phantomsection
\section*{References}
% Consolidated bibliography for standalone and master builds
\begin{thebibliography}{99}

\bibitem{Einstein1915}
A.~Einstein,
``Die Feldgleichungen der Gravitation,''
Sitzungsber.\ Preuss.\ Akad.\ Wiss.\ Berlin, 844 (1915).

\bibitem{Will2014}
C.~M.~Will,
``The Confrontation between General Relativity and Experiment,''
Living Rev.\ Relativ.\ \textbf{17}, 4 (2014).

\bibitem{Will2018}
C.~M.~Will,
\textit{Theory and Experiment in Gravitational Physics},
2nd ed.\ (Cambridge University Press, 2018).

\bibitem{LIGO2016}
B.~P.~Abbott \textit{et al.}\ (LIGO/Virgo),
``Observation of Gravitational Waves from a Binary Black Hole Merger,''
Phys.\ Rev.\ Lett.\ \textbf{116}, 061102 (2016).

\bibitem{LIGO_GW170817}
B.~P.~Abbott \textit{et al.}\ (LIGO/Virgo/Fermi/INTEGRAL),
``Gravitational Waves and Gamma-Rays from a Binary Neutron Star Merger: GW170817 and GRB~170817A,''
Astrophys.\ J.\ Lett.\ \textbf{848}, L13 (2017).

\bibitem{TullyFisher1977}
R.~B.~Tully and J.~R.~Fisher,
``A New Method of Determining Distances to Galaxies,''
Astron.\ Astrophys.\ \textbf{54}, 661 (1977).

\bibitem{McGaugh2016}
S.~S.~McGaugh, F.~Lelli, and J.~M.~Schombert,
``Radial Acceleration Relation in Rotationally Supported Galaxies,''
Phys.\ Rev.\ Lett.\ \textbf{117}, 201101 (2016).

\bibitem{Milgrom1983a}
M.~Milgrom,
``A Modification of the Newtonian Dynamics as a Possible Alternative to the Hidden Mass Hypothesis,''
Astrophys.\ J.\ \textbf{270}, 365 (1983).

\bibitem{Milgrom1999}
M.~Milgrom,
``The Modified Dynamics as a Vacuum Effect,''
Phys.\ Lett.\ A \textbf{253}, 273 (1999).

\bibitem{Famaey2012}
B.~Famaey and S.~S.~McGaugh,
``Modified Newtonian Dynamics (MOND): Observational Phenomenology and Relativistic Extensions,''
Living Rev.\ Relativ.\ \textbf{15}, 10 (2012).

\bibitem{Angus2008}
G.~W.~Angus, B.~Famaey, and H.~S.~Zhao,
``Can MOND take a bullet? Analytical comparisons of three versions of MOND beyond spherical symmetry,''
Mon.\ Not.\ R.\ Astron.\ Soc.\ \textbf{371}, 138 (2006).

\bibitem{Bekenstein2004}
J.~D.~Bekenstein,
``Relativistic Gravitation Theory for the Modified Newtonian Dynamics Paradigm,''
Phys.\ Rev.\ D \textbf{70}, 083509 (2004).

\bibitem{Sotiriou2010}
T.~P.~Sotiriou and V.~Faraoni,
``$f(R)$ Theories of Gravity,''
Rev.\ Mod.\ Phys.\ \textbf{82}, 451 (2010).

\bibitem{BrunetonEsposito2007}
J.-P.~Bruneton and G.~Esposito-Far\`ese,
``Field-Theoretical Formulations of MOND-like Gravity,''
Phys.\ Rev.\ D \textbf{76}, 124012 (2007).

\bibitem{Penrose1965}
R.~Penrose,
``Gravitational Collapse and Space-Time Singularities,''
Phys.\ Rev.\ Lett.\ \textbf{14}, 57 (1965).

\bibitem{Hawking1970}
S.~W.~Hawking and R.~Penrose,
``The Singularities of Gravitational Collapse and Cosmology,''
Proc.\ R.\ Soc.\ Lond.\ A \textbf{314}, 529 (1970).

\bibitem{DeWitt1967}
B.~S.~DeWitt,
``Quantum Theory of Gravity. I. The Canonical Theory,''
Phys.\ Rev.\ \textbf{160}, 1113 (1967).

\bibitem{tHooft1974}
G.~'t~Hooft and M.~Veltman,
``One-Loop Divergences in the Theory of Gravitation,''
Ann.\ Inst.\ Henri Poincar\'e A \textbf{20}, 69 (1974).

\bibitem{BransDicke1961}
C.~Brans and R.~H.~Dicke,
``Mach's Principle and a Relativistic Theory of Gravitation,''
Phys.\ Rev.\ \textbf{124}, 925 (1961).

\bibitem{Horndeski1974}
G.~W.~Horndeski,
``Second-Order Scalar-Tensor Field Equations in a Four-Dimensional Space,''
Int.\ J.\ Theor.\ Phys.\ \textbf{10}, 363 (1974).

\bibitem{BelliniSawicki2014}
E.~Bellini and I.~Sawicki,
``Maximal Freedom at Minimum Cost: Linear Large-Scale Structure in General Modifications of Gravity,''
J.\ Cosmol.\ Astropart.\ Phys.\ \textbf{2014}(07), 050.

\bibitem{Bertotti2003}
B.~Bertotti, L.~Iess, and P.~Tortora,
``A Test of General Relativity Using Radio Links with the Cassini Spacecraft,''
Nature \textbf{425}, 374 (2003).

\bibitem{Williams2004}
J.~G.~Williams, S.~G.~Turyshev, and D.~H.~Boggs,
``Progress in Lunar Laser Ranging Tests of Relativistic Gravity,''
Phys.\ Rev.\ Lett.\ \textbf{93}, 261101 (2004).

\bibitem{Everitt2011}
C.~W.~F.~Everitt \textit{et al.},
``Gravity Probe B: Final Results of a Space Experiment to Test General Relativity,''
Phys.\ Rev.\ Lett.\ \textbf{106}, 221101 (2011).

\bibitem{Shapiro1964}
I.~I.~Shapiro,
``Fourth Test of General Relativity,''
Phys.\ Rev.\ Lett.\ \textbf{13}, 789 (1964).

\bibitem{Blanchet1993}
L.~Blanchet and G.~Sch\"afer,
``Gravitational Wave Tails and Binary Star Systems,''
Classical Quantum Gravity \textbf{10}, 2699 (1993).

\bibitem{BodennerWill2003}
J.~Bodenner and C.~M.~Will,
``Deflection of light to second order: A tool for illustrating principles of general relativity,''
Am.\ J.\ Phys.\ \textbf{71}, 770--773 (2003).

\bibitem{EpsteinShapiro1980}
R.~Epstein and I.~I.~Shapiro,
``Post-post-Newtonian deflection of light by the Sun,''
Phys.\ Rev.\ D \textbf{22}, 2947--2949 (1980).

\bibitem{KeetonPetters2005}
C.~R.~Keeton and A.~O.~Petters,
``Formalism for testing theories of gravity using lensing by compact objects,''
Phys.\ Rev.\ D \textbf{72}, 104006 (2005).

\bibitem{SKA2015}
M.~Kramer \textit{et al.},
``Strong-Field Tests of Gravity Using Pulsars and Black Holes,''
in \textit{Advancing Astrophysics with the Square Kilometre Array} (2015).

\bibitem{Carroll2003}
S.~M.~Carroll, M.~Hoffman, and M.~Trodden,
``Can the Dark Energy Equation-of-State Parameter $w$ Be Less Than $-1$?''
Phys.\ Rev.\ D \textbf{68}, 023509 (2003).

\bibitem{Cline2004}
J.~M.~Cline, S.~Jeon, and G.~D.~Moore,
``The Phantom Menaced: Constraints on Low-Energy Effective Ghosts,''
Phys.\ Rev.\ D \textbf{70}, 043543 (2004).

\bibitem{Sbisa2015}
F.~Sbis\`a,
``Classical and Quantum Ghosts,''
Eur.\ J.\ Phys.\ \textbf{36}, 015009 (2015).

\bibitem{Papallo2018}
G.~Papallo and H.~S.~Reall,
``On the Local Well-Posedness of Lovelock and Horndeski Theories,''
Phys.\ Rev.\ D \textbf{96}, 044019 (2017).

\bibitem{Kovacs2020}
A.~D.~Kov\'acs and H.~S.~Reall,
``Well-Posed Formulation of Scalar-Tensor Effective Field Theory,''
Phys.\ Rev.\ Lett.\ \textbf{124}, 221101 (2020).

\bibitem{Afshordi2007}
N.~Afshordi, D.~J.~H.~Chung, and G.~Geshnizjani,
``Cuscuton: A Causal Field Theory with an Infinite Speed of Sound,''
Phys.\ Rev.\ D \textbf{75}, 083513 (2007).

\bibitem{Chamseddine2013}
A.~H.~Chamseddine and V.~Mukhanov,
``Mimetic Dark Matter,''
J.\ High Energy Phys.\ \textbf{2013}(11), 135.

\bibitem{ArkaniHamed2004}
N.~Arkani-Hamed, H.-C.~Cheng, M.~A.~Luty, and S.~Mukohyama,
``Ghost Condensation and a Consistent Infrared Modification of Gravity,''
J.\ High Energy Phys.\ \textbf{2004}(05), 074.

\bibitem{SageManifolds}
E.~Gourgoulhon, M.~Bejger, and others,
``SageManifolds: Differential Geometry and Tensor Calculus with SageMath,''
\url{https://sagemanifolds.obspm.fr/} (2023).

\bibitem{EHT2022}
Event Horizon Telescope Collaboration,
``First Sagittarius~A* Event Horizon Telescope Results. I.,''
Astrophys.\ J.\ Lett.\ \textbf{930}, L12 (2022).

\bibitem{EHT2019M87}
Event Horizon Telescope Collaboration,
``First M87 Event Horizon Telescope Results. I. The Shadow of the Supermassive Black Hole,''
Astrophys.\ J.\ Lett.\ \textbf{875}, L1 (2019),
DOI: 10.3847/2041-8213/ab0ec7.

\bibitem{EHT2024M87}
Event Horizon Telescope Collaboration,
``The persistent shadow of the supermassive black hole of M87. I. Observations, calibration, imaging, and analysis,''
Astron.\ Astrophys.\ \textbf{681}, A79 (2024),
DOI: 10.1051/0004-6361/202347932.

\bibitem{Cardoso2019}
V.~Cardoso and P.~Pani,
``Testing the Nature of Dark Compact Objects: A Status Report,''
Living Rev.\ Relativ.\ \textbf{22}, 4 (2019).

\bibitem{Miani2023}
A.~Miani \emph{et al.},
``Constraints on the amplitude of gravitational wave echoes from black hole ringdown using minimal assumptions,''
Phys.\ Rev.\ D \textbf{108}, 064018 (2023),
DOI: 10.1103/PhysRevD.108.064018.

\bibitem{LVK2025TGR}
R.~Abbott \emph{et al.}\ (LIGO Scientific, Virgo, and KAGRA Collaborations),
``Tests of general relativity with GWTC-3,''
Phys.\ Rev.\ D \textbf{112}, 084080 (2025),
DOI: 10.1103/PhysRevD.112.084080.

\bibitem{Gogberashvili2018}
M.~Gogberashvili and D.~Singleton,
``Quantum particles trapped in a Schwarzschild black hole can escape in finite time,''
Phys.\ Lett.\ B \textbf{784}, 17--21 (2018).

\bibitem{Gogberashvili2019}
M.~Gogberashvili, A.~Sakharov, and E.~Sarkisyan-Grinbaum,
``Black holes as inner boundaries of the universe,''
Int.\ J.\ Mod.\ Phys.\ D \textbf{28}, 1950178 (2019).

\bibitem{McGaugh2011}
S.~S.~McGaugh,
``Novel Test of Modified Newtonian Dynamics with Gas Rich Galaxies,''
Phys.\ Rev.\ Lett.\ \textbf{106}, 121303 (2011).

\bibitem{Lotishvili2026}
G.~Lotishvili,
``Bi-Conformal Scalar--Tensor Gravity: Exact PPN Equivalence with General Relativity and Falsifiable Predictions,''
Zenodo (2026).

\bibitem{Lotishvili2026MOND}
G.~Lotishvili,
``Deriving MOND from Bi-Conformal Gravity: Vortex Closure, the Interpolating Function, and the Dark-Energy Connection,''
Zenodo (2026).

\bibitem{Lotishvili2026Q}
G.~Lotishvili,
``Quantum Extension of Bi-Conformal Scalar--Tensor Gravity: Canonical Quantization, Particle Spectrum, and Falsifiable Predictions,''
Zenodo (2026).

\bibitem{Hawking1976}
S.~W.~Hawking,
``Breakdown of Predictability in Gravitational Collapse,''
Phys.\ Rev.\ D \textbf{14}, 2460 (1976).

\bibitem{Kibble1976}
T.~W.~B.~Kibble,
``Topology of Cosmic Domains and Strings,''
J.\ Phys.\ A \textbf{9}, 1387 (1976).

\bibitem{DESI2024}
DESI Collaboration,
``DESI 2024 VI: Cosmological constraints from the measurements
of baryon acoustic oscillations,''
Phys.\ Rev.\ D \textbf{110} (2024) 023507.

\bibitem{DESIDR2}
DESI Collaboration, A.~Abdul-Karim \textit{et al.},
``DESI DR2 Results II: Measurements of Baryon Acoustic
Oscillations and Cosmological Constraints,''
Phys.\ Rev.\ D \textbf{112} (2025) 083515.

\bibitem{JWST2023}
S.~L.~Finkelstein \textit{et al.},
``A Long Time Ago in a Galaxy Far, Far Away: A Candidate $z \sim 12$ Galaxy,''
Astrophys.\ J.\ Lett.\ \textbf{940}, L55 (2022).

\bibitem{Euclid2024}
Euclid Collaboration,
``Euclid. I. Overview of the Euclid Mission,''
Astron.\ Astrophys.\ \textbf{662}, A112 (2022).

\bibitem{MICROSCOPE2022}
P.~Touboul \textit{et al.}\ (MICROSCOPE Collaboration),
``MICROSCOPE Mission: Final Results of the Test of the Equivalence Principle,''
Phys.\ Rev.\ Lett.\ \textbf{129}, 121102 (2022).

\bibitem{PoundRebka1960}
R.~V.~Pound and G.~A.~Rebka, Jr.,
``Apparent Weight of Photons,''
Phys.\ Rev.\ Lett.\ \textbf{4}, 337 (1960).

\bibitem{VessotLevine1980}
R.~F.~C.~Vessot and M.~W.~Levine,
``A Test of the Equivalence Principle Using a Space-Borne Clock,''
Gen.\ Relativ.\ Gravit.\ \textbf{10}, 181 (1979);
R.~F.~C.~Vessot \textit{et al.},
``Test of Relativistic Gravitation with a Space-Borne Hydrogen Maser,''
Phys.\ Rev.\ Lett.\ \textbf{45}, 2081 (1980).

\bibitem{MichelsonMorley1887}
A.~A.~Michelson and E.~W.~Morley,
``On the Relative Motion of the Earth and the Luminiferous Ether,''
Am.\ J.\ Sci.\ \textbf{34}, 333 (1887).

\bibitem{LambHydro}
H.~Lamb,
\textit{Hydrodynamics},
6th ed.\ (Cambridge University Press, 1932).

\bibitem{Fisher1948}
I.~Z.~Fisher,
``Scalar Mesostatic Field with Regard for Gravitational Effects,''
Zh.\ Eksp.\ Teor.\ Fiz.\ \textbf{18}, 636 (1948)
[arXiv:gr-qc/9911008 (English translation)].

\bibitem{JNW1968}
A.~I.~Janis, E.~T.~Newman, and J.~Winicour,
``Reality of the Schwarzschild Singularity,''
Phys.\ Rev.\ Lett.\ \textbf{20}, 878 (1968).

\bibitem{Bronnikov1973}
K.~A.~Bronnikov,
``Scalar-Tensor Theory and Scalar Charge,''
Acta Phys.\ Pol.\ B \textbf{4}, 251 (1973).

\bibitem{Ellis1973}
H.~G.~Ellis,
``Ether Flow through a Drainhole: A Particle Model in General Relativity,''
J.\ Math.\ Phys.\ \textbf{14}, 104 (1973).

\bibitem{Papapetrou1954}
A.~Papapetrou,
``Eine Theorie des Gravitationsfeldes mit einer Feldfunktion,''
Z.\ Phys.\ \textbf{139}, 518 (1954).

\bibitem{Yilmaz1958}
H.~Yilmaz,
``New Approach to General Relativity,''
Phys.\ Rev.\ \textbf{111}, 1417 (1958).

\bibitem{Boonserm2020}
P.~Boonserm, T.~Ngampitipan, A.~Simpson, and M.~Visser,
``Triple Path to the Exponential Metric,''
Found.\ Phys.\ \textbf{50}, 1368 (2020).

\bibitem{Bronnikov2012}
K.~A.~Bronnikov and S.~G.~Rubin,
``Instabilities of Wormholes and Regular Black Holes Supported by a Phantom Scalar Field,''
Phys.\ Rev.\ D \textbf{86}, 024028 (2012).

\bibitem{LISA2017}
P.~Amaro-Seoane \textit{et al.},
``Laser Interferometer Space Antenna,''
arXiv:1702.00786 [astro-ph.IM] (2017).

\bibitem{Adelberger2003}
E.~G.~Adelberger, B.~R.~Heckel, and A.~E.~Nelson,
``Tests of the Gravitational Inverse-Square Law,''
Annu.\ Rev.\ Nucl.\ Part.\ Sci.\ \textbf{53}, 77 (2003).

\bibitem{Clowe2006}
D.~Clowe, M.~Brada\v{c}, A.~H.~Gonzalez, M.~Markevitch,
S.~W.~Randall, C.~Jones, and D.~Zaritsky,
``A Direct Empirical Proof of the Existence of Dark Matter,''
Astrophys.\ J.\ Lett.\ \textbf{648}, L109 (2006).

\bibitem{Freire2012}
P.~C.~C.~Freire \textit{et al.},
``The relativistic pulsar--white dwarf binary PSR~J1738+0333~II.\ The most stringent test of scalar-tensor gravity,''
Mon.\ Not.\ R.\ Astron.\ Soc.\ \textbf{423}, 3328 (2012).

\bibitem{Archibald2018}
A.~M.~Archibald \textit{et al.},
``Universality of free fall from the orbital motion of a pulsar in a stellar triple system,''
Nature \textbf{559}, 73 (2018).

\bibitem{Voisin2025}
G.~Voisin \textit{et al.},
``Explanation of the exceptionally strong timing noise of PSR~J0337+1715 by a circum-ternary planet and consequences for gravity tests,''
Astron.\ Astrophys.\ \textbf{693}, A143 (2025),
DOI: 10.1051/0004-6361/202452100.

\bibitem{Clarke1993}
C.~J.~S.~Clarke,
\emph{The Analysis of Space-Time Singularities}
(Cambridge University Press, 1993).

\bibitem{Racz2023}
I.~R\'acz,
``Spacetime singularities and curvature blow-ups,''
Gen.\ Relativ.\ Gravit.\ \textbf{55}, 3 (2023),
DOI: 10.1007/s10714-022-03053-9.

\bibitem{Abbott2017GW}
B.~P.~Abbott \emph{et al.} (LIGO Scientific and Virgo Collaborations),
\emph{GW170817: Observation of Gravitational Waves from a Binary Neutron Star Inspiral},
Phys.\ Rev.\ Lett.\ \textbf{119}, 161101 (2017).

\bibitem{SkordisZlosnik2021}
C.~Skordis and T.~Z\l{}o\'snik,
\emph{New relativistic theory for Modified Newtonian Dynamics},
Phys.\ Rev.\ Lett.\ \textbf{127}, 161302 (2021).

\bibitem{Barbour1999}
J.~Barbour,
\emph{The End of Time: The Next Revolution in Physics}
(Oxford University Press, 1999).

\bibitem{ConnesRovelli1994}
A.~Connes and C.~Rovelli,
``Von~Neumann algebra automorphisms and time-thermodynamics relation
in generally covariant quantum theories,''
Class.\ Quantum Grav.\ \textbf{11}, 2899 (1994).

\bibitem{Whitehead1929}
A.~N.~Whitehead,
\emph{Process and Reality}
(Macmillan, 1929).

\bibitem{Penrose2010}
R.~Penrose,
\emph{Cycles of Time: An Extraordinary New View of the Universe}
(Bodley Head, 2010).

% --- Particle / Mathieu / holography supplement (FrequencyToMass, quantum sector) ---

\bibitem{Koide1982}
Y.~Koide,
``Fermion--boson two-body bound systems and the fermion mass hierarchy,''
Lett.\ Nuovo Cim.\ \textbf{34}, 201 (1982).

\bibitem{DLMFMathieu}
NIST Digital Library of Mathematical Functions,
Chap.~28, \textit{Mathieu Functions and Hill's Equation},
\url{https://dlmf.nist.gov/28},
NIST, Gaithersburg, MD (release 1.2.1, 2024).

\bibitem{McLachlan1947}
N.~W.~McLachlan,
\emph{Theory and Application of Mathieu Functions}
(Oxford University Press, 1947).

\bibitem{PDG2024}
R.~L.~Workman \textit{et al.}\ (Particle Data Group),
``Review of Particle Physics,''
Prog.\ Theor.\ Exp.\ Phys.\ \textbf{2024}, 083C01 (2024).

\bibitem{RyuTakayanagi2006}
S.~Ryu and T.~Takayanagi,
``Holographic Derivation of the Entanglement Entropy from AdS/CFT,''
Phys.\ Rev.\ Lett.\ \textbf{96}, 181602 (2006).

\bibitem{NishiokaTakayanagi2009}
T.~Nishioka, S.~Ryu, and T.~Takayanagi,
``Holographic Entanglement Entropy: An Overview,''
J.\ Phys.\ A \textbf{42}, 504008 (2009).

\bibitem{Bombelli1986}
L.~Bombelli, R.~K.~Koul, J.~Lee, and R.~D.~Sorkin,
``A Quantum Source of Entropy for Black Holes,''
Phys.\ Rev.\ D \textbf{34}, 373 (1986).

\bibitem{MaldacenaSusskind2013}
J.~Maldacena and L.~Susskind,
``Cool horizons for entangled black holes,''
Fortschr.\ Phys.\ \textbf{61}, 781 (2013).

\bibitem{Maldacena1998}
J.~M.~Maldacena,
``The Large $N$ limit of superconformal field theories and supergravity,''
Adv.\ Theor.\ Math.\ Phys.\ \textbf{2}, 231 (1998).

\bibitem{Leggett2001}
A.~J.~Leggett,
``Bose--Einstein condensation in the alkali gases: Some fundamental concepts,''
Rev.\ Mod.\ Phys.\ \textbf{73}, 307 (2001).

\bibitem{PethickSmith2008}
C.~J.~Pethick and H.~Smith,
\emph{Bose--Einstein Condensation in Dilute Gases},
2nd ed.\ (Cambridge University Press, 2008).

\bibitem{GleiserWatkins1993}
I.~Gleiser and R.~Watkins,
``Towards a theory of oscillons,''
Phys.\ Rev.\ D \textbf{47}, 5197 (1993).

\bibitem{Copeland1995}
E.~J.~Copeland, M.~Gleiser, and H.-R.~Muller,
``Oscillons: Resonant configurations during damping in one dimension,''
Phys.\ Rev.\ D \textbf{52}, 1920 (1995).

\bibitem{Nicolis2009}
A.~Nicolis, R.~Rattazzi, and E.~Trincherini,
``The Galileon as a local modification of gravity,''
Phys.\ Rev.\ D \textbf{79}, 064036 (2009).

\bibitem{Deffayet2009}
C.~Deffayet, G.~Esposito-Far\`ese, and A.~Vikman,
``Covariantizing Galileons,''
Phys.\ Rev.\ D \textbf{79}, 084003 (2009).

\bibitem{Goldstone1961}
J.~Goldstone,
``Field theories with ``superconductor'' solutions,''
Nuovo Cimento \textbf{19}, 154 (1961).

\bibitem{ColemanWeinberg1973}
S.~Coleman and E.~Weinberg,
``Radiative Corrections as the Origin of Spontaneous Symmetry Breaking,''
Phys.\ Rev.\ D \textbf{7}, 1888 (1973).

\bibitem{Higgs1964}
P.~W.~Higgs,
``Broken Symmetries and the Masses of Gauge Bosons,''
Phys.\ Rev.\ Lett.\ \textbf{13}, 508 (1964).

\bibitem{EnglertBrout1964}
F.~Englert and R.~Brout,
``Broken Symmetry and the Mass of Gauge Vector Mesons,''
Phys.\ Rev.\ Lett.\ \textbf{13}, 321 (1964).

\bibitem{Guralnik1964}
G.~S.~Guralnik, C.~R.~Hagen, and T.~W.~B.~Kibble,
``Global Conservation Laws and Massless Particles,''
Phys.\ Rev.\ Lett.\ \textbf{13}, 585 (1964).

\bibitem{Hopkins2010}
P.~F.~Hopkins \textit{et al.},
``Mergers, Minor Mergers, and the Origin of Disklike Galaxies,''
Astrophys.\ J.\ \textbf{715}, 202 (2010).

\bibitem{Stewart2008}
K.~R.~Stewart \textit{et al.},
``Merger Histories of Galaxy Halos and Implications for Disk Survival,''
Astrophys.\ J.\ \textbf{683}, 597 (2008).

\end{thebibliography}


\end{document}
